\documentclass[11pt]{book}
\usepackage[utf8]{inputenc}
\usepackage{amsmath}
\usepackage{amsfonts}
\usepackage{amssymb}
\usepackage{amsthm}
\usepackage{url}
\usepackage{hyperref}
\usepackage[spanish]{babel}
\usepackage{booktabs}
\usepackage{subfigure}
\usepackage{enumitem}
\usepackage{multicol}

%\usepackage{amsmath, amssymb, amsthm}
\usepackage{graphicx}
\usepackage{float}
%\usepackage{booktabs}
\usepackage{multirow}
\usepackage{caption}
\usepackage{subcaption}
%\usepackage{hyperref}


\usepackage[letterpaper,top=3cm,bottom=3cm,left=3cm,right=3cm,marginparwidth=1.75cm]{geometry}
\usepackage{tikz}
\usetikzlibrary{decorations.pathmorphing}

\newcounter{cuent}
%\newcommand{\proba}[1]{\stepcounter{cuent}{\alph{cuent})\quad}
%\displaystyle#1\qquad}
\newcommand{\proba}[1]{\stepcounter{cuent}{\alph{cuent})\ }
\displaystyle#1\hfill}
\newcommand{\cuento}{\setcounter{cuent}{0}}
\newcommand{\R}{\mathbb{R}}


\usepackage{listings}
\usepackage{xcolor}

% Define a Python style
\lstdefinestyle{pythonstyle}{
    language=Python,
    basicstyle=\ttfamily\footnotesize,
    keywordstyle=\color{blue}\bfseries,
    commentstyle=\color{gray}\itshape,
    stringstyle=\color{red},
    numberstyle=\tiny\color{gray},
    numbers=left,
    stepnumber=1,
    breaklines=true,
    frame=single,
    backgroundcolor=\color{black!5},
    captionpos=b
}

\lstdefinestyle{juliastyle}{
    %language=Julia,
    basicstyle=\ttfamily\footnotesize,
    keywordstyle=\color{blue}\bfseries,
    commentstyle=\color{gray}\itshape,
    stringstyle=\color{red},
    numberstyle=\tiny\color{gray},
    numbers=left,
    stepnumber=1,
    breaklines=true,
    frame=single,
    backgroundcolor=\color{black!5},
    captionpos=b
}

\lstdefinestyle{matlabstyle}{
  language=Matlab,
  basicstyle=\ttfamily\footnotesize,
  keywordstyle=\color{blue},
  commentstyle=\color{green!50!black},
  stringstyle=\color{red},
  numbers=left,
  numberstyle=\tiny,
  stepnumber=1,
  breaklines=true,
  frame=single,
  captionpos=b
}

%%%%
%\usepackage[utf8]{inputenc}
%\usepackage{amsmath}
%\usepackage{amssymb}
%\usepackage{graphicx}
%\usepackage{listings}
%\usepackage{xcolor}
%\usepackage{hyperref}

\definecolor{codegreen}{rgb}{0,0.6,0}
\definecolor{codegray}{rgb}{0.5,0.5,0.5}
\definecolor{codepurple}{rgb}{0.58,0,0.82}
\definecolor{backcolour}{rgb}{0.95,0.95,0.92}

\lstdefinestyle{mystyle}{
    backgroundcolor=\color{backcolour},
    commentstyle=\color{codegreen},
    keywordstyle=\color{magenta},
    numberstyle=\tiny\color{codegray},
    stringstyle=\color{codepurple},
    basicstyle=\ttfamily\footnotesize,
    breakatwhitespace=false,
    breaklines=true,
    captionpos=b,
    keepspaces=true,
    numbers=left,
    numbersep=5pt,
    showspaces=false,
    showstringspaces=false,
    showtabs=false,
    tabsize=2
}

\lstset{style=mystyle}

%%%%

%\documentclass{article}
%\usepackage[utf8]{inputenc}
%\usepackage{amsmath}
%\usepackage{amsthm}
%\usepackage{amssymb}
%\usepackage{enumitem}
%\usepackage{tikz}

% Estilos para los entornos
\newtheoremstyle{definicion}
    {12pt}   % Espacio arriba
    {12pt}   % Espacio abajo
    {\itshape} % Fuente del cuerpo
    {}        % Sangría
    {\bfseries} % Fuente del encabezado
    {.}       % Puntuación después del encabezado
    { }       % Espacio después del encabezado
    {}        % Notas al final

\newtheoremstyle{teorema}
    {12pt}
    {12pt}
    {\itshape}
    {}
    {\bfseries}
    {.}
    { }
    {}

\newtheoremstyle{ejemplo}
    {12pt}
    {12pt}
    {}
    {}
    {\bfseries}
    {.}
    { }
    {}

% Definición de entornos
\theoremstyle{definicion}
\newtheorem{definicion}{Definición}[section]

\theoremstyle{teorema}
\newtheorem{teorema}{Teorema}[section]

\theoremstyle{ejemplo}
\newtheorem{ejemplo}{Ejemplo}[section]


\newtheorem{theorem}{Theorem}
\newtheorem{cor}{Corollary}
\newtheorem{lemma}{Lemma}
\newtheorem{proposition}{Proposition}
%\newtheorem{definition}{Definition}

\newtheorem{definition}{Definición}
\newtheorem{example}{Ejemplo}
%\newtheorem{definicion}{Definición}
%\newtheorem{ejemplo}{Ejemplo}

\title{C\'alculo diferencial e integral II}
%ODEs and PDEs}
\author{Franz Chouly}
\date{\today}

\begin{document}

\maketitle

\tableofcontents

\chapter{Introducci\'on}

%La idea de este curso es dar herramientas basicas para la fisica clasica: mecanica de solidos y de fluidos, termica, electromagnetismo, relatividad, etc.
\noindent Para la motivaci\'on, vease por ejemplo \cite{chorin1993}.
El libro de base es el de Cálculo de varias variables de James Stewart~\cite{stewartmanyvariables} (versi\'on en espa\~nol), ver tambi\'en
\url{https://www.stewartcalculus.com/}.
%Referencia en electromagnestismo.
%Ver el capitulo de Zuqi y preguntar. Libro de Monk pero demasiado avanzado. Preguntar a Chiara, Alessandra. 
%Chorin Marsden. Campo de velocidades de un fluido. Rio por ejemplo. Divergencia nula. Lineas de corriente. Integral en una superficie para medir la fuerza
%
%Stewart ...

\chapter{Conjuntos en $\mathbb{R}^n$}

\section{Temario del curso}

\begin{itemize}
    \item Producto escalar, norma y distancia.
    \item Representación gráfica de conjuntos dados por ecuaciones e inecuaciones. 
\end{itemize}

\section{Notas}

%Geometria del plano y del espacio.
\noindent
Aqui volvemos a ver nociones de base para ubicarse en el plano y el espacio.

\subsection{El plano}

El plano es el conjunto de todos los puntos $p=(x,y)$ cuando $x,y$ son dos reales. Lo notamos $\mathbb{R}^2$. Tambi\'en un vector $v=(v_x,v_y)$ es definido sin ambiguidad por sus coordenadas.  Dos puntos $q$ y $p$ definen un vector $v=q-p$. Los vectores de base se notan $e_x = (1,0)$ y $e_y = (0,1)$ y asi todo vector $v$ se puede escribir
\[
v = 
(v_x,v_y)
= 
v_x (1,0) + v_y (0,1) =
v_x e_x + v_y e_y.
\]

\subsection{Operaciones con vectores}
\begin{definition}
La suma de dos vectores \(u = (u_x, u_y)\) y \(v = (v_x, v_y)\) en \(\mathbb{R}^2\) se define como \(u + v = (u_x + v_y, u_x + v_y)\).
\end{definition}

\begin{example}
Si \(u = (1, 2)\) y \(v = (3, 4)\), entonces \(u + v = (1 + 3, 2 + 4) = (4, 6)\).
\end{example}

Tambi\'en se puede multiplicar un vector por un escalar (tenemos una estructura de espacio vectorial).

\subsection{Producto escalar y norma}

\noindent 
El producto escalar y la norma se pueden definir como lo siguiente:
\begin{definition}
El producto escalar entre dos vectores $u=(u_x,u_y)$ y $v=(v_x,v_y)$ es definido mediante la siguiente formula
\[
u \cdot v = \langle u , v \rangle = u_x v_x + u_y v_y 
\]
(se puede notar de dos formas)
y la norma de un vector $u$ es dada
por
\[
\| u \| = \sqrt{\langle u , u \rangle} = \sqrt{u_x^2 + u_y^2}.
\]
%
%La norma de un vector \(u = (u_1, u_2)\) en \(\mathbb{R}^2\) se define como \(\|u\| = \sqrt{u_1^2 + u_2^2}\).
\end{definition}

\begin{example}
Si \(u = (3, 4)\), entonces \(\|u\| = \sqrt{3^2 + 4^2} = 5\).
\end{example}

Sea ahora un vector $u$ con $\|u\|=1$ y que hace un angulo $\theta$ con $e_x$. Por definici\'on de $\cos(\theta)$:
\[
\cos(\theta) = u_x 
= 1 u_x + 0 u_y = \langle u , e_x \rangle
\]
y sean ahora $u$ con $\|u\|=1$ y $v$ con $\|v\|=1$. Se puede mostrar que la formula arriba se generaliza, osea % por invariancia, cambio de coordenadas
\[
\cos(\theta) = u_x v_x + u_y v_y = \langle u , v \rangle,
\]
con esta vez $\theta$ que define el angulo entre los dos vectores.
Ahora sean $u$ y $v$ dos vectores cualquiera:
\[
\langle u , v \rangle = 
\left \langle \|u\| \frac{u}{\|u\|} , \|v\| \frac{v}{\|v\|} \right \rangle
=
\|u\| \|v\| \left \langle \frac{u}{\|u\|} ,  \frac{v}{\|v\|} \right \rangle
=
\|u\| \|v\| \cos(\theta).
\]

\subsection{Distancia}

\noindent
La distancia entre dos puntos $p=(p_x,p_y)$ y $q=(q_x,q_y)$ es la norma del vector $p-q$.

\subsection{Recta y distancia a una recta}

Una recta $\Delta$ puede ser vista como el conjunto de puntos que pasa por un punto $m_0=(x_0,y_0)$ y ortogonal a un vector $v=(a,b)$ con $a$ o $b$ diferente de $0$. Cualquier punto $m=(x,y)$ de la recta verifica
\[
\langle m - m_0 , v \rangle = 0
\]
o con coordenadas
\[
(x-x_0)a + (y-y_0)b = 0,
\]
definiendo $c=a x_0 + b y_0$ obtenemos
la ecuaci\'on de la recta $\Delta$:
\[
a x + b y = c.
\]
Tambi\'en la recta se puede definir a traves del vector director $u=(-b,a)$:
\[
x = x_0 - b t, \quad y = y_0 + a t,
\]
para todo real $t$. Volvemos a encontrar
\[
a x + by = a x_0 + b y0 -abt + abt = c. 
\]
Resumimos:
%\subsection{Ecuación de una recta en \(\mathbb{R}^2\)}
\begin{definition}
La ecuación de una recta en \(\mathbb{R}^2\) puede ser dada en la forma \(ax + by = c\), donde \(a\) y \(b\) no son ambos cero.
\end{definition}

\begin{example}
La ecuación \(2x + 3y = 6\) representa una recta en \(\mathbb{R}^2\).
\end{example}

Ahora veamos como calcular la distancia de un punto del plano $p=(p_x,p_y)$ a la recta $\Delta$ de ecuacion $ax + by = c$. Definimos $\tilde v = v / \| v \|$ el vector unitario ortogonal a $\Delta$. Tenemos
\[
\tilde v = \left ( 
\frac{a}{\sqrt{a^2+b^2}},
\frac{b}{\sqrt{a^2+b^2}} \right ).
\]
Por propriedad de la proyecci\'on ortogonal, tenemos
\[
d(p,\Delta)= 
\left | \langle \tilde v, p - m_0 \rangle \ \right |
= \frac{1}{\sqrt{a^2+b^2}}
\left | a (p_x - x_0) + b (p_y - x_0 ) \right |
\]
y entonces
\[
d(p,\Delta)= 
\frac{1}{\sqrt{a^2+b^2}}
\left | a p_x + b p_y - c \right |.
\]
La distancia obviamente no depende de $x_0,y_0$, y volvemos a encontrar que si $p$ esta en la recta, su distancia a la recta es $0$.

\subsection{Intersección de rectas}
\begin{definition}
La intersección de dos rectas en \(\mathbb{R}^2\) es el conjunto de puntos \((x, y)\) que satisfacen ambas ecuaciones de las rectas.
\end{definition}
Se obtiene solucionando el sistema
\begin{eqnarray*}
    a x + by &=& c\\
    a' x + b'y &=& c'.
\end{eqnarray*}
Esta intersecci\'on puede ser vacia, reducida a un punto, o una recta entera. Veamos un ejemplo.

\begin{example}
Las rectas (\(x + y = 2\)) y (\(x - y = 0\)) se intersectan en el punto \((1, 1)\).
\end{example}


\subsection{Un ejemplo detallado}

\noindent Dados los vectores $u=(0,2),\ v=(2,0)$ de $\R^2$ hallamos
\begin{enumerate}
    \item $\parallel u + v \parallel = \sqrt{2^2 + 2^2} = \sqrt{8} = 2 \sqrt{2} \simeq 2,828$.
    \item ${d(u,v)} =
    \parallel u - v \parallel
    = 2 \sqrt{2} \simeq 2,828$. 
    \item $\left<u,v\right> = 0\times 2 + 2 \times 0 = 0$.
    \item ${\mbox{el \'angulo entre } u \mbox{ y } v}$
    \item ${\mbox{la recta } r \mbox{ paralela a }u\mbox{ por }(0,2)}$: 
    $y=2$  
    \item ${\mbox{la recta } s \mbox{ perpendicular a }u\mbox{ por }(2,0)}$:
    $x=2$.
    \item ${r\cap s} = (2,2).$

\end{enumerate}

\subsection{Conjuntos}

\noindent Veamos primero una definici\'on y un ejemplo.
\begin{definition}
Un conjunto en \(\mathbb{R}^2\) es una colección de pares ordenados \((x, y)\) que satisfacen una o más ecuaciones o inecuaciones.
\end{definition}

\begin{example}
El conjunto \(A = \{(x, y) \in \mathbb{R}^2 : y = \sin(x), -\frac{\pi}{2} \leq x \leq \frac{\pi}{2}\}\) representa una curva senoidal en el plano.
\end{example}

Veamos otro respeto a la representación gráfica de conjuntos dados por ecuaciones e inecuaciones:
$A=\{(x,y)\in\R^2:\ -2\leq x \leq y^2,\ =1\leq x \leq 1\}$ y
$A=\{(x,y)\in\R^2:\ x-y=-1\}$.

\subsection{Dimension tres y m\'as}

Todas las nociones previas se extenden sin (mayor) problema a la dimension tres o a mas dimensi\'on, es cosa de agregar una (o mas) coordenadas. 

%subsection{Diferenciabilidad}
%\begin{definition}
%Una función \(f: \mathbb{R}^2 \to \mathbb{R}\) es diferenciable en un punto \((a, b)\) si existe un plano tangente a la gráfica de \(f\) en ese punto.
%\end{definition}

%\begin{example}
%La función \(f(x, y) = x^2 + y^2\) es diferenciable en todo \(\mathbb{R}^2\).
%\end{example}

\subsection{Conclusión}
Este curso ha cubierto las nociones básicas de álgebra lineal y geometría analítica. Se recomienda practicar con más ejemplos y ejercicios para una mejor comprensión.

\newpage

\subsection{Complemento sobre las cónicas}
%***
%\documentclass[12pt,a4paper]{article}
%\usepackage[utf8]{inputenc}
%\usepackage[spanish]{babel}
%\usepackage{amsmath}
%\usepackage{amsfonts}
%\usepackage{amssymb}
%\usepackage{graphicx}
%\usepackage{enumitem}
%\usepackage{geometry}
%\geometry{left=2cm,right=2cm,top=2cm,bottom=2cm}

%\title{Notas de Clase: Las Cónicas}
%\author{Franz Chouly}
%\date{\today}

%\begin{document}

%\maketitle

%\section{Introducción}
Las \textbf{cónicas} son curvas que se obtienen al intersectar un plano con un cono de revolución. Dependiendo del ángulo de intersección, se obtienen diferentes tipos de cónicas: \textbf{circunferencia}, \textbf{elipse}, \textbf{parábola} e \textbf{hipérbola}.

\subsubsection{Definición General}
Una cónica es el lugar geométrico de los puntos \( P \) del plano tales que la relación de sus distancias a un punto fijo (foco) y a una recta fija (directriz) es constante. Esta relación se denomina \textbf{excentricidad} (\( e \)):
\[
\frac{d(P, \text{foco})}{d(P, \text{directriz})} = e.
\]
\begin{itemize}
    \item Si \( e = 1 \): Parábola.
    \item Si \( e < 1 \): Elipse (si \( e = 0 \), es una circunferencia).
    \item Si \( e > 1 \): Hipérbola.
\end{itemize}

\subsubsection{La Circunferencia}
\paragraph{Definición}
La circunferencia es el conjunto de puntos del plano que equidistan de un punto fijo llamado \textbf{centro}.

\paragraph{Ecuación Canónica}
La ecuación canónica de una circunferencia con centro en \( (h, k) \) y radio \( r \) es:
\[
(x - h)^2 + (y - k)^2 = r^2
\]

\paragraph{Propiedades}
\begin{itemize}
    \item Todos los puntos están a la misma distancia del centro.
    \item Es un caso especial de la elipse donde los dos focos coinciden.
\end{itemize}

\subsubsection{La Elipse}
\paragraph{Definición}
La elipse es el conjunto de puntos del plano cuya suma de distancias a dos puntos fijos (focos) es constante.

\paragraph{Ecuación Canónica}
Para una elipse centrada en el origen con eje mayor \( 2a \) y eje menor \( 2b \):
\[
\frac{x^2}{a^2} + \frac{y^2}{b^2} = 1
\]

\paragraph{Propiedades}
\begin{itemize}
    \item Los focos están ubicados en \( (\pm c, 0) \), donde \( c^2 = a^2 - b^2 \).
    \item La excentricidad es \[ e = \frac{c}{a}.\]
\end{itemize}

\subsubsection{La Parábola}
\paragraph{Definición}
La parábola es el conjunto de puntos del plano que equidistan de un punto fijo (foco) y una recta fija (directriz).

\paragraph{Ecuación Canónica}
Para una parábola con vértice en el origen y foco en \( (a, 0) \):
\[
y^2 = 4ax.
\]
\paragraph{Propiedades}
\begin{itemize}
    \item Tiene un solo foco y una directriz.
    \item Su excentricidad es \( e = 1 \).
\end{itemize}

\subsubsection{La Hipérbola}
\paragraph{Definición}
La hipérbola es el conjunto de puntos del plano cuya diferencia de distancias a dos puntos fijos (focos) es constante.

\paragraph{Ecuación Canónica}
Para una hipérbola centrada en el origen con eje transversal \( 2a \):
\[
\frac{x^2}{a^2} - \frac{y^2}{b^2} = 1.
\]

\paragraph{Propiedades}
\begin{itemize}
    \item Tiene dos focos y dos asíntotas.
    \item La excentricidad es \( e = \frac{c}{a} \), donde \( c^2 = a^2 + b^2 \).
\end{itemize}

\subsubsection{Aplicaciones}
Las cónicas tienen múltiples aplicaciones en la vida real:
\begin{itemize}
    \item Las órbitas de los planetas son elipses.
    \item Los espejos parabólicos se utilizan en telescopios y antenas.
    \item Las hipérbolas se usan en sistemas de navegación (LORAN).
\end{itemize}

\subsubsection{Conclusión}
Las cónicas son fundamentales en matemáticas y física. Su estudio permite modelar fenómenos naturales y diseñar tecnologías avanzadas.
Complementos interesantes son disponibles en la web, particularmente:
\url{http://uruguayeduca.anep.edu.uy/recursos-educativos/3700} y
\url{https://www.bibmath.net/dico/index.php?action=affiche&quoi=./c/conique.html&special=imprimable}.

%\end{document}

%\section{Mistral (tentative)}

%\documentclass{article}
%\usepackage{amsmath}
%\usepackage{amssymb}
%\usepackage{graphicx}


%\begin{document}

%\section*{Apuntes de Curso}

%\subsection*{Conjuntos en $\mathbb{R}^2$}

%A continuación se presentan varios conjuntos en $\mathbb{R}^2$:

%\begin{enumerate}
%    \item $A=\{(x,y)\in\mathbb{R}^2: y=\sin(x), \frac{-\pi}{2}\leq x \leq \frac{\pi}{2}\}$.
%    \item $A=\{(x,y)\in\mathbb{R}^2: -1\leq x \leq y^2, 2\leq y < 3\}$.
%    \item $A=\{(x,y)\in\mathbb{R}^2: -\sqrt{1-x^2}\leq y \leq \sqrt{1-x^2}, -1\leq x \leq 1\}$.
%    \item $A=\{(x,y)\in\mathbb{R}^2: 3x+2y=-1\}$.
%    \item $A=\{(x,y)\in\mathbb{R}^2: x-5y=-1, -2\leq x < 2\}$.
%    \item $A=\{(x,y)\in\mathbb{R}^2: x^2+y^2-2x-4y+5=3\}$.
%    \item $A=\{(x,y)\in\mathbb{R}^2: x=y^2-2y+1, -1\leq x \leq 1\}$.
%    \item $A=\{(x,y)\in\mathbb{R}^2: 5x^2+2y^2=2\}$.
%    \item $A=\{(x,y)\in\mathbb{R}^2: -3x^2+y^2=1\}$.
%    \item $B((1,-1),\frac{1}{2})$.
%\end{enumerate}

%\subsection*{Ejercicios con Vectores en $\mathbb{R}^2$}

%Dados los vectores $u=(1,1), v=(2,1)$ de $\mathbb{R}^2$, hallar:

%\begin{enumerate}[label=(\alph*)]
%    \item $\|u + v\|$
%    \item $d(u,v)$
%    \item $\langle u,v \rangle$
%    \item El ángulo entre $u$ y $v$
%    \item La recta $r$ paralela a $u$ por $(0,1)$
%    \item La recta $s$ perpendicular a $v$ por $(0,0)$
%    \item $r \cap s$
%\end{enumerate}

%\subsection*{Proyección sobre un Vector}
%
%Muestra que en general, dado un vector $v$ del plano $\mathbb{R}^2$, el conjunto $K_v$ de puntos cuya proyección sobre $v$ es cero es una recta que pasa por el origen. Encuentra dicha recta para $v=(2,3)$. Muestra que dados dos vectores $v$ y $v'$, dichas rectas coinciden si y solo si $v$ y $v'$ son colineales.

%\subsection*{Conjuntos en $\mathbb{R}^3$}

%A continuación se presentan varios conjuntos en $\mathbb{R}^3$:

%\begin{enumerate}
%    \item $A=\{(x,y,z)\in\mathbb{R}^3: 3x+2y-z=1\}$.
%    \item $A=\{(x,y,z)\in\mathbb{R}^3: x+y-z=0, -2x-2y+2z=1\}$.
%    \item $A=\{(x,y,z)\in\mathbb{R}^3: x^2+y^2=1\}$.
%    \item $A=\{(x,y,z)\in\mathbb{R}^3: 3x^2+y^2=2, -1\leq z \leq 2\}$.
%    \item $A=\{(x,y,z)\in\mathbb{R}^3: x^2-y^2=1, z=1\}$.
%    \item $A=\{(x,y,z)\in\mathbb{R}^3: x^2+y^2+(z-2)^2=10\}$.
%    \item $A=\{(x,y,z)\in\mathbb{R}^3: x^2+y^2+z^2=1, (x-1)^2+y^2=1\}$.
%    \item $A=\{(x,y,z)\in\mathbb{R}^3: 2x^2+y^2+z^2=2, x+y=1\}$.
%    \item $A=\{(x,y,z)\in\mathbb{R}^3: x^2+y^2-4z^2=1\}$.
%\end{enumerate}

%\subsection*{Ejercicios con Vectores en $\mathbb{R}^3$}

%Dados los vectores $u=(1,1,1), v=(2,1,0)$ de $\mathbb{R}^3$, hallar:

%\begin{enumerate}[label=(\alph*)]
%    \item $\|u + v\|$
%    \item $d(u,v)$
%    \item $\langle u,v \rangle$
%    \item El ángulo entre $u$ y $v$
%    \item La recta $r$ paralela a $u$ por $(0,0,1)$
%    \item El plano $\Pi$ perpendicular a $v$ por $(0,0,0)$
%    \item $r \cap s$
%\end{enumerate}

%\subsection*{Proyección sobre un Vector en $\mathbb{R}^3$}

%Muestra que en general, dado un vector $v$ del espacio $\mathbb{R}^3$, el conjunto $K_v$ de puntos cuya proyección sobre $v$ es cero es un plano que pasa por el origen. Encuentra dicho plano para $v=(1,-1,0)$. Muestra que dados dos vectores $v$ y $v'$, dichos planos coinciden si y solo si $v$ y $v'$ son colineales.

%\subsection*{Análisis de un Subconjunto en $\mathbb{R}^3$}

%Considera el subconjunto $A=\{(x,y,z)\in\mathbb{R}^3: x^2+y^2-4z^2=1\}$ de $\mathbb{R}^3$.

%\begin{enumerate}[label=(\alph*)]
%    \item Describe la intersección de $A$ con el plano $\{z=k\}$ en función de $k$.
%    \item Ídem con las intersecciones con los planos $\{x=k\}$ e $\{y=k\}$ en función de $k$.
%    \item Deduzca la forma aproximada que tiene el conjunto $A$.
%\end{enumerate}

%\end{document}

\newpage 
\section{Pr\'actico 1}

\begin{enumerate}

\item Representar gr\'aficamente los siguientes conjuntos de $\R^2$. Cuando sea posible, reconocer geom\'etricamente el conjunto.

\begin{enumerate}

\item $A=\{(x,y)\in\R^2:\ y=\sen(x),\ \frac{-\pi}{2}\leq x \leq \frac{\pi}{2}\}$.

\item $A=\{(x,y)\in\R^2:\ -1\leq x \leq y^2,\ 2\leq y < 3\}$.

\item $A=\{(x,y)\in\R^2:\ -\sqrt{1-x^2}\leq y \leq \sqrt{1-x^2},\ -1\leq x \leq 1\}$.

\item $A=\{(x,y)\in\R^2:\ 3x+2y=-1\}$.

\item $A=\{(x,y)\in\R^2:\ x-5y=-1,\ -2\leq x < 2\}$.

\item $A=\{(x,y)\in\R^2:\ x^2+y^2-2x-4y+5=3\}$.

\item $A=\{(x,y)\in\R^2:\ x=y^2-2y+1,\ -1\leq x \leq 1\}$.

\item $A=\{(x,y)\in\R^2:\ 5x^2+2y^2=2\}$.

\item $A=\{(x,y)\in\R^2:\ -3x^2+y^2=1\}$.

\item $B((1,-1),\frac{1}{2})$.

\end{enumerate}

\item Dados los vectores $u=(1,1),\ v=(2,1)$ de $\R^2$ hallar

$\proba{\parallel u + v \parallel} \proba{d(u,v)}
\proba{\left<u,v\right>} \proba{\mbox{ el \'angulo entre } u \mbox{ y } v}$\\
$\proba{\mbox{ la recta } r \mbox{ paralela a }u\mbox{ por }(0,1)}  \proba{\mbox{ la recta } s \mbox{ perpendicular a }v\mbox{ por }(0,0)}$\\
$\proba{r\cap s}.$


\item Se considera la funci\'on 
$$f(x,y)=\left\{ 
\begin{array}{ccc}
\frac{2 xy^2}{x^{2}+y^{2}} & \text{si} & \left( x,y\right) \neq \left(
0,0\right)  \\ 
0 & \text{si} & \left( x,y\right) =\left( 0,0\right) 
\end{array}%
\right. $$
\begin{enumerate}
\item Probar que $f$ es diferenciable en $(-1,1)$ y calcular el gradiente en este punto.
\item Calcular las derivadas parciales de $f$ en $(0,0)$. Demostrar que $f$ no es diferenciable en $(0,0)$.
\item Determinar si $f$ es continua en los puntos $(0,0)$  y $(-1,1)$.
\end{enumerate}

\end{enumerate}

\begin{enumerate}

\item 
\begin{enumerate}

\item $A=\{(x,y)\in\R^2:\ y=\sen(x),\ \frac{-\pi}{2}\leq x \leq \frac{\pi}{2}\}$.

\item $A=\{(x,y)\in\R^2:\ -1\leq x \leq y^2,\ 2\leq y < 3\}$.

\item $A=\{(x,y)\in\R^2:\ -\sqrt{1-x^2}\leq y \leq \sqrt{1-x^2},\ -1\leq x \leq 1\}$.

\item $A=\{(x,y)\in\R^2:\ 3x+2y=-1\}$.

\item $A=\{(x,y)\in\R^2:\ x-5y=-1,\ -2\leq x < 2\}$.

\item $A=\{(x,y)\in\R^2:\ x^2+y^2-2x-4y+5=3\}$.

\item $A=\{(x,y)\in\R^2:\ x=y^2-2y+1,\ -1\leq x \leq 1\}$.

\item $A=\{(x,y)\in\R^2:\ 5x^2+2y^2=2\}$.

\item $A=\{(x,y)\in\R^2:\ -3x^2+y^2=1\}$.

\item $B((1,-1),\frac{1}{2})$.

\end{enumerate}

\item Dados los vectores $u=(1,1),\ v=(2,1)$ de $\R^2$ hallar

%$\proba{\parallel u + v \parallel} \proba{d(u,v)}
%\proba{\left<u,v\right>} \proba{\mbox{ el angulo entre } u \mbox{ y } v}$\\
%$\proba{\mbox{ la recta } r \mbox{ paralela a }u\mbox{ por }(0,1)}  \proba{\mbox{ la recta } s \mbox{ perpendicular a }v\mbox{ por }(0,0)}$\\
%$\proba{r\cap s}.$

\cuento

\vspace{0,4cm}

\item Muestre que en general dado un vector $v$ del plano $\R^2$, el conjunto $K_v$ de puntos cuya proyecci\'on sobre $v$ es cero es una recta que pasa por el origen.
Encuentre dicha recta para $v=(2,3)$. Muestre que dados dos vectores $v$ y $v'$, dichas rectas coinciden si y solo si $v$ y $v'$ son colineales. 



\item Represente gr\'aficamente los siguientes conjuntos del espacio $\R^3$. Cuando sea posible, reconocer geom\'etricamente el conjunto. 

\begin{enumerate}

\item $A=\{(x,y,z)\in\R^3:\ 3x+2y-z=1\}$.

\item $A=\{(x,y,z)\in\R^3:\ x+y-z=0,\ -2x-2y+2z=1\}$.

\item $A=\{(x,y,z)\in\R^3:\ x^2+y^2=1\}$.

\item $A=\{(x,y,z)\in\R^3:\ 3x^2+y^2=2,\ -1\leq z \leq 2\}$.

\item $A=\{(x,y,z)\in\R^3:\ x^2-y^2=1,\ z=1\}$.

\item $A=\{(x,y,z)\in\R^3:\ x^2+y^2+(z-2)^2=10\}$.

\item $A=\{(x,y,z)\in\R^3:\ x^2+y^2+z^2=1, (x-1)^2+y^2=1\}$.

\item $A=\{(x,y,z)\in\R^3:\ 2x^2+y^2+z^2=2, x+y=1\}$.

\item $A=\{(x,y,z)\in\R^3:\ x^2+y^2-4z^2=1\}$.

\end{enumerate}

\item Dados los vectores $u=(1,1,1),\ v=(2,1,0)$ de $\R^3$ hallar

$\proba{\parallel u + v \parallel} \proba{d(u,v)}
\proba{\left<u,v\right>} \proba{\mbox{ el \'angulo entre } u \mbox{ y } v}$\\
$\proba{\mbox{ la recta } r \mbox{ paralela a }u\mbox{ por }(0,0,1)}  \proba{\mbox{ el plano } \Pi \mbox{ perpendicular a }v\mbox{ por }(0,0,0)}$\\
$\proba{r\cap s}.$

\cuento

\vspace{0,4cm}

\item Muestre que en general dado un vector $v$ del espacio $\R^3$, el conjunto $K_v$ de puntos cuya proyecci\'on sobre $v$ es cero es un plano que pasa por el origen.
Encuentre dicho plano para $v=(1,-1,0)$. Muestre que dados dos vectores $v$ y $v'$, dichos planos coinciden si y solo si $v$ y $v'$ son colineales. 

\item Considere el subconjunto $A=\{(x,y,z)\in\R^3:\ x^2+y^2-4z^2=1\}$ de $\R^3$.

  \begin{enumerate}
  \item Describa la intersecci\'on de $A$ con el plano $\{z=k\}$ en funci\'on de $k$ (hay que observar que el tipo de conjunto que se obtiene puede cambiar para ciertos valores del parámetro $k$).
\item Ídem con las intersecciones con los planos $\{x=k\}$ e $\{y=k\}$ en funci\'on de $k$.    
\item Deduzca la forma aproximada que tiene el conjunto $A$.
\end{enumerate}

\end{enumerate}

%\section{Ejemplos(practico)}

\chapter{Funciones escalares I de V}

\section{Temario} 

\begin{itemize}
    \item Dominio, imágen, gráfico y conjuntos de nivel.
    \item Límites y continuidad.
%    \item Derivadas (derivadas parciales/direccionales, diferenciabilidad, gradiente y espacio tangente).
%    \item Derivadas de orden superior (Taylor, criterio de la Hessiana).
%    \item Teorema de la Función Implícita.
%    \item Extremos condicionados (Multiplicadores de Lagrange).
\end{itemize}

\section{Notas}

%***
%\documentclass{article}
%\usepackage{amsmath, amsthm, amssymb}
%\usepackage[utf8]{inputenc}
%\usepackage[T1]{fontenc}
%\usepackage[spanish]{babel}
%\usepackage{graphicx}

%\newtheorem{definition}{Definición}
%\newtheorem{example}{Ejemplo}

%\title{Curso sobre funciones de dos variables}
%\author{}
%\date{}

%\begin{document}

%\maketitle

%\section{Introducción}
En este curso, estudiaremos las funciones de dos variables. Estas funciones son de la forma \( f(x, y) \) y juegan un papel importante en matemáticas y en aplicaciones, particularmente en física, economía e ingeniería.

\section{Dominio}
\begin{definition}
El dominio de una función de dos variables \( f(x, y) \) es el conjunto de pares \((x, y) \in \mathbb{R}^2\) para los cuales la función está definida. Se puede notar $\mathrm{Dom}(f)$.
\end{definition}

El dominio de una función de dos variables es un subconjunto de \(\mathbb{R}^2\). Puede estar definido por restricciones sobre \(x\) y \(y\), tales como desigualdades o ecuaciones.

\begin{example}
Consideremos la función \( f(x, y) = \sqrt{x^2 + y^2 - 1} \). El dominio de \( f \) es el conjunto de puntos \((x, y)\) tales que \(x^2 + y^2  \geq 1\), es decir, el exterior del círculo unitario abierto.
\end{example}

\begin{example}
Consideremos la función \( f(x, y) = \frac{1}{x - y} \). El dominio de \( f \) es el conjunto de puntos \((x, y)\) tales que \(x \neq y\), es decir, todos los puntos excepto aquellos en la línea \(y = x\).
\end{example}

\section{Imagen}
\begin{definition}
La imagen de una función de dos variables \( f(x, y) \) es el conjunto de valores que la función puede tomar, es decir,
\[ \{ f(x, y) \mid (x, y) \in \mathrm{Dom}(f) \}. \]
\end{definition}

La imagen es un subconjunto de \(\mathbb{R}\) que representa todos los valores posibles que la función puede devolver.

\begin{example}
Consideremos la función \( f(x, y) = x^2 + y^2 \). La imagen de \( f \) es \([0, +\infty)\) porque \(x^2 + y^2\) es siempre no negativo y puede tomar cualquier valor positivo.
\end{example}

\begin{example}
Consideremos la función \( f(x, y) = \sin(x) + \cos(y) \). La imagen de \( f \) es \([-2, 2]\) porque los valores máximo y mínimo de \(\sin(x)\) y \(\cos(y)\) son 1 y -1, respectivamente.
\end{example}

\section{Gráfico}
\begin{definition}
El gráfico de una función de dos variables \( z = f(x, y) \) es el conjunto de puntos \((x, y, z) \in \mathbb{R}^3\) tales que \( z = f(x, y) \).
\end{definition}

El gráfico es una superficie en el espacio tridimensional. Para cada punto \((x, y)\) en el dominio de \(f\), hay un punto correspondiente \((x, y, z)\) en el gráfico donde \(z\) es el valor de la función en \((x, y)\).

\begin{example}
Consideremos la función \( f(x, y) = x^2 + y^2 \). El gráfico de \( f \) es un paraboloide de revolución.
\end{example}

\begin{example}
Consideremos la función \( f(x, y) = \sqrt{1 - x^2 - y^2} \). El gráfico de \( f \) es un hemisferio superior de radio 1, centrado en el origen.
\end{example}

\section{Conjuntos de nivel}
\begin{definition}
Para una función \( f(x, y) \) y una constante \( c \), el conjunto de nivel \( c \) es el conjunto de puntos \((x, y)\) en el dominio de \(f\) tales que \( f(x, y) = c \).
\end{definition}

Los conjuntos de nivel son curvas en el plano \(xy\) donde la función toma un valor constante. Son útiles para visualizar el comportamiento de la función.

\begin{example}
Consideremos la función \( f(x, y) = x^2 + y^2 \). Los conjuntos de nivel para \( c > 0 \) son círculos concéntricos centrados en el origen, con un radio de \(\sqrt{c}\).
\end{example}

\begin{example}
Consideremos la función \( f(x, y) = xy \). Los conjuntos de nivel para \( c \neq 0 \) son hipérbolas.
\end{example}

\section{Límite}
\begin{definition}
El límite de \( f(x, y) \) cuando \((x, y)\) se aproxima a \((a, b)\) es \( L \), denotado por
\[ \lim_{(x, y) \to (a, b)} f(x, y) = L, \]
si para todo \(\epsilon > 0\), existe un \(\delta > 0\) tal que para todo \((x, y)\) en el dominio de \(f\) con \(0 < \sqrt{(x - a)^2 + (y - b)^2} < \delta\), se tiene \(|f(x, y) - L| < \epsilon\).
\end{definition}

El límite generaliza la noción de límite para funciones de una variable. Para funciones de dos variables, el límite debe ser el mismo sin importar la dirección de aproximación al punto \((a, b)\).

\begin{example}
Mostremos que \[\lim_{(x, y) \to (0, 0)} \frac{xy}{x^2 + y^2} = 0\] a lo largo de cualquier línea que pasa por el origen. Sin embargo, el límite no existe porque depende de la dirección de aproximación. Por ejemplo, si nos acercamos a lo largo de la línea \(y = x\), el límite es \( \frac{1}{2} \), mientras que si nos acercamos a lo largo de \(y = 2x\), el límite es \( \frac{2}{5} \).
\end{example}

\begin{example}
Calculemos \[\lim_{(x, y) \to (0, 0)} \frac{x^2 y}{x^2 + y^2}.\] Usando coordenadas polares %o mostrando que el límite a lo largo de cualquier línea que pasa por el origen es 0, y que la función está acotada, 
podemos concluir que el límite es 0.
\end{example}

\section{Continuidad}
\begin{definition}
Una función \( f(x, y) \) es continua en un punto \((a, b)\) de su dominio si
\[ \lim_{(x, y) \to (a, b)} f(x, y) = f(a, b). \]
Es continua en su dominio si es continua en cada punto del mismo.
\end{definition}

La continuidad para funciones de dos variables significa que no hay "saltos" o "agujeros" en la superficie representada por el gráfico de la función.

\begin{example}
La función \( f(x, y) = x^2 + y^2 \) es continua en todo \(\mathbb{R}^2\) porque es una función polinómica.
\end{example}

\begin{example}
Consideremos la función
\[ f(x, y) = \begin{cases}
\frac{xy}{x^2 + y^2} & \text{si } (x, y) \neq (0, 0) \\
0 & \text{si } (x, y) = (0, 0)
\end{cases} \]
Esta función no es continua en \((0, 0)\) porque el límite no existe en ese punto.
\end{example}

\begin{example}
    Mas sorprendente aun:
    \[ f(x, y) = \begin{cases}
\frac{x^2y}{x^4 + y^2} & \text{si } (x, y) \neq (0, 0) \\
0 & \text{si } (x, y) = (0, 0)
\end{cases} \]
En este caso, encontramos limites distintos $(0,0)$ en cero si tomamos un camino $y=\alpha x$ y otro camino $y = \alpha x^2$.
\end{example}

%\end{document}

%***

%\noindent Veamos unas primeras nociones sobre funciones de dos variables.

%\subsection{Introducción a las funciones de dos variables}

%Una función de dos variables es una función que asigna a cada par ordenado \((x, y)\) en un dominio un único valor real \(f(x, y)\). Por ejemplo, la función \(f(x, y) = x^2 + y^2\) es una función de dos variables.

%\documentclass{article}
%\usepackage{tikz}
%\usepackage{amsmath}

%\begin{document}

%\section*{Curso: Funciones de Dos Variables}

%\subsection*{Introducción}

%Las funciones de dos variables son fundamentales en el estudio del cálculo multivariable. Estas funciones asignan un valor único a cada par de variables independientes \((x, y)\). En este curso, exploraremos los conceptos básicos de dominio, imagen, gráfico y conjuntos de nivel.

%\subsection*{Dominio de una Función de Dos Variables}

%\textbf{Definición:}

%El dominio de una función de dos variables \( f(x, y) \) es el conjunto de todos los pares ordenados \((x, y)\) para los cuales la función está definida.

%\textbf{Ejemplo 1:}

%Consideremos la función \( f(x, y) = \sqrt{x^2 + y^2} \). El dominio de esta función es todo \(\mathbb{R}^2\), ya que \(x^2 + y^2\) es siempre no negativo, y la raíz cuadrada está definida para todos los números no negativos.

%\textbf{Ejemplo 2:}

%Para la función \( f(x, y) = \frac{1}{x - y} \), el dominio es el conjunto de todos los pares \((x, y)\) tales que \(x \neq y\), ya que el denominador no puede ser cero.

%\subsection*{Imagen de una Función de Dos Variables}

%\textbf{Definición:}

%La imagen de una función de dos variables \( f(x, y) \) es el conjunto de todos los valores posibles que \( f(x, y) \) puede tomar.

%\textbf{Ejemplo 1:}

%Para la función \( f(x, y) = x^2 + y^2 \), la imagen es el conjunto de todos los números reales no negativos, \([0, \infty)\), ya que \(x^2 + y^2\) siempre da un resultado no negativo.

%\textbf{Ejemplo 2:}

%Consideremos la función \( f(x, y) = \sin(x + y) \). La imagen de esta función es el intervalo \([-1, 1]\), ya que la función seno oscila entre -1 y 1 para cualquier valor real de su argumento.

%\subsection*{Gráfico de una Función de Dos Variables}

%\textbf{Definición:}

%El gráfico de una función de dos variables \( f(x, y) \) es el conjunto de todos los puntos \((x, y, z)\) en \(\mathbb{R}^3\) tales que \( z = f(x, y) \).

%\textbf{Ejemplo 1:}

%El gráfico de la función \( f(x, y) = x^2 + y^2 \) es un paraboloide de revolución. Este gráfico se puede visualizar como una superficie en forma de cuenco que se extiende hacia arriba desde el origen.

%\begin{tikzpicture}[scale=1.5]
%    \draw[->] (-2,0,0) -- (2,0,0) node[right]{$x$};
%    \draw[->] (0,-2,0) -- (0,2,0) node[left]{$y$};
%    \draw[->] (0,0,0) -- (0,0,4) node[above]{$z$};
%    \draw[domain=-1.5:1.5, variable=\x, smooth] plot ({\x}, {0}, {\x*\x});
%    \draw[domain=-1.5:1.5, variable=\y, smooth] plot ({0}, {\y}, {\y*\y});
%\end{tikzpicture}

%\textbf{Ejemplo 2:}

%Para la función \( f(x, y) = \sqrt{1 - x^2 - y^2} \), el gráfico es el hemisferio superior de una esfera de radio 1 centrada en el origen.

%\subsection*{Conjuntos de Nivel}

%\textbf{Definición:}

%Un conjunto de nivel de una función de dos variables \( f(x, y) \) es el conjunto de todos los puntos \((x, y)\) en el dominio de \( f \) tales que \( f(x, y) = c \), donde \( c \) es una constante.

%\textbf{Ejemplo 1:}

%Para la función \( f(x, y) = x^2 + y^2 \), los conjuntos de nivel son círculos concéntricos centrados en el origen. Por ejemplo, el conjunto de nivel para \( c = 1 \) es el círculo \( x^2 + y^2 = 1 \).

%\begin{tikzpicture}[scale=1.5]
%    \draw[->] (-2,0) -- (2,0) node[right]{$x$};
%    \draw[->] (0,-2) -- (0,2) node[above]{$y$};
%    \draw (0,0) circle (1);
%    \node at (0.5,0.5) {$x^2 + y^2 = 1$};
%\end{tikzpicture}

%\textbf{Ejemplo 2:}

%Consideremos la función \( f(x, y) = x + y \). Los conjuntos de nivel de esta función son líneas rectas de la forma \( x + y = c \). Por ejemplo, el conjunto de nivel para \( c = 2 \) es la línea \( x + y = 2 \).

%\begin{tikzpicture}[scale=1.5]
%    \draw[->] (-2,0) -- (2,0) node[right]{$x$};
%    \draw[->] (0,-2) -- (0,2) node[above]{$y$};
%    \draw[domain=-1:1, variable=\x, smooth] plot ({\x}, {2 - \x});
%    \node at (0.5,1) {$x + y = 2$};
%\end{tikzpicture}

%\subsection*{Conclusión}

%Estos conceptos son fundamentales para entender y trabajar con funciones de dos variables. El dominio y la imagen nos ayudan a entender dónde está definida la función y qué valores puede tomar, mientras que el gráfico y los conjuntos de nivel nos proporcionan una visualización útil de la función.

%\end{document}

%\subsection{Dominio, im\'agen, gr\'afico, conjuntos de nivel}

%\begin{definition}
%    El dominio \dots % MISTRAL
%\end{definition}

%\begin{example}
%Para la siguiente funciones hallamos el dominio m\'as grande posible,
%representar curvas de nivel e intentar bosquejar el gr\'afico.
%$${f(x,y)={\sqrt{x+y}}}.$$
%\end{example}

%\subsection{Límites} %y continuidad}

%\begin{example}
%Hallamos el l\'{\i}mite (si existe), o probamos que no existe:
%$\underset{\left( x,y\right) \rightarrow \left( 0,0\right) }{\lim }\frac{%
%y^{2}}{x^{2}+y^{2}}$.
%\end{example}

%\subsection{Continuidad de funciones de dos variables}

%Una función \(f(x, y)\) es continua en un punto \((a, b)\) si se satisfacen las siguientes condiciones:
%\begin{enumerate}
%    \item \(f(a, b)\) está definida.
%    \item \(\lim_{(x, y) \to (a, b)} f(x, y)\) existe.
%    \item \(\lim_{(x, y) \to (a, b)} f(x, y) = f(a, b)\).
%\end{enumerate}

%\subsection{Ejemplo de función continua}
%Consideremos la función \(f(x, y) = x^2 + y^2\). Esta función es continua en todo \(\mathbb{R}^2\) porque es una combinación de funciones polinómicas, que son continuas.

%\subsection{Contraejemplo de función no continua}
%Consideremos la función definida por partes:
%\[ f(x, y) = \begin{cases}
%\frac{xy}{x^2 + y^2} & \text{si } (x, y) \neq (0, 0), \\
%0 & \text{si } (x, y) = (0, 0).
%\end{cases} \]

%Esta función es discontinua en \((0, 0)\) porque el límite no existe. Para ver esto, consideremos dos caminos diferentes hacia \((0, 0)\):
%\begin{itemize}
%    \item Si nos acercamos a lo largo del eje \(x\) (es decir, \(y = 0\)), entonces \(f(x, 0) = 0\) para todo \(x\).
%    \item Si nos acercamos a lo largo de la recta \(y = x\), entonces \(f(x, x) = \frac{x^2}{2x^2} = \frac{1}{2}\).
%\end{itemize}

%Como el límite depende del camino tomado, el límite no existe y la función es discontinua en \((0, 0)\).

%{Curso breve sobre derivadas parciales, diferencial y derivada direccional}

%\subsection{Continuidad (bis)}

%\begin{example}
%Veamos que la función siguiente sea continua
%$f(x,y)=\left\{ 
%\begin{array}{ccc}
%\frac{xy}{\sqrt{x^{2}+y^{2}}} & \text{si} & (x,y)\neq \left(
%0,0\right)  \\ 
%0 & \text{si} & \left( x,y\right) =\left( 0,0\right) 
%\end{array}%
%\right. $.
%\end{example}

%\begin{example}
%Funciones no continuas:
%\[ 
%f(x,y) = |x+y|, \quad
%f(x,y) = \frac{xy}{x^2 y^2}, \quad  
%f(x,y) = \frac{x y^2}{x^2 y^4}.
%\]
%\end{example}

%(Método de caminos / caracterización sequencial de la continuidad)
%(coordenadas polares)

\newpage

\section{Practico 2}

\begin{enumerate}

\item Para las siguientes funciones hallar el dominio m\'as grande posible,
representar curvas de nivel e intentar bosquejar el gr\'afico.\\

$\proba{f(x,y)=e^{\sqrt{x-y}}} \proba{\ln \left(\frac{1}{x^2-y^2+1}\right)} \proba{\sen(x^2+y^2)}$\\
$\proba{f(x,y)=\cos(y)+x} \proba{f(x,y)=\sqrt{x^2+y^2}\mbox{ si }x^2+y^2\geq 1,\ \frac{1}{x^2+y^2}\mbox{ si }x^2+y^2\leq 1}$\\
$\proba{f(x,y)=\arctg(\ln(x)+2y)} \proba{f(x,y)=e^{\cos(x)+y}} $

\cuento

\vspace{0,4cm}

\item Para los siguietnes casos hallar dominio m\'aximo, conjunto im\'agen y superficies de nivel:


a) $f(x,y,z)=2x+y-3z$, \ \ b) $f(x,y,z)=x^2-y^2+z^2$, \ \ c) $f(x,y,z)=e^{x+y}-z$ \\
d) $f(x,y,z)= x\cos(y)$, \ \ e) $f(x_1,\dots,x_n)=x_1+\dots+x_n$ \ \ f) $f(x_1,\dots,x_n)=x_1^2+\dots+x_n^2$.  




\item Hallar el l\'{\i}mite (si existe), o probar que no existe, en los siguientes
casos:

a) $\underset{\left( x,y\right) \rightarrow \left( 0,0\right) }{\lim }\frac{%
x^{2}}{x^{2}+y^{2}}$ b) $\underset{\left(
x,y\right) \rightarrow \left( 0,0\right) }{\lim }\frac{xy}{\sqrt{x^{2}+y^{2}}%
}$ c) $\underset{\left( x,y\right) \rightarrow
\left( 0,0\right) }{\lim }\frac{(x+y)^{2}}{x^{2}+y^{2}}$ d) $\underset{%
\left( x,y\right) \rightarrow \left( 0,0\right) }{\lim }\frac{\left(
x+y\right) ^{2}}{\sqrt{x^{2}+y^{2}}}$ e) $\underset{\left( x,y\right)
\rightarrow \left( 0,0\right) }{\lim }\frac{y^{3}+x^{2}y}{x^{2}+y^{2}}$ \ f) $\underset{\left( x,y\right) \rightarrow \left(
0,0\right) }{\lim }(x^{2}+y^{2})\ln (x^{2}+y^{2})$\, g) $\underset{\left( x,y\right)
\rightarrow \left( 0,0\right) }{\lim }\frac{\sin \left( x^{3}+y^{2}\right) }{%
x^{2}+y^{2}}$ h) $\underset{\left( x,y\right) \rightarrow \left( 0,0\right) 
}{\lim }\frac{1-\cos \left( x+y^{{}}\right) }{\sqrt{x^{2}+y^{2}}}$ i) $%
\underset{\left( x,y\right) \rightarrow \left( 0,0\right) }{\lim }\frac{%
e^{xy}-1}{\sqrt{x^{2}+y^{2}}}$

j) $\underset{\left( x,y\right) \rightarrow \left( 0,0\right) }{\lim }\frac{%
\ln \left( 1+x+y^{{}}\right) }{x+y}$

\medskip
\item Investigar si la función definida en cada caso es continua

a) $f(x,y)=\left\{ 
\begin{array}{ccc}
\frac{x^{2}y^{3}}{x^{2}+y^{2}} & \text{si} & (x,y)\neq \left( 0,0\right)  \\ 
1 & \text{si} & \left( x,y\right) =\left( 0,0\right) 
\end{array}%
\right. 
 $ b) $f(x,y)=\left\{ 
\begin{array}{ccc}
\frac{xy}{x^{2}+xy+y^{2}} & \text{si} & (x,y)\neq \left( 0,0\right)  \\ 
0 & \text{si} & \left( x,y\right) =\left( 0,0\right) 
\end{array}%
\right. $ \\d) $f(x,y)=\left\{ 
\begin{array}{ccc}
\frac{e^{x^{2}y}-1}{\sqrt{x^{2}+y^{2}}} & \text{si} & (x,y)\neq \left(
0,0\right)  \\ 
0 & \text{si} & \left( x,y\right) =\left( 0,0\right) 
\end{array}%
\right. $.




\item Se considera el potencial $V(x) = kx^2$. La energ\'ia total de una part\'icula unidimensional que experimenta este potencial, depende solamente de la posici\'on y de la velocidad, y es $E(x,v) = \frac{1}{2}m v^2 + V(x)$. Determinar la im\'agen de la funci\'on $E$ y dibujar las curvas de nivel.



\item  Considerar el potencial 
$$V(r) = \frac{L^2}{2mr^2} - \frac{Gm}{r},$$
donde $m,L$ y $G$ son constantes.

Determinar la im\'agen de la funci\'on $E(r,\dot{r}) = \frac{1}{2}m\dot{r}^2 + V(r)$ y bosquejar las curvas de nivel.




\end{enumerate}

\chapter{Funciones escalares II de V}

\section{Temario}

\begin{itemize}
%    \item Dominio, imágen, gráfico y conjuntos de nivel.
%    \item Límites y continuidad.
    \item Derivadas: 
        \begin{itemize}
            \item derivadas parciales,
            \item derivadas direccionales,
            \item espacio tangente.
        \end{itemize}
        %(/direccionales, diferenciabilidad, gradiente y espacio tangente).
 %   \item Derivadas de orden superior (Taylor, criterio de la Hessiana).
 %   \item Teorema de la Función Implícita.
 %   \item Extremos condicionados (Multiplicadores de Lagrange).
\end{itemize}

\section{Notas}

%\documentclass[12pt,a4paper]{article}
%\usepackage[utf8]{inputenc}
%\usepackage[spanish]{babel}
%\usepackage{amsmath}
%\usepackage{amsthm}
%\usepackage{amsfonts}
%\usepackage{amssymb}
%\usepackage{enumitem}
%\usepackage{geometry}
%\geometry{left=2cm,right=2cm,top=2cm,bottom=2cm}

%\newtheorem{definition}{Definición}[section]
%\newtheorem{example}{Ejemplo}[section]

%\title{Notas de Clase: Derivadas en Varias Variables}
%\author{Franz Chouly}
%\date{\today}

%\begin{document}

%\maketitle

\subsection{Derivadas Parciales}

%\begin{definition}
%Sea \( f: \mathbb{R}^n \to \mathbb{R} \) una función y \( \mathbf{a} = (a_1, a_2, \ldots, a_n) \) un punto en el dominio de \( f \). La \textbf{derivada parcial} de \( f \) con respecto a la variable \( x_i \) en \( \mathbf{a} \) se define como:
%\[
%\frac{\partial f}{\partial x_i}(\mathbf{a}) = \lim_{h \to 0} \frac{f(a_1, \ldots, a_i + h, \ldots, a_n) - f(a_1, \ldots, a_n)}{h},
%\]
%si este límite existe.
%\end{definition}

\begin{definition}
%    
%\subsection{Derivadas parciales de funciones de dos variables}
Dada una función de dos variables \(f(x, y)\), las derivadas parciales con respecto a \(x\) y \(y\) se definen como:
\[ f_x(x, y) = \frac{\partial f}{\partial x}(x,y) = \lim_{h \to 0} \frac{f(x + h, y) - f(x, y)}{h}, \]
\[ f_y(x, y) = \frac{\partial f}{\partial y}(x,y)= \lim_{h \to 0} \frac{f(x, y + h) - f(x, y)}{h}. \]
\end{definition}


\begin{example}
Consideremos la función \( f(x, y) = x^2 y + \sin(y) \).
\begin{itemize}
    \item Derivada parcial respecto a \( x \): \( \frac{\partial f}{\partial x} = 2xy \).
    \item Derivada parcial respecto a \( y \): \( \frac{\partial f}{\partial y} = x^2 + \cos(y) \).
\end{itemize}
\end{example}

\begin{example}
Para la función \( g(x, y, z) = e^{xyz} \):
\begin{itemize}
    \item \( \frac{\partial g}{\partial x} = yz e^{xyz} \),
    \item \( \frac{\partial g}{\partial y} = xz e^{xyz} \),
    \item \( \frac{\partial g}{\partial z} = xy e^{xyz} \).
\end{itemize}
\end{example}

\begin{example}
La funcion
    $$ %\underset{\left( x,y\right) \rightarrow \left( 0,0\right) }{\lim }
    f(x,y)=\frac{xy}{x^{2}+y^{2}}$$
tiene derivadas parciales $\frac{\partial f}{\partial x}$ y $\frac{\partial f}{\partial y}$ nulas en $(0,0)$, a pesar de que no es continua en $(0,0)$.
\end{example}

\begin{example}
    La funci\'on cono 
    \[
    f(x,y) = \sqrt{x^2+y^2}
    \]
    es continua en $(0,0)$ pero no admite derivadas parciales.
\end{example}

\begin{example}
    La funci\'on  
    \[
    f(x,y) = |x|
    \]
    es continua, y en $(0,0)$ admite una derivada parcial en $y$, pero no en $x$.
\end{example}

%OTRO EJEMPLO {Derivadas parciales}

%$f(x,y)=\left\{ 
%\begin{array}{ccc}
%xy & \text{si} & x\leq y \\ 
%x^{2} & \text{si} & x>y%
%\end{array}%
%\right. $ en $\left(1,1\right)$ y $\left( 1,2\right) $


\subsection{Derivadas Direccionales}

\begin{definition}
Sea \( f: \mathbb{R}^n \to \mathbb{R} \) y \( \mathbf{v} \) un vector unitario en \( \mathbb{R}^n \). La \textbf{derivada direccional} de \( f \) en la dirección de \( \mathbf{v} \) en el punto \( \mathbf{a} \) es:
\[
D_{\mathbf{v}} f(\mathbf{a}) = \lim_{h \to 0} \frac{f(\mathbf{a} + h\mathbf{v}) - f(\mathbf{a})}{h}.
\]
\end{definition}

\begin{example}
Sea \( f(x, y) = x^2 + y^2 \) y \( \mathbf{v} = \left( \frac{1}{\sqrt{2}}, \frac{1}{\sqrt{2}} \right) \). La derivada direccional en \( (1, 1) \) es:
\[
D_{\mathbf{v}} f(1, 1) = %\nabla f(1,1) \cdot \mathbf{v} = (2, 2) \cdot \left( \frac{1}{\sqrt{2}}, \frac{1}{\sqrt{2}} \right) = \frac{4}{\sqrt{2}}
2\sqrt{2}.
\]
\end{example}

\begin{example}
Para \( g(x, y) = xy \) y \( \mathbf{v} = (0, 1) \), la derivada direccional en \( (2, 3) \) es:
\[
D_{\mathbf{v}} g(2, 3) = \nabla g(2,3) \cdot \mathbf{v} = (3, 2) \cdot (0, 1) = 2.
\]
\end{example}

\begin{example}
    La funci\'on cono 
    \[
    f(x,y) = \sqrt{x^2+y^2}
    \]
     no admite ninguna derivada direccional en $(0,0)$. 
\end{example}

\begin{example}
    La funci\'on  
    \[
    f(x,y) = |2x+y|
    \]
    en $(0,0)$ admite una derivada direccional en la direcci\'on de $v$ si y solo si $v=(1,2)$.
%
%*** Detallar aqui ***
%    
\end{example}

\begin{example}
La funcion $f(x,y)=x+y$
admite una derivada direccional $\partial_v f(x,y)$ para todo $v=(v_1,v_2)\in\R^2$.    
\end{example}
 
\subsection{Plano Tangente}

%\begin{definition}
%Sea \( S \) una superficie en \( \mathbb{R}^3 \) definida por \( F(x, y, z) = 0 \). El \textbf{espacio tangente} a \( S \) en un punto \( \mathbf{p} = (x_0, y_0, z_0) \) es el plano dado por:
%\[
%\nabla F(\mathbf{p}) \cdot (x - x_0, y - y_0, z - z_0) = 0.
%\]
%\end{definition}

%\begin{example}
%Consideremos la esfera \( %F(x, y, z) = x^2 + y^2 + z^2 - 1 = 0 \). El espacio tangente en \( (1, 0, 0) \) es:
%\[
%(2, 0, 0) \cdot (x - 1, y, z) = 0 \implies x = 1.
%\]
%\end{example}

%\subsection{Nota SEB ***}

\begin{definition}
El plano tangente a $f(x,y)$ en el punto $(a,b)$ es dado por:
\begin{equation}
    z = f(a,b) + f_x(a,b) (x-a) + f_y(a,b) (y-b).
\end{equation}    
\end{definition}
\noindent 
Se deduce facilmente un vector normal al plano tangente
\[
n = ( -f_x , -f_y , 1).
\]
%
%**** Detaillar aqui ***

\begin{example}
El plano tangente de $z=5-x^2-2y^2$ en $(1,1,2)$ es dado por:
\[
z=2-2(x-1)-4(y-1).
\]
\end{example}

\begin{example}
Para el paraboloide \( %F(x, y, z) = 
z - x^2 - y^2 = 0 \), el plano tangente en \( (1, 1, 2) \) es:
\[
%(-2, -2, 1) \cdot (x - 1, y - 1, z - 2) = 0 \implies 
-2(x-1) - 2(y-1) + (z-2) = 0.
\]
\end{example}

%VER PRACTICO 3 , etc : EDPs

%\end{document}
%***


%\subsection{Espacio tangente}

%$f(x,y)=x^2+y^2$ en $(0,0,0)$ y $(1,-1,2)$.


%Ejemplo cono y $|x+y|$ etc
\newpage

\section{Práctico 3}

\begin{enumerate}


\item Calcular, en caso de que existan, las derivadas parciales de las
  siguientes funciones en los puntos que se indican:
  \begin{itemize}
  \item[a)] $f(x,y)=\left\{ 
\begin{array}{ccc}
\frac{x^{2}y^{3}}{x^{2}+y^{2}} & \text{si} & \left( x,y\right) \neq \left(
0,0\right)  \\ 
0 & \text{si} & \left( x,y\right) =\left( 0,0\right) 
\end{array}%
\right. $ en $\left( 0,0\right) $ y $\left( 1,2\right) $

\item[b)] $f(x,y)=\left\{ 
\begin{array}{ccc}
xy & \text{si} & x\neq 1 \\ 
y^{2} & \text{si} & x=1%
\end{array}%
\right. $ en $\left( 1,1\right) $ y $\left( 2,3\right) 
$
% \item[c)] $f(x,y)=\left\{ 
% \begin{array}{ccc}
% 2x-3y & \text{si} & xy\neq 0 \\ 
% 3x-4y & \text{si} & xy=0%
% \end{array}%
% \right. $ en $\left( 0,0\right) $; $\left( 1,0\right) ;$ $\left( 0,1\right) $
% y $\left( 1,2\right) $
% \item[d)] $f(x,y)=\left\{ 
% \begin{array}{ccc}
% \frac{e^{\left( x-y\right) ^{2}}-1}{\sqrt{x^{2}+y^{2}}} & \text{si} & \left(
% x,y\right) \neq \left( 0,0\right)  \\ 
% 0 & \text{si} & \left( x,y\right) =\left( 0,0\right) 
% \end{array}%
% \right. $ en $\left( 0,0\right) $

\item[c)] $f(x,y)=\left\{ 
\begin{array}{ccc}
\frac{x^{2}y^{3}}{x^{2}+y^{2}} & \text{si} & \left( x,y\right) \neq \left(
0,0\right)  \\ 
0 & \text{si} & \left( x,y\right) =\left( 0,0\right) 
\end{array}%
\right. $ en $\left( 0,0\right) $ y $\left( 1,2\right) $
\item[d)] $f(x,y)=\left\{ 
\begin{array}{ccc}
xy & \text{si} & x\leq y \\ 
x^{2} & \text{si} & x>y%
\end{array}%
\right. $ en $\left(1,1\right)$ y $\left( 1,2\right) $
% \item[g)] $
% f(x,y)=\left\vert x-y^{2}\right\vert $ en $\left( 0,0\right) ;$ $\left(
% 1,1\right) $ y $\left( 1,2\right) .$
\end{itemize}


\item Hallar $\partial_{(1,1)}f(x,y)$ con $f(x,y)=2x+y$.

  \item Si $f(x,y)=ax+by$, donde $(a,b)$ es un vector fijo de $\R^2$,
hallar  $\partial_v f(x,y)$ para todo $v=(v_1,v_2)\in\R^2$. 

\smallskip
\item Analizar en qué puntos de $\R^2$ existe derivada direccional de la siguiente función:
   \[f(x,y)=\left\{\begin{array}{l}
                           0\ \mbox{ si }x\neq y^2\\
                           1\mbox{ si } x=y^2.\end{array}\right.\]

\item En los siguientes casos calcular la $\partial_v f(a)$, en el punto $a\in\R^n$ y vector $v\in\R^n$ indicados. 
\begin{itemize}
\item[a)] $f(x,y)=\sen \left( \sqrt{x^{2}+y^{2}}\right), a=(0,0), v=(v_1,v_2)$.
\item[b)]  $f(x,y)=e^{2x+y}-1, a=(0,0), v=(1,-1)$.
\item[c)]  $f(x,y,z)=\sqrt{x^2+y^2+z^2}, a=(x,y,z), v=(1,0,-1)$.
\end{itemize}


\item En los siguientes casos investigar si existe cada una de las
  derivadas parciales en el origen:
  \begin{itemize}
    \item[a)] $f(x,y,z)=\sen \left( \sqrt{x^{2}+y^{2}+z^2}
      \right) $
    \item[b)] $f(x,y,z)=x\sen \left( \sqrt{x^{2}+y^{2}+z^2}\right) $
      \item[c)] $f(x,y,z)=
        \sqrt[3]{x^{3}+y^{3}+z^3}$
      \end{itemize}
Observar que para $a)$ y $c)$, por un argumento de simetría, basta con hacer el cálculo para una sola de las derivadas parciales. En el caso $b)$ sin embargo, hay diferencias entre lo que ocurre con $\partial f/\partial x(0,0,0)$ y las otras dos derivads parciales (¿puede explicar cuál es la diferencia?) 

\item Calcular la ecuación del plano tangente al gráfico de las siguientes funciones, en los puntos indicados:

  \begin{itemize}
  \item[a)] $f(x,y)=x^2+y^2$ en $(0,0,0)$ y $(1,-1,2)$.

  \item[b)] $f(x,y)=xy-x^2y^2$ en $(1,1,0)$.
    \end{itemize}


  \item En cada caso, calcular la derivada parcial de orden superior indicada:
%    \begin{multicols}{2}
%      \item[a)] $f(x,y)=x^{2}y^{3}-2x^{4}y;$ $
%f_{xx}$ \item[b)] $f(x,y)=e^{xy^{2}};$ $f_{xy},\,f_{yx}$  \item[c)] $
%f(x,y)=x^{5}+x^{4}y^{4}+y^{2};$ $f_{xyx}$ \item[d)] $f(x,y)=e^{xy^2};$ $
%f_{yxy}$.
%\end{multicols}

  \item Determinar si alguna de las siguientes funciones es solución de la llamada \emph{Ecuación de Laplace} (bidimensional)  $u_{xx}+u_{yy}=0$:
     \begin{itemize}
  \item[a)] $u=x^2+y^2$, 

  \item[b)] $u=x^2-y^2$

  \item[c)] $u=\ln\sqrt{x^2+y^2}$

    \item[d)] $u=e^{-x}\cos y-e^{-y}\cos x$ 
    \end{itemize}
\item Verificar que la función $u=e^{-\alpha^2k^2t}\sen(kx)$ es solución de la \emph{Ecuación del calor} (bidimensional) $u_t=\alpha^2u_{xx}$, donde $\alpha$ y $k$ son constantes. 

\item \textbf{Relaciones de Maxwell:} Para un gas monoatómico ideal, tenemos las ecuación $PV = k_B NT$, y una expresión para la entropía

\[
S = k_B N \ln \left( \frac{V T^{3/2}}{N} \right)
\]

Verificar las siguientes igualdades (despejando siempre que sea necesario una variable en función de las otras):

\begin{enumerate}
\item $\frac{\partial S}{\partial V} = \frac{\partial P}{\partial T}$
\item $-\frac{\partial S}{\partial P} = \frac{\partial V}{\partial T}$
\item $\frac{\partial T}{\partial V} = - \frac{\partial P}{\partial S}$
\item $\frac{\partial T}{\partial P} = \frac{\partial V}{\partial S}$
\end{enumerate}

Cada una de las igualdades anteriores tiene un significado físico preciso en términos de distintos tipos de energía. Interpretar 
la tercera de las igualdades en términos de la energía interna del gas ideal, recordando que $U = \frac{3}{2} k_B N T$. 
(Sugerencia: calcular $\frac{\partial^2 U}{\partial S \partial V}$)
    
\end{enumerate}

\chapter{Funciones escalares III de V}

\section{Temario}

\begin{itemize}
%    \item Dominio, imágen, gráfico y conjuntos de nivel.
%    \item Límites y continuidad.
    \item Derivadas: %(derivadas parciales/direccionales, diferenciabilidad, gradiente y espacio tangente).
%
\begin{itemize}
    \item diferenciabilidad,
    \item gradiente,
    \item regla de la cadena.
\end{itemize}
    
%    \item Derivadas de orden superior (Taylor, criterio de la Hessiana).
%    \item Teorema de la Función Implícita.
%    \item Extremos condicionados (Multiplicadores de Lagrange).
\end{itemize}

\section{Notas}

\noindent Nociones sobre la diferencial.

\subsection{Diferencial de una función de dos variables}

\begin{definicion}
Sea \( f: D \subset \mathbb{R}^2 \to \mathbb{R} \) una función de dos variables. La \textbf{diferencial total} de \( f \) en el punto \( (x_0, y_0) \), denotada por \( df(x_0, y_0) \), es una aplicación lineal que aproxima el cambio en \( f \) cerca de \( (x_0, y_0) \). Si \( f \) es diferenciable en \( (x_0, y_0) \), entonces la diferencial se define como:
\[
df(x_0, y_0)(h, k) = \frac{\partial f}{\partial x}(x_0, y_0)h + \frac{\partial f}{\partial y}(x_0, y_0)k,
\]
donde \( h \) y \( k \) son incrementos en \( x \) y \( y \), respectivamente.
\end{definicion}

\begin{ejemplo}
Consideremos la función \( f(x, y) = x^2 + y^2 \). Las derivadas parciales son:
\[
\frac{\partial f}{\partial x} = 2x, \quad \frac{\partial f}{\partial y} = 2y.
\]
La diferencial de \( f \) en el punto \( (1, 2) \) es:
\[
df(1, 2)(h, k) = 2(1)h + 2(2)k = 2h + 4k.
\]
\end{ejemplo}

\begin{ejemplo}
Sea \( f(x, y) = \sin(xy) \). Las derivadas parciales son:
\[
\frac{\partial f}{\partial x} = y\cos(xy), \quad \frac{\partial f}{\partial y} = x\cos(xy).
\]
La diferencial de \( f \) en el punto \( \left(\frac{\pi}{2}, 1\right) \) es:
\[
df\left(\frac{\pi}{2}, 1\right)(h, k) = 1 \cdot \cos\left(\frac{\pi}{2} \cdot 1\right)h + \frac{\pi}{2} \cdot \cos\left(\frac{\pi}{2} \cdot 1\right)k = 0 \cdot h + \frac{\pi}{2} \cdot 0 \cdot k = 0.
\]
\end{ejemplo}

\noindent Una función diferenciable es continua. 

\subsection{Gradiente de una función de dos variables}

\begin{definicion}
Sea \( f: D \subset \mathbb{R}^2 \to \mathbb{R} \) una función de dos variables. El \textbf{gradiente} de \( f \) en el punto \( (x_0, y_0) \), denotado por \( \nabla f(x_0, y_0) \), es el vector formado por las derivadas parciales de \( f \) en ese punto:
\[
\nabla f(x_0, y_0) = \left( \frac{\partial f}{\partial x}(x_0, y_0), \frac{\partial f}{\partial y}(x_0, y_0) \right).
\]
El gradiente indica la dirección de máximo crecimiento de la función.
\end{definicion}

\begin{ejemplo}
Para la función \( f(x, y) = x^2 + y^2 \), el gradiente en \( (1, 2) \) es:
\[
\nabla f(1, 2) = \left( 2(1), 2(2) \right) = (2, 4).
\]
\end{ejemplo}

\begin{ejemplo}
Para la función \( f(x, y) = e^{x} \ln(y) \), las derivadas parciales son:
\[
\frac{\partial f}{\partial x} = e^{x} \ln(y), \quad \frac{\partial f}{\partial y} = \frac{e^{x}}{y}.
\]
El gradiente en \( (0, 1) \) es:
\[
\nabla f(0, 1) = \left( e^{0} \ln(1), \frac{e^{0}}{1} \right) = (0, 1).
\]
\end{ejemplo}

%\subsection{La diferencial}
%La diferencial de una función \(f(x, y)\) se define como:
%\[ df = f_x(x, y) dx + f_y(x, y) dy \]
%
%La diferencial se utiliza para aproximar cambios pequeños en la función.

%Diferencial

%
%\subsection{Gradiente}

\noindent
El gradiente generaliza la nocion de derivada para funciones de varias variables. Permite evaluar las variaciones de la funcion, sus amplitudes tan como sus direcciones. Permite tambien ubicar los puntos criticos (minimo y maximo en particular).
%Generaliza la derivada, permite ver como el grafico sube, baja, y ver donde estan los puntos criticos (min, max, etc).

\medskip

\noindent
Si \( f \) es diferenciable, entonces \( df({v})(\mathbf{a}) = \nabla f(\mathbf{a}) \cdot \mathbf{v} \).

\subsection{Gradiente y derivadas parciales, derivada direccional}

\noindent Una funcion diferenciable siempre admite gradiente, derivadas parciales y derivada direccional.

\paragraph{Derivada direccional}
La derivada direccional de \(f(x, y)\) en la dirección de un vector unitario \(\mathbf{u} = (a, b)\) es:
\[ D_{\mathbf{u}} f(x, y) = f_x(x, y) a + f_y(x, y) b = \nabla f(x,y) \cdot \mathbf{u}.\]

\paragraph{Ejemplos de funciónes con derivadas direccionales pero no diferenciable}
Consideremos la función: 
\[ f(x, y) = \begin{cases}
\displaystyle \frac{xy}{x^2 + y^2} & \text{si } (x, y) \neq (0, 0), \\
0 & \text{si } (x, y) = (0, 0).
\end{cases} \]
Hemos visto que no es continua en $(0,0)$, entonces tampoco es diferenciable en $(0,0)$. 
%Su gradiente es
%\[
%\nabla f(x, y) 
%= \left( \frac{\partial f}{\partial x}(x, y), \frac{\partial f}{\partial y}(x, y) \right)
%= \left( 
%\frac{3x^2 (x^2 + y^2) - 2x^4}{x^2+y^2}, 
%\frac{-2yx^3}{x^2+y^2}
% \right)
%.
%\]
%Entonces tenemos
%\[
%f\nabla f(x, y) 
%= \left( \frac{\partial f}{\partial x}(x, y), \frac{\partial f}{\partial y}(x, y) \right)
%= \left( 
%\frac{x^4 + 3 x^2 y^2}{x^2+y^2}, 
%\frac{-2yx^3}{x^2+y^2}
% \right)
%.
%\]
Sin embargo, admite derivadas parciales (entonces derivadas direccionales en las direcciones $(0,1)$ y $(1,0)$).

\medskip

\noindent Sea ahora
\[ f(x, y) = \begin{cases}
\displaystyle \frac{x^3}{x^2 + y^2} & \text{si } (x, y) \neq (0, 0), \\
0 & \text{si } (x, y) = (0, 0).
\end{cases} \]

\paragraph{Derivadas direccionales en \((0, 0)\)}
Para un vector unitario \(\mathbf{u} = (a, b)\), la derivada direccional en \((0, 0)\) es:
\[ D_{\mathbf{u}} f(0, 0) = 
\lim_{t \rightarrow 0}
\frac{f(0+ta,0+tb) - f(0,0)}{t} = 
\lim_{t \rightarrow 0}
\frac{(ta)^3}{t( (ta)^2 + (tb)^2)} = 
\frac{a^3}{a^2 + b^2}. \]
Esto muestra que la derivada direccional existe en todas las direcciones.

\paragraph{No diferenciabilidad en \((0, 0)\)}
Para verificar que \(f\) no es diferenciable en \((0, 0)\), consideramos el límite:
\[ \lim_{(x, y) \to (0, 0)} \frac{f(x, y) - f(0, 0) - f_x(0, 0) x - f_y(0, 0) y}{\sqrt{x^2 + y^2}} \]
Calculamos
\[
\frac{f(\delta x,0)-f(0,0)}{\delta x} = \frac{(\delta x)^3}{\delta x (\delta x)^2} =1
\]
y
\[
\frac{f(0, \delta y)-f(0,0)}{\delta y} = \frac{0}{\delta y (\delta y)^2} = 0.
\]
Deducimos que \(f_x(0, 0) = 1\) y \(f_y(0, 0) = 0\), así que el límite se convierte en:
\[ \lim_{(x, y) \to (0, 0)} \frac{\frac{x^3}{x^2 + y^2} - x}{\sqrt{x^2 + y^2}} 
= \frac{- x y^2}{(x^2 + y^2) \sqrt{x^2 + y^2}}
\]
Este límite no existe, como se puede ver al acercarse a lo largo de diferentes trayectorias: 
\[ \lim_{(x, 0) \to (0, 0)} %\frac{\frac{x^3}{x^2 + y^2} - x}{\sqrt{x^2 + y^2}} 
%= 
\frac{0}{(x^2) \sqrt{x^2}} = 0
\]
y
\[ \lim_{(x, x) \to (0, 0)} %\frac{\frac{x^3}{x^2 + y^2} - x}{\sqrt{x^2 + y^2}} 
%= 
\frac{- x x^2}{(x^2 + x^2) \sqrt{x^2 + x^2}} = \frac{-\sqrt{2}}{2},
\]
lo que demuestra que \(f\) no es diferenciable en \((0, 0)\).

\subsection{Relación entre gradiente y curvas de nivel}

\begin{definicion}
Las \textbf{curvas de nivel} de una función \( f: \mathbb{R}^2 \to \mathbb{R} \) son las curvas definidas por \( f(x, y) = c \), donde \( c \) es una constante. El gradiente \( \nabla f(x, y) \) es perpendicular a la curva de nivel que pasa por \( (x, y) \). Esto significa que, si \( \mathbf{r}(t) = (x(t), y(t)) \) es una parametrización de la curva de nivel, entonces:
\[
\nabla f(x(t), y(t)) \cdot \mathbf{r}'(t) = 0.
\]
\end{definicion}

\begin{ejemplo}
Para la función \( f(x, y) = x^2 + y^2 \), las curvas de nivel son circunferencias centradas en el origen. El gradiente en \( (1, 1) \) es \( \nabla f(1, 1) = (2, 2) \), que es perpendicular a la tangente de la circunferencia \( x^2 + y^2 = 2 \) en ese punto.
\end{ejemplo}

\subsection{Regla de la cadena para funciones de dos variables}

\begin{definicion}
Sea \( f: \mathbb{R}^2 \to \mathbb{R} \) una función diferenciable y sean \( x = g(t) \) e \( y = h(t) \) funciones diferenciables de una variable \( t \). Entonces, la derivada de la función compuesta \( F(t) = f(g(t), h(t)) \) está dada por:
\[
F'(t) = \frac{\partial f}{\partial x}(g(t), h(t)) \cdot g'(t) + \frac{\partial f}{\partial y}(g(t), h(t)) \cdot h'(t).
\]
\end{definicion}

\begin{ejemplo}
Sea \( f(x, y) = x^2 y \), \( x(t) = t^2 \), e \( y(t) = \sin(t) \). Entonces:
\[
F(t) = f(t^2, \sin(t)) = t^4 \sin(t).
\]
La derivada de \( F \) es:
\[
F'(t) = 2xy \cdot 2t + x^2 \cdot \cos(t) = 2(t^2)(\sin(t)) \cdot 2t + (t^2)^2 \cdot \cos(t) = 4t^3 \sin(t) + t^4 \cos(t).
\]
\end{ejemplo}

\subsection{Algoritmo de descenso de gradiente \emph{(bonus track)}}

\begin{definicion}
El \textbf{algoritmo de descenso de gradiente} es un método iterativo para encontrar el mínimo de una función \( f: \mathbb{R}^2 \to \mathbb{R} \). Partiendo de un punto inicial \( (x_0, y_0) \), se actualiza la posición según la regla:
\[
(x_{n+1}, y_{n+1}) = (x_n, y_n) - \alpha \nabla f(x_n, y_n),
\]
donde \( \alpha > 0 \) es el tamaño del paso. El algoritmo se detiene cuando \( \|\nabla f(x_n, y_n)\| \) es menor que una tolerancia prefijada.
\end{definicion}

\begin{ejemplo}
Para minimizar \( f(x, y) = x^2 + y^2 \) con \( \alpha = 0.1 \) y punto inicial \( (1, 1) \):
\[
\nabla f(x, y) = (2x, 2y).
\]
La primera iteración es:
\[
(x_1, y_1) = (1, 1) - 0.1(2, 2) = (0.8, 0.8).
\]
\end{ejemplo}

%\subsection{Vuelta al espacio tangente (bonus)}

%\begin{definition}
%Sea \( S \) una superficie en \( \mathbb{R}^3 \) definida por \( F(x, y, z) = 0 \). El \textbf{espacio tangente} a \( S \) en un punto \( \mathbf{p} = (x_0, y_0, z_0) \) es el plano dado por:
%\[
%\nabla F(\mathbf{p}) \cdot (x - x_0, y - y_0, z - z_0) = 0.
%\]
%\end{definition}

%\begin{example}
%Consideremos la esfera \( F(x, y, z) = x^2 + y^2 + z^2 - 1 = 0 \). El espacio tangente en \( (1, 0, 0) \) es:
%\[
%(2, 0, 0) \cdot (x - 1, y, z) = 0 \implies x = 1.
%\]
%\end{example}

\newpage

\section{Practico 4 (diferenciabilidad)}

\begin{enumerate}

\item Estudiar la diferenciabilidad de las siguientes funciones en los puntos
indicados:
\begin{itemize}
\item[a)] $f(x,y)=\left\{ 
\begin{array}{ccc}
\frac{x^{2}y^{3}}{x^{2}+y^{2}} & \text{si} & \left( x,y\right) \neq \left(
0,0\right)  \\ 
0 & \text{si} & \left( x,y\right) =\left( 0,0\right) 
\end{array}%
\right. $ en $\left( 0,0\right) $

\item[b)] $f(x,y)=\left\{ 
\begin{array}{ccc}
2x-y & \text{si} & xy\neq 0 \\ 
x+y & \text{si} & xy=0%
\end{array}%
\right. $ en $\left( 0,0\right) $, $\left( 1,0\right) $, $\left( 0,1\right) $
y $\left( 1,1\right) $





% \item[e)] $f(x,y)=\left\{ 
% \begin{array}{ccc}
% \frac{\ln \left( 1+x^{2}y\right) }{x^{2}+y^{2}} & \text{si} & \left(
% x,y\right) \neq \left( 0,0\right)  \\ 
% 0 & \text{si} & \left( x,y\right) =\left( 0,0\right) 
% \end{array}%
% \right. $ en $\left( 0,0\right) $

% \item[f)] $f(x,y)=\left\{ 
% \begin{array}{ccc}
% 4x^{2} & \text{si} & y=2 \\ 
% 2xy & \text{si} & y\neq 2%
% \end{array}%
% \right. $ en $\left( 0,0\right) $

% \item[g)] $f(x,y)=|y-2x^{2}|$ en $(1,1)$ y $(0,0)$.
\end{itemize}

\item Probar que las siguientes funciones son diferenciables en todo punto de su dominio:
\begin{multicols}{2}
\item[a)] $f(x,y)=\frac{2x+y}{x-y}$
\item[b)] $f(x,y,z)=\ln(x^2-y^2+z^2)$

  \item[c)] $f(x,y,z)=e^{x+y}-z^2$
  \item[d)] $f(x,y)= x\cos(y)$
\end{multicols}


\item Hallar la aproximaci\'{o}n lineal de $f(x,y)=\sqrt{20-x^{2}-7y^{2}}$ en $
\left( 2,1\right) $ y usarla para calcular de manera aproximada $f(1{.}95,1{.}08)$ (estamos escribiendo los decimales de la manera anglosajona $1,95=1{.}95$).

\item Hallar y representar gráficamente el campo de gradientes y las curvas de nivel para las siguientes funciones en su correspondiente dominio. Verificar en cada caso la relación geométrica esperada entre la curva de nivel y el gradiente para los puntos indicados, y deducir la ecuación de la recta tangente a dicha curva en tales puntos.
\begin{multicols}{2}
\item[a)] $f(x,y)=\frac{x+y}{x-y},\ (1,-1)$ 

\item[b)] $f(x,y)=\ln(x^2+y^2),\ (1.0)\ \mbox{y}\ (1.1)$ 

\item[c)] $f(x,y)=4x^2+9y^2,\ (\sqrt{3}/2,\sqrt{2}/3)$ 

\item[d)] $f(x,y)=4x^2-9y^2,\ (\sqrt{3}/2,\sqrt{2}/3) \ \mbox{y}\ (1,1)$

%d) $E(\theta,\dot{\theta})=\dot{\theta}^2+(1-\cos(\theta))$ (funci\'on energía del péndulo).
\end{multicols}

% \item Hallar el espacio tangente y la direcci\'on perpendicular a las siguientes curvas de nivel en los puntos indicados:

% a) $x-y=1$ en $(1,0)$ b) $2x^2+y^2=3$ en $(1,1)$ c)$\ln(x^2+y^2)=0$ en $(\frac{1}{\sqrt{2}},-\frac{1}{\sqrt{2}})$

% c) $x\sen(y)+1=3$ en $(2,\frac{\pi}{2})$ d) $1=\sen(x^2-y^2)$ en $(\sqrt{\frac{\pi}{2}},0)$.

  \item Calcular la ecuación de la recta tangente a las curvas de nivel de las siguientes funciones, en el punto indicado:
    \begin{enumerate}
    \item $f(x,y)=\sen(x)\cos(y), p=(\pi/4,\pi/4)$.
      \item $f(x,y)=4x^2+9y^2, p=(1/2,1/3)$.
      \end{enumerate}

\item Sea $g_0(x,y)$ una función diferenciable en $\R^2$. 
\begin{itemize}
\item[a)] Determinar qué condiciones tiene que verificar una función $f(x,y)$ para que sus curvas de nivel intersecten ortogonmente las curvas de nivel de $g_0(x,y)$; entendemos que dos curvas son ortogonales en un punto $(a,b)$ si sus rectas tangentes en dicho punto son perpendiculares.
\item[b)] Si $g_0(x,y)=x^2+y^2$, encontrar $f(x,y)$ que satisfaga la condición determinada en $a)$ (observar que en este caso uno puede prever cómo tienen que ser las curvas de nivel de $f(x,y)$ de manera geométrica).

\item[c)] Ídem con $g_0(x,y)=xy$.
\end{itemize}


% \item En los siguientes casos (del Pr. 2 Ej. 3), estudiar la diferenciabilidad de $f$ y bosquejar el campo de gradientes:

% a) $f(x,y,z)=2x+y-3z$, \ \ b) $f(x,y,z)=x^2-y^2+z^2$, \ \ c) $f(x,y,z)=e^{x+y}-z$ \\
% d) $f(x,y,z)= x\cos(y)$, \ \ e) $f(x_1,\dots,x_n)=x_1+\dots+x_n$ \ \ f) $f(x_1,\dots,x_n)=x_1^2+\dots+x_n^2$.  

% \item En los siguientes casos hallar el espacio tangente y el plano tangante. Cuando sea posible, representar graficamente.

% a) $x^2+y^2+z^2=6$ en $(1,1,2)$ \ \ b) $x^2-3y^2+z^2=2$ en $(1,1,2)$ \ \ c) $e^{xy}\cos(z)=0$ en $(0,1,\frac{\pi}{2})$\\
% d) $x_1^2+\dots+x_n^2=n$ en $(1,\dots,1)$ \ \ e) $x_1x_4-x_2x_3x_1+x_3x_2=1$ en $(1,1,1,1)$. 


\item 
La presi\'{o}n, el volumen y la temperatura de un mol de gas ideal se
relacionan mediante la ecuaci\'{o}n (ley de los gases) $PV=8.31T$ donde $P$ se mide en
kilopascales, $V$ en litros y $T$ en grados kelvin. Calcular una aproximación lineal de la variación de presi\'{o}n en el gas, si el volumen se incrementa
de $12$ para $12{.}3$ litros y la temperatura disminuye de $310$ a $305$ grados kelvin.


%\item En los siguientes casos hallar el espacio tangente y el plano tangante. Cuando sea posible, representar graficamente.


% \item Para el potencial

% \[
% V (x,y) = \left( - \frac{\alpha}{(x^2 + y^2)^{1/2}} + \frac{\beta}{x^2 + y^2 } \right) \left( \frac{x^2}{x^2 + y^2} + 1 \right)
% \]

% bosquejar el campo de gradiente.

% \item En un problema electrostático con simetría cilíndrica, el potencial no depende de la coordenada $z$, y se lo puede tratar como una función solamente de $x$ e $y$. 
% Hallar la expresión de cada uno de los siguientes potenciales en función de $x,y$ y calcular las componentes $x$ e $y$ del campo eléctrico correspondiente 
% ($r$ y $\theta$ son coordenadas polares):
%
% \begin{enumerate}
% \item $V(r,\theta) = \ln(r)$
% \item $V(r,\theta) = r + \frac{2}{r}$
% \item $V(r,\theta) = \left( r + \frac{2}{r} \right)(\cos (\theta) + 2\sin (\theta))$
% \item $V(r,\theta) = \left( r^2 + \frac{2}{r^2} \right)(3\cos (2\theta) - \sin (2\theta))$
% \item $V(r,\theta) = \left( r^2 + \frac{2}{r^2} \right)(\cos (3\theta) + \sin (3\theta))$
% \end{enumerate}

% Esbozar en el plano, en cada caso, el campo eléctrico.

% %\textbf{Para pensar:} ¿el último potencial puede corresponder a un sistema electrostático real?

% \item Se tienen tres cargas puntuales $q_1,q_2,q_3$ en el plano. Probar que la única configuración de equilibrio es cuando están alineadas, que dos tienen el mismo signo y la otra tiene signo distinto y se encuentra ubicada entre las dos primeras. Calcular la distancia relativa entre las cargas en función de sus cargas eléctricas.

% \textbf{Nota}: el potencia de una carga puntual $q$ ubicada en el origen es $V(x,y,z) = \frac{kq}{(x^2 + y^2 + z^2)^{1/2}}$. Cuando nos restringimos a $z=0$, 
% el potencial resulta una función que depende sólo de $x$ e $y$.

% %\item ¿Es posible ubicar cuatro cargas puntuales, dos positivas y dos negativas, en el plano de manera que estén en equilibrio?
\item En los siguientes casos usar la regla de la cadena para calcular $(f\circ \gamma)'(t_0)$.

\begin{enumerate}
\item $f(x,y)=e^{xy}\cos(x)$, $\gamma(t)=(2\sen(t),3\cos(t))$, $t_0=\frac{\pi}{2}$.
\item $f(x,y)=x^2-y^2$, $\gamma(t)=(t,\arctg(t)+1)$, $t_0=0$.
\item $f(x,y,z)=z+\ln(x+y)$, $\gamma(t)=(\cos(t),\sen(t), t)$, $t=\frac{\pi}{4}$. 
\end{enumerate}

%\item Calcular las derivadas parciales de $f\circ g(p)$, donde $f$ es como en el ejercicio anterior, y en cada caso $g$ y $p$ es uno de los siguientes, respectivamente:
%
%  \begin{enumerate}
%  \item $g(x,y,z)=(x^3,yz^2)$, $p=(1,1,-1)$.
%  \item $g(r,\theta)=(r\cos\theta, r\sen\theta)$, $p=(\sqrt{3},\pi)$.
%    \item $g(x,y)=(x+y, xy, y^2), p=(0,-1)$.
%    \end{enumerate}

\end{enumerate}

\chapter{Funciones escalares IV de V}

\section{Temario}

\begin{itemize}
%    \item Dominio, imágen, gráfico y conjuntos de nivel.
%    \item Límites y continuidad.
%    \item Derivadas (derivadas parciales/direccionales, diferenciabilidad, gradiente y espacio tangente).
    \item Derivadas de orden superior (Taylor, criterio de la Hessiana).
%    \item Teorema de la Función Implícita.
%    \item Extremos condicionados (Multiplicadores de Lagrange).
\end{itemize}

%\title{Cursillo: Derivadas de Orden Superior y Extremos de Funciones de Dos Variables}
%\author{Franz Chouly}
%\date{\today}

%\begin{document}

%\maketitle

\section{Notas}

\subsection{Derivadas de Orden Superior para Funciones de Dos Variables}

\begin{definicion}
    Sea \( f: \mathbb{R}^2 \to \mathbb{R} \) una función diferenciable. Las \textbf{derivadas parciales de segundo orden} de \( f \) en un punto \( (a, b) \) son:
    \[
    \frac{\partial^2 f}{\partial x^2}(a, b), \quad \frac{\partial^2 f}{\partial y^2}(a, b), \quad \frac{\partial^2 f}{\partial x \partial y}(a, b), \quad \frac{\partial^2 f}{\partial y \partial x}(a, b).
    \]
    Si estas derivadas existen y son continuas en un entorno de \( (a, b) \), se dice que \( f \) es de \textbf{clase \( C^2 \)} en ese punto.
\end{definicion}

\begin{ejemplo}
    Sea \( f(x, y) = x^2 y + \sin(xy) \). Calcular las derivadas parciales de segundo orden.
    \[
    \frac{\partial f}{\partial x} = 2xy + y \cos(xy), \quad \frac{\partial f}{\partial y} = x^2 + x \cos(xy).
    \]
    \[
    \frac{\partial^2 f}{\partial x^2} = 2y - y^2 \sin(xy), \quad \frac{\partial^2 f}{\partial y^2} = -x^2 \sin(xy).
    \]
    \[
    \frac{\partial^2 f}{\partial x \partial y} = 2x + \cos(xy) - xy \sin(xy), \quad \frac{\partial^2 f}{\partial y \partial x} = 2x + \cos(xy) - xy \sin(xy).
    \]
\end{ejemplo}

\begin{ejemplo}
    Sea \( f(x, y) = e^{x+y} \). Calcular las derivadas parciales de segundo orden.
    \[
    \frac{\partial f}{\partial x} = e^{x+y}, \quad \frac{\partial f}{\partial y} = e^{x+y}.
    \]
    \[
    \frac{\partial^2 f}{\partial x^2} = e^{x+y}, \quad \frac{\partial^2 f}{\partial y^2} = e^{x+y}.
    \]
    \[
    \frac{\partial^2 f}{\partial x \partial y} = e^{x+y}, \quad \frac{\partial^2 f}{\partial y \partial x} = e^{x+y}.
    \]
\end{ejemplo}

\section{Teorema de Schwarz}

\begin{teorema}[Teorema de Schwarz]
    Si \( f: \mathbb{R}^2 \to \mathbb{R} \) es de clase \( C^2 \) en un entorno de \( (a, b) \), entonces:
    \[
    \frac{\partial^2 f}{\partial x \partial y}(a, b) = \frac{\partial^2 f}{\partial y \partial x}(a, b).
    \]
\end{teorema}

\begin{ejemplo}
    Para \( f(x, y) = x^2 y + \sin(xy) \), se verificó en el ejemplo anterior que:
    \[
    \frac{\partial^2 f}{\partial x \partial y} = \frac{\partial^2 f}{\partial y \partial x}.
    \]
\end{ejemplo}

\begin{ejemplo}
    Para \( f(x, y) = \ln(x^2 + y^2) \), calcular las derivadas mixtas:
    \[
    \frac{\partial f}{\partial x} = \frac{2x}{x^2 + y^2}, \quad \frac{\partial f}{\partial y} = \frac{2y}{x^2 + y^2}.
    \]
    \[
    \frac{\partial^2 f}{\partial x \partial y} = \frac{-4xy}{(x^2 + y^2)^2}, \quad \frac{\partial^2 f}{\partial y \partial x} = \frac{-4xy}{(x^2 + y^2)^2}.
    \]
\end{ejemplo}

\begin{ejemplo}
    Aquí un ejemplo explícito y sencillo donde las derivadas cruzadas \( \frac{\partial^2 f}{\partial x \partial y} \) y \( \frac{\partial^2 f}{\partial y \partial x} \) no coinciden en un punto, debido a que no son continuas en ese punto.
%### Función de ejemplo:
\[
f(x, y) = \begin{cases}
xy \cdot \frac{x^2 - y^2}{x^2 + y^2} & \text{si } (x, y) \neq (0, 0), \\
0 & \text{si } (x, y) = (0, 0).
\end{cases}
\]
%
%---
%
%### Paso 1: 
Verificamos que \( f \) es diferenciable en \( (0, 0) \)
Sus derivadas parciales en \( (0, 0) \) son:
  \[
  \frac{\partial f}{\partial x}(0, 0) = \lim_{h \to 0} \frac{f(h, 0) - f(0, 0)}{h} = \lim_{h \to 0} \frac{0 - 0}{h} = 0,
  \]
  \[
  \frac{\partial f}{\partial y}(0, 0) = \lim_{h \to 0} \frac{f(0, h) - f(0, 0)}{h} = \lim_{h \to 0} \frac{0 - 0}{h} = 0.
  \]
La función es diferenciable en \( (0, 0) \) porque las derivadas parciales existen y el límite de la diferencial coincide con el plano tangente.
%### Paso 2: 
Calculamos las derivadas cruzadas en \( (0, 0) \).
%#### Derivada \( \frac{\partial^2 f}{\partial x \partial y}(0, 0) \):
Calculamos \( \frac{\partial f}{\partial y}(x, 0) \) para \( x \neq 0 \):
   \[
   f(x, y) = xy \cdot \frac{x^2 - y^2}{x^2 + y^2} \implies f(x, 0) = 0.
   \]
   \[
   \frac{\partial f}{\partial y}(x, 0) = \lim_{h \to 0} \frac{f(x, h) - f(x, 0)}{h} = \lim_{h \to 0} \frac{xh \cdot \frac{x^2 - h^2}{x^2 + h^2}}{h} = \lim_{h \to 0} x \cdot \frac{x^2 - h^2}{x^2 + h^2} = x.
   \]
Ahora, calculamos \( \frac{\partial^2 f}{\partial x \partial y}(0, 0) \):
   \[
   \frac{\partial^2 f}{\partial x \partial y}(0, 0) = \lim_{h \to 0} \frac{\frac{\partial f}{\partial y}(h, 0) - \frac{\partial f}{\partial y}(0, 0)}{h} = \lim_{h \to 0} \frac{h - 0}{h} = 1.
   \]
%#### Derivada \( \frac{\partial^2 f}{\partial y \partial x}(0, 0) \):
Calculamos \( \frac{\partial f}{\partial x}(0, y) \) para \( y \neq 0 \):
   \[
   f(0, y) = 0 \implies \frac{\partial f}{\partial x}(0, y) = \lim_{h \to 0} \frac{f(h, y) - f(0, y)}{h} = \lim_{h \to 0} \frac{hy \cdot \frac{h^2 - y^2}{h^2 + y^2}}{h} = y \cdot \frac{-y^2}{y^2} = -y.
   \]
Ahora, calculamos \( \frac{\partial^2 f}{\partial y \partial x}(0, 0) \):
   \[
   \frac{\partial^2 f}{\partial y \partial x}(0, 0) = \lim_{h \to 0} \frac{\frac{\partial f}{\partial x}(0, h) - \frac{\partial f}{\partial x}(0, 0)}{h} = \lim_{h \to 0} \frac{-h - 0}{h} = -1.
   \]
%
%---
%
%### 
Conclusión:
En el punto \( (0, 0) \):
\[
\frac{\partial^2 f}{\partial x \partial y}(0, 0) = 1 \quad \text{y} \quad \frac{\partial^2 f}{\partial y \partial x}(0, 0) = -1.
\]
Las derivadas cruzadas no coinciden porque no son continuas en \( (0, 0) \). Esto ilustra que el Teorema de Schwarz no se aplica si las derivadas parciales mixtas no son continuas.
\end{ejemplo}

\section{Fórmula de Taylor de Orden 2}

\begin{definicion}
    Sea \( f: \mathbb{R}^2 \to \mathbb{R} \) de clase \( C^2 \) en un entorno de \( (a, b) \). La \textbf{fórmula de Taylor de orden 2} de \( f \) en \( (a, b) \) es:
    \[
    f(x, y) = f(a, b) + \nabla f(a, b) \cdot (x - a, y - b) + \frac{1}{2} (x - a, y - b)^T H_f(a, b) (x - a, y - b) + R_2(x, y),
    \]
    donde \( \nabla f(a, b) \) es el gradiente, \( H_f(a, b) \) es la matriz hessiana, y \( R_2(x, y) \) es el resto.
\end{definicion}

%Sea \( f: \mathbb{R}^2 \to \mathbb{R} \) una función de clase \( C^3 \) en un entorno del punto \( (a, b) \). El **polinomio de Taylor de orden 2** de \( f \) centrado en \( (a, b) \) es:
\noindent
Otra forma de escribirlo es la siguiente:
\[
\begin{aligned}
f(x, y) \simeq & \ f(a, b) \\
& + \frac{\partial f}{\partial x}(a, b)(x - a) + \frac{\partial f}{\partial y}(a, b)(y - b) \\
& + \frac{1}{2} \left[ \frac{\partial^2 f}{\partial x^2}(a, b)(x - a)^2 + 2 \frac{\partial^2 f}{\partial x \partial y}(a, b)(x - a)(y - b) + \frac{\partial^2 f}{\partial y^2}(a, b)(y - b)^2 \right].
\end{aligned}
\]
Donde
 \( \frac{\partial f}{\partial x}(a, b) \) y \( \frac{\partial f}{\partial y}(a, b) \) son las derivadas parciales de primer orden en \( (a, b) \);
\( \frac{\partial^2 f}{\partial x^2}(a, b) \), \( \frac{\partial^2 f}{\partial y^2}(a, b) \), y \( \frac{\partial^2 f}{\partial x \partial y}(a, b) \) son las derivadas parciales de segundo orden en \( (a, b) \).


\begin{ejemplo}
    Para \( f(x, y) = e^{x + y} \), desarrollar la fórmula de Taylor de orden 2 en \( (0, 0) \):
    \[
    f(x, y) \simeq 1 + (x + y) + \frac{1}{2}(x^2 + 2xy + y^2).
    \]
\end{ejemplo}

%\begin{ejemplo}
%    Para \( f(x, y) = \sin(x) \cos(y) \), desarrollar la fórmula de Taylor de orden 2 en \( (0, 0) \):
%    \[
%    f(x, y) \simeq 0 + x - \frac{x^3}{6} - xy^2 + \cdots
%    \]
%    (Nota: Se simplifica para mostrar solo los términos de orden 2).
%\end{ejemplo}

%### Ejemplo 1: 

\begin{ejemplo}
Desarrollar el polinomio de Taylor de orden 2 de \( f(x, y) = e^{x} \sin(y) \) en \( (0, 0) \).
1. Derivadas parciales de primer orden:
   \[
   \frac{\partial f}{\partial x}(0, 0) = e^{0} \sin(0) = 0, \quad \frac{\partial f}{\partial y}(0, 0) = e^{0} \cos(0) = 1.
   \]
2. Derivadas parciales de segundo orden:
   \[
   \frac{\partial^2 f}{\partial x^2}(0, 0) = e^{0} \sin(0) = 0, \quad \frac{\partial^2 f}{\partial y^2}(0, 0) = -e^{0} \sin(0) = 0,
   \]
   \[
   \frac{\partial^2 f}{\partial x \partial y}(0, 0) = e^{0} \cos(0) = 1.
   \]
3. Polinomio de Taylor:
   \[
   f(x, y) \simeq 0 + 0 \cdot x + 1 \cdot y + \frac{1}{2} \left[ 0 \cdot x^2 + 2 \cdot 1 \cdot x y + 0 \cdot y^2 \right] = y + xy.
   \]
\end{ejemplo}

\begin{ejemplo}
%### Ejemplo 2: 
Sea \( f(x, y) = \ln(1 + x + y) \).
Desarrollar el polinomio de Taylor de orden 2 en \( (0, 0) \):
%
1. Derivadas parciales de primer orden:
   \[
   \frac{\partial f}{\partial x}(0, 0) = \frac{1}{1 + 0 + 0} = 1, \quad \frac{\partial f}{\partial y}(0, 0) = \frac{1}{1 + 0 + 0} = 1.
   \]
%
2. Derivadas parciales de segundo orden:
   \[
   \frac{\partial^2 f}{\partial x^2}(0, 0) = -\frac{1}{(1 + 0 + 0)^2} = -1, \quad \frac{\partial^2 f}{\partial y^2}(0, 0) = -\frac{1}{(1 + 0 + 0)^2} = -1,
   \]
   \[
   \frac{\partial^2 f}{\partial x \partial y}(0, 0) = -\frac{1}{(1 + 0 + 0)^2} = -1.
   \]
%
3. Polinomio de Taylor:
   \[
   f(x, y) \simeq 0 + 1 \cdot x + 1 \cdot y + \frac{1}{2} \left[ -1 \cdot x^2 + 2 \cdot (-1) \cdot x y - 1 \cdot y^2 \right] = x + y - \frac{1}{2}(x^2 + 2xy + y^2).
   \]

\end{ejemplo}



\section{Caracterización de los Extremos de una Función}

\begin{definicion}
    Sea \( f: \mathbb{R}^2 \to \mathbb{R} \) de clase \( C^2 \) en un entorno de \( (a, b) \), y supongamos que \( \nabla f(a, b) = (0, 0) \). Definimos el \textbf{determinante hessiano} como:
    \[
    D = f_{xx}(a, b) f_{yy}(a, b) - [f_{xy}(a, b)]^2.
    \]
    \begin{itemize}
        \item Si \( D > 0 \) y \( f_{xx}(a, b) > 0 \), entonces \( (a, b) \) es un \textbf{mínimo local}.
        \item Si \( D > 0 \) y \( f_{xx}(a, b) < 0 \), entonces \( (a, b) \) es un \textbf{máximo local}.
        \item Si \( D < 0 \), entonces \( (a, b) \) es un \textbf{punto de silla}.
        \item Si \( D = 0 \), no se puede concluir.
    \end{itemize}
\end{definicion}

\noindent Observamos que también se puede utilisar la traza de la matriz hessiana.

\begin{ejemplo}
    Sea \( f(x, y) = x^2 + y^2 \). El punto crítico es \( (0, 0) \):
    \[
    f_{xx} = 2, \quad f_{yy} = 2, \quad f_{xy} = 0.
    \]
    \[
    D = 4 > 0 \quad \text{y} \quad f_{xx} > 0 \implies \text{Mínimo local en } (0, 0).
    \]
\end{ejemplo}

\begin{ejemplo}
    Sea \( f(x, y) = x^2 - y^2 \). El punto crítico es \( (0, 0) \):
    \[
    f_{xx} = 2, \quad f_{yy} = -2, \quad f_{xy} = 0.
    \]
    \[
    D = -4 < 0 \implies \text{Punto de silla en } (0, 0).
    \]
\end{ejemplo}

%\end{document}

\begin{ejemplo} Veamos

\textbf{Función: } \[ f(x, y) = \frac{x^3}{3} - y^2 - x + 1 \]

\textbf{Dominio: }
La función está definida para todo \( (x, y) \in \mathbb{R}^2 \).

\textbf{Derivadas parciales:}
\[
\frac{\partial f}{\partial x} = x^2 - 1, \quad \frac{\partial f}{\partial y} = -2y.
\]

\textbf{Puntos críticos:}
Para encontrar los puntos críticos, resolvemos el sistema:
\[
x^2 - 1 = 0 \quad \text{y} \quad -2y = 0.
\]

Las soluciones son \( (1, 0) \) y \( (-1, 0) \).

\textbf{Matriz hessiana:}
\[
H_f(x, y) = \begin{bmatrix}
\frac{\partial^2 f}{\partial x^2} & \frac{\partial^2 f}{\partial x \partial y} \\
\frac{\partial^2 f}{\partial y \partial x} & \frac{\partial^2 f}{\partial y^2}
\end{bmatrix}
= \begin{bmatrix}
2x & 0 \\
0 & -2
\end{bmatrix}.
\]

\textbf{Análisis de los puntos críticos:}
\begin{itemize}
\item \textbf{Punto crítico \( (1, 0) \):}
La matriz hessiana es:
\[
H_f(1, 0) = \begin{bmatrix}
2 & 0 \\
0 & -2
\end{bmatrix}.
\]
El determinante es \( \det(H_f(1, 0)) = -4 < 0 \), por lo que \( (1, 0) \) es un \textbf{punto silla}.

\item \textbf{Punto crítico \( (-1, 0) \):}
La matriz hessiana es:
\[
H_f(-1, 0) = \begin{bmatrix}
-2 & 0 \\
0 & -2
\end{bmatrix}.
\]
El determinante es \( \det(H_f(-1, 0)) = 4 > 0 \) y \( \frac{\partial^2 f}{\partial x^2} = -2 < 0 \), por lo que \( (-1, 0) \) es un \textbf{máximo local}.
\end{itemize}
\end{ejemplo}


\begin{ejemplo} Veamos:

\textbf{Función: } \[ f(x, y) = x^4 + y^4 - 4xy + 1 \]

\textbf{Dominio: }
La función está definida para todo \( (x, y) \in \mathbb{R}^2 \).

\textbf{Derivadas parciales:}
\[
\frac{\partial f}{\partial x} = 4x^3 - 4y, \quad \frac{\partial f}{\partial y} = 4y^3 - 4x.
\]

\textbf{Puntos críticos:}
Resolvemos el sistema:
\[
4x^3 - 4y = 0 \quad \text{y} \quad 4y^3 - 4x = 0.
\]
Las soluciones son \( (0, 0) \), \( (1, 1) \) y \( (-1, -1) \).

\textbf{Matriz hessiana:}
\[
H_f(x, y) = \begin{bmatrix}
12x^2 & -4 \\
-4 & 12y^2
\end{bmatrix}.
\]

\textbf{Análisis de los puntos críticos:}
\begin{itemize}
\item \textbf{Punto crítico \( (0, 0) \):}
La matriz hessiana es:
\[
H_f(0, 0) = \begin{bmatrix}
0 & -4 \\
-4 & 0
\end{bmatrix}.
\]
El determinante es \( \det(H_f(0, 0)) = -16 < 0 \), por lo que \( (0, 0) \) es un \textbf{punto silla}.

\item \textbf{Punto crítico \( (1, 1) \):}
La matriz hessiana es:
\[
H_f(1, 1) = \begin{bmatrix}
12 & -4 \\
-4 & 12
\end{bmatrix}.
\]
El determinante es \( \det(H_f(1, 1)) = 128 > 0 \) y \( \frac{\partial^2 f}{\partial x^2} = 12 > 0 \), por lo que \( (1, 1) \) es un \textbf{mínimo local}.

\item \textbf{Punto crítico \( (-1, -1) \):}
La matriz hessiana es:
\[
H_f(-1, -1) = \begin{bmatrix}
12 & -4 \\
-4 & 12
\end{bmatrix}.
\]
El determinante es \( \det(H_f(-1, -1)) = 128 > 0 \) y \( \frac{\partial^2 f}{\partial x^2} = 12 > 0 \), por lo que \( (-1, -1) \) es un \textbf{mínimo local}.
\end{itemize}
\end{ejemplo}

\begin{ejemplo} Veamos

\textbf{Función: } \[ f(x, y) = -y^2 + x^2y + \ln(xy) \]

\textbf{Dominio: }
La función está definida para \( xy > 0 \) 

(es decir, \( x > 0 \) y \( y > 0 \), o \( x < 0 \) y \( y < 0 \)).

\textbf{Derivadas parciales:}
\[
\frac{\partial f}{\partial x} = 2xy + \frac{y}{xy} = 2xy + \frac{1}{x}, \quad \frac{\partial f}{\partial y} = -2y + x^2 + \frac{x}{xy} = -2y + x^2 + \frac{1}{y}.
\]

\textbf{Punto crítico:}
Resolvemos el sistema:
\[
2xy + \frac{1}{x} = 0 \quad \text{y} \quad -2y + x^2 + \frac{1}{y} = 0.
\]

La solución es \( (-1, -\frac{1}{2}) \).

\textbf{Matriz hessiana:}
\[
H_f(x, y) = \begin{bmatrix}
2y - \frac{1}{x^2} & 2x \\
2x & -2 - \frac{1}{y^2}
\end{bmatrix}.
\]

\textbf{Análisis del punto crítico \( (-1, -\frac{1}{2}) \):}
La matriz hessiana en este punto es:
\[
H_f(-1, -\frac{1}{2}) = \begin{bmatrix}
-1 - 1 & -2 \\
-2 & -2 - 4
\end{bmatrix}
= \begin{bmatrix}
-2 & -2 \\
-2 & -6
\end{bmatrix}.
\]
El determinante es \( \det(H_f(-1, -\frac{1}{2})) = (-2)(-6) - (-2)(-2) = 12 - 4 = 8 > 0 \) y \( \frac{\partial^2 f}{\partial x^2} = -2 < 0 \), por lo que \( (-1, -\frac{1}{2}) \) es un \textbf{máximo local}.
\end{ejemplo}


%\chapter{In progress}

%\newpage 

\chapter{Funciones escalares IVbis de V}

\noindent
Seguimos y terminamos el curso 4 con 
un ejemplo en tres dimensiones y extremos de una funcion en un compacto (extremos absolutos).

\section{Notas}

\noindent
Primero veamos los extremos en tres dimensiones. Lo dejamos como un ejercicio.

\subsection{Ejemplo divertido en 3D}

%\documentclass{article}
%\usepackage{amsmath}

%\begin{document}

\paragraph{Problema: Volumen de una caja con restricción}
%
Queremos maximizar el volumen \( V = xyz \) bajo la restricción \( xy + 2xz + 2yz = 12 \).

\paragraph{Paso 1: Expresar \( z \) en función de \( x \) e \( y \)}
Despejamos \( z \) de la restricción:
\[
xy + 2xz + 2yz = 12 \implies xy + z( 2x + 2y) = 12 \implies z = \frac{12-xy}{ 2x + 2y}.
\] 

\paragraph{Paso 2: Sustituir \( z \) en la función de volumen}

El volumen \( V \) se expresa ahora como función de \( x \) e \( y \):
\[
V(x, y) = xyz = xy \left( \frac{12-xy}{2x + 2y} \right) = \frac{xy(12-xy)}{2x + 2y}.
\]

%\subsection*{Paso 3: Calcular las derivadas parciales de \( V(x, y) \)}

%\dots
%Calculamos las derivadas parciales respecto a \( x \) e \( y \):
%\dots
%\[
%\frac{\partial V}{\partial x} = \frac{12y(2x + 2y) - 12xy(y + 2)}{(2x + 2y)^2} = 0,
%\]
%\[
%\frac{\partial V}{\partial y} = \frac{12x(xy + 2x + 2y) - 12xy(x + 2)}{(xy + 2x + 2y)^2} = 0.
%\]

%\subsection*{Paso 4: Resolver el sistema de ecuaciones}
%
%Igualamos los numeradores a cero y simplificamos:

%1. Para \( \frac{\partial V}{\partial x} = 0 \):
%\[
%12y(xy + 2x + 2y) - 12xy(y + 2) = 0 \implies y(xy + 2x + 2y) - xy(y + 2) = 0.
%\]
%Desarrollando:
%\[
%xy^2 + 2xy + 2y^2 - xy^2 - 2xy = 0 \implies 2y^2 = 0 \implies y = 0.
%\]
%Sin embargo, \( y = 0 \) no es una solución válida en este contexto físico. Revisemos el cálculo:

%\[
%\frac{\partial V}{\partial x} = \frac{12y(xy + 2x + 2y) - 12xy(y + 2)}{(xy + 2x + 2y)^2} = 0.
%\]
%El numerador debe ser cero:
%\[
%12y(xy + 2x + 2y) - 12xy(y + 2) = 0.
%\]
%Simplificando:
%\[
%y(xy + 2x + 2y) - xy(y + 2) = 0,
%\]
%\[
%xy^2 + 2xy + 2y^2 - xy^2 - 2xy = 0 \implies 2y^2 = 0.
%\]
%Esto sugiere que hay un error en el enfoque. En su lugar, resolvamos el sistema de manera alternativa.

%\subsection*{Paso 5: Simplificar el problema}

%Observamos que la función \( V(x, y) \) es simétrica en \( x \) e \( y \). Supongamos \( x = y \). Sustituyendo \( y = x \) en la restricción:
%\[
%x^2 + 2xz + 2xz = 12 \implies x^2 + 4xz = 12 \implies z = \frac{12 - x^2}{4x}.
%\]
%El volumen se convierte en:
%\[
%V(x) = x \cdot x \cdot \frac{12 - x^2}{4x} = \frac{x(12 - x^2)}{4}.
%\]

%\subsection*{Paso 6: Encontrar el punto crítico de \( V(x) \)}

%Derivamos \( V(x) \) respecto a \( x \):
%\[
%V(x) = \frac{12x - x^3}{4},
%\]
%\[
%V'(x) = \frac{12 - 3x^2}{4}.
%\]
%Igualamos \( V'(x) = 0 \):
%\[
%12 - 3x^2 = 0 \implies x^2 = 4 \implies x = 2 \quad (\text{ya que } x > 0).
%\]
%Por lo tanto, \( y = x = 2 \), y sustituyendo en la expresión para \( z \):
%\[
%z = \frac{12 - (2)^2}{4 \cdot 2} = \frac{12 - 4}{8} = 1.
%\]

%\subsection*{Paso 7: Verificar la solución}
%
%Sustituyendo \( x = 2 \), \( y = 2 \), y \( z = 1 \) en la restricción:
%\[
%(2)(2) + 2(2)(1) + 2(2)(1) = 4 + 4 + 4 = 12.
%\]
%La restricción se satisface.
%
%\subsection*{Paso 8: Calcular el volumen máximo}
%
%\[
%V = xyz = (2)(2)(1) = 4.
%\]

%\subsection*{Conclusión}
%
%El volumen máximo de la caja bajo la restricción dada es \( V = 4 \), y se alcanza cuando \( x = 2 \), \( y = 2 \), y \( z = 1 \).

%\end{document}
%

%\noindent
%Volumen de un carton: $V=xyz$ bajo la condicion $xy + 2xz + 2yz = 12$. 
Minimizando $V(x,y)$ se encuentra $x=2$, $y=2$ y $z=3$.

\subsection{Extremos en un compacto}

%\documentclass{article}
%\usepackage{amsmath}

%\begin{document}

\noindent
Veamos ahora extremos en un compacto.

\subsubsection{Ejemplo 1: Extremos de \( f(x, y) = x^2 + y^2 + 2x + 1 \) en un rectángulo}

\paragraph{Dominio}
\[
-2 \leq x \leq 2, \quad -3 \leq y \leq 3.
\]

\paragraph{Paso 1: Puntos críticos en el interior}

Calculamos las derivadas parciales de \( f(x, y) \):
\[
\frac{\partial f}{\partial x} = 2x + 2, \quad \frac{\partial f}{\partial y} = 2y.
\]
Igualamos a cero para encontrar puntos críticos:
\[
2x + 2 = 0 \implies x = -1, \quad 2y = 0 \implies y = 0.
\]
El único punto crítico en el interior es \( (-1, 0) \). 

\paragraph*{Paso 2: Evaluar \( f(x, y) \) en los bordes}

\noindent
El borde del dominio es un rectángulo. Evaluamos \( f(x, y) \) en los cuatro lados:

\noindent
1. Lado \( x = -2 \):
   \[
   f(-2, y) = (-2)^2 + y^2 + 2(-2) + 1 = 4 + y^2 - 4 + 1 = y^2 + 1.
   \]
   Extremos en \( y = -3 \) y \( y = 3 \):
   \[
   f(-2, -3) = (-3)^2 + 1 = 10, \quad f(-2, 3) = 3^2 + 1 = 10.
   \]

\noindent
2. Lado \( x = 2 \):
   \[
   f(2, y) = 2^2 + y^2 + 2(2) + 1 = 4 + y^2 + 4 + 1 = y^2 + 9.
   \]
   Extremos en \( y = -3 \) y \( y = 3 \):
   \[
   f(2, -3) = (-3)^2 + 9 = 18, \quad f(2, 3) = 3^2 + 9 = 18.
   \]

\noindent
3. Lado \( y = -3 \):
   \[
   f(x, -3) = x^2 + (-3)^2 + 2x + 1 = x^2 + 9 + 2x + 1 = x^2 + 2x + 10.
   \]
   Extremos en \( x = -2 \) y \( x = 2 \):
   \[
   f(-2, -3) = (-2)^2 + 2(-2) + 10 = 4 - 4 + 10 = 10,
   \]
   \[
   f(2, -3) = 2^2 + 2(2) + 10 = 4 + 4 + 10 = 18.
   \]

\noindent
4. Lado \( y = 3 \):
   \[
   f(x, 3) = x^2 + 3^2 + 2x + 1 = x^2 + 9 + 2x + 1 = x^2 + 2x + 10.
   \]
   Extremos en \( x = -2 \) y \( x = 2 \):
   \[
   f(-2, 3) = (-2)^2 + 2(-2) + 10 = 4 - 4 + 10 = 10,
   \]
   \[
   f(2, 3) = 2^2 + 2(2) + 10 = 4 + 4 + 10 = 18.
   \]

\paragraph*{Paso 3: Evaluar en las esquinas}

\noindent
Las esquinas del rectángulo son:
\[
(-2, -3), (-2, 3), (2, -3), (2, 3).
\]
Evaluamos \( f(x, y) \) en cada una:
\[
f(-2, -3) = 10, \quad f(-2, 3) = 10, \quad f(2, -3) = 18, \quad f(2, 3) = 18.
\]

\paragraph*{Paso 4: Comparar valores}

En el interior: \( f(-1, 0) = (-1)^2 + 0^2 + 2(-1) + 1 = 1 - 2 + 1 = 0 \).
En los bordes: Los valores máximos son \( 18 \) y los mínimos son \( 0 \).

\paragraph*{Conclusión}

El \emph{mínimo} de \( f(x, y) \) en el dominio es \( 0 \) en \( (-1, 0) \), y el \emph{máximo} es \( 18 \) en \( (2, -3) \) y \( (2, 3) \).

%\end{document}

%\documentclass{article}
%\usepackage{amsmath}

%\begin{document}

%\chapter{Tmp}

\subsubsection*{Ejemplo 2: Extremos de \( f(x, y) = x^2 y \) en la elipse \( 3x^2 + 4y^2 \leq 12 \)}

\paragraph*{Dominio}
\[
3x^2 + 4y^2 \leq 12.
\]

\paragraph*{Paso 1: Puntos críticos en el interior}

Calculamos las derivadas parciales de \( f(x, y) \):
\[
\frac{\partial f}{\partial x} = 2xy, \quad \frac{\partial f}{\partial y} = x^2.
\]
Igualamos a cero:
\[
2xy = 0 \implies x = 0 \text{ o } y = 0,
\]
\[
x^2 = 0 \implies x = 0.
\]
El único punto crítico en el interior es \( (0, y) \), pero \( f(0, y) = 0 \) para cualquier \( y \).

\paragraph*{Paso 2: Parametrizar la frontera}

La frontera es la elipse \( 3x^2 + 4y^2 = 12 \). Parametrizamos usando coordenadas polares modificadas:
\[
x = 2\sqrt{2} \cos \theta, \quad y = \sqrt{3} \sin \theta, \quad \theta \in [0, 2\pi].
\]
Sustituimos en \( f(x, y) \):
\[
f(x, y) = (2\sqrt{2} \cos \theta)^2 (\sqrt{3} \sin \theta) = 8 \cos^2 \theta \cdot \sqrt{3} \sin \theta = 8\sqrt{3} \cos^2 \theta \sin \theta.
\]

\paragraph*{Paso 3: Encontrar extremos de \( f(\theta) \)}

Definimos \( g(\theta) = 8\sqrt{3} \cos^2 \theta \sin \theta \). Buscamos sus puntos críticos:
\[
g'(\theta) = 8\sqrt{3} [2 \cos \theta (-\sin \theta) \sin \theta + \cos^2 \theta \cos \theta] = 8\sqrt{3} \cos \theta [ -2 \sin^2 \theta + \cos^2 \theta ].
\]
Igualamos \( g'(\theta) = 0 \):
\[
\cos \theta = 0 \quad \text{o} \quad -2 \sin^2 \theta + \cos^2 \theta = 0.
\]
- Si \( \cos \theta = 0 \), entonces \( \theta = \frac{\pi}{2}, \frac{3\pi}{2} \), y \( f(x, y) = 0 \).
- Si \( -2 \sin^2 \theta + \cos^2 \theta = 0 \):
  \[
  -2 \sin^2 \theta + \cos^2 \theta = 0 \implies -2(1 - \cos^2 \theta) + \cos^2 \theta = 0 \implies -2 + 3\cos^2 \theta = 0 \implies \cos^2 \theta = \frac{2}{3}.
  \]
  Por lo tanto, \( \cos \theta = \pm \sqrt{\frac{2}{3}} \), y \( \sin \theta = \pm \sqrt{\frac{1}{3}} \).

Sustituyendo en \( f(x, y) \):
\[
f(x, y) = 8\sqrt{3} \left( \frac{2}{3} \right) \left( \pm \sqrt{\frac{1}{3}} \right) = \pm \frac{16\sqrt{3}}{3} \cdot \frac{1}{\sqrt{3}} = \pm \frac{16}{3}.
\]

\paragraph*{Paso 4: Evaluar en puntos críticos}

Los valores extremos en la frontera son \( \frac{16}{3} \) y \( -\frac{16}{3} \).

\paragraph*{Paso 5: Comparar con el interior}

En el interior, \( f(x, y) = 0 \). Por lo tanto, los extremos absolutos son \( \frac{16}{3} \) y \( -\frac{16}{3} \).

\paragraph*{Puntos críticos en la frontera}

Para \( \cos \theta = \sqrt{\frac{2}{3}} \) y \( \sin \theta = \sqrt{\frac{1}{3}} \):
\[
x = 2\sqrt{2} \cdot \sqrt{\frac{2}{3}} = \frac{4}{\sqrt{3}} = \frac{4\sqrt{3}}{3}, \quad y = \sqrt{3} \cdot \sqrt{\frac{1}{3}} = 1.
\]
Por lo tanto, los puntos críticos en la frontera son \( \left( \frac{4\sqrt{3}}{3}, 1 \right) \) y \( \left( -\frac{4\sqrt{3}}{3}, -1 \right) \).

\paragraph*{Conclusión}

El \emph{máximo} de \( f(x, y) \) en el dominio es \( \frac{16}{3} \) en \( \left( \frac{4\sqrt{3}}{3}, 1 \right) \), y el \emph{mínimo} es \( -\frac{16}{3} \) en \( \left( -\frac{4\sqrt{3}}{3}, -1 \right) \).

%Extremos Absolutos Clase 18
%Teorema en un compacto
%\noindent
%Por ejemplo 
%$x^2+y^2+2x+1$ en $-2 \leq  x \leq 2$,  $-3 \leq y \leq 3$. Se encuentra un minimo en $(-1,0)$ y un maximo en $(2,-3)$.

%\noindent 
%Otro ejemplo 
%$x^2y$ (simetrias) en el elipse $3x^2+4y^2 \leq 12$. Se parametriza el elipse. Puntos criticos en $$\sqrt{8}/3,1.$$

\newpage

\section{Practico 5 puntos criticos}

\begin{enumerate}

\item Para las siguientes funciones calcular los desarrollos de Taylor de orden 2 en el punto indicado.

a) $f(x,y)=\sen(x)\sen(y)$ en $(0,0)$  \ \ \  b) $f(x,y)=e^{x^2+y(y+1)}$ en $(0,0)$.

Y calcular los siguientes l\'imites:

a) $\underset{\left( x,y\right) \rightarrow \left( 0,0\right) }{\lim}\frac{xy-\sen(x)\sen(y)}{x^2+y^2}$\ \ 
b) $\underset{\left( x,y\right) \rightarrow \left( 0,0\right) }{\lim}\frac{e^{x^2+y(y+1)}-(1 + y)}{x^2+y^2}$.


\item
Para cada una de las siguientes funciones primero bosqueje el gráfico y prediga si es acotada, si alcanza un máximo (y un mínimo) en $\R^2$, así como si tiene puntos críticos. Luego, verifique hallando y clasificando los puntos críticos.
\begin{enumerate}
\item $f(x,y)=1+x^2 -y^2$.
\item $f(x,y)=x^2 + 5y^2$.
\item  $f(x,y)=\left( x^2+y^2 \right) \ e^{-(x^2+y^2)}$.

\end{enumerate}

\item
En los casos en que existan, hallar el m\'{a}ximo y el mínimo absolutos en todo $\R^2$
de cada una de las siguientes funciones:
\begin{enumerate}
\item  $f(x,y)=1-y^2-x^4$.
\item  $f(x,y)=\left( x^2-y^2 \right) \ e^{-(x^2+y^2)}$.
\end{enumerate}

\item  
Sean $f(x,y)=(3-x)(3-y)(x+y-3)$ y $g(x,y)=x^3 +y^3 -3xy$. Hallar todos los puntos cr\'\i ticos para cada una y
clasificarlos. ?`Tienen extremos absolutos en todo $\R^2$?

\item  
Verificar que la funci\'on $f(x,y,z)=x^4 + y^4 +z^4 -4xyz$ tiene un
punto crítico en $(1,1,1)$ y clasificarlo.% determinar la naturaleza de dicho punto.

\item 
Determinar los extremos absolutos y relativos, y los puntos de silla de la funci\'on
\begin{equation*}
f(x,y)=xy(1-x^2-y^2) \text{ en }\  [0,1]\times [0,1].
\end{equation*}


\end{enumerate}

\chapter{Funciones escalares V de V}

\section{Temario}

\begin{itemize}
%    \item Dominio, imágen, gráfico y conjuntos de nivel.
%    \item Límites y continuidad.
%    \item Derivadas (derivadas parciales/direccionales, diferenciabilidad, gradiente y espacio tangente).
%    \item Derivadas de orden superior (Taylor, criterio de la Hessiana).
    \item Extremos condicionados (Multiplicadores de Lagrange).
    \item Teorema de la Función Implícita.
\end{itemize}

\section{Notas}

%\subsection{Lagrange}

\noindent
Veamos primero como podemos extender el calculo de extremos cuando agregamos una restriccion.

\subsection{Extremos Condicionados (Multiplicadores de Lagrange)}

El método de los multiplicadores de Lagrange se utiliza para encontrar los máximos y mínimos de una función sujeta a restricciones. Para encontrar los extremos de \( f(x, y) \) sujeta a la restricción \( g(x, y) = 0 \), se resuelve el sistema de ecuaciones:
\[ \nabla f = \lambda \nabla g. \]
Esto es:
\[ \frac{\partial f}{\partial x} = \lambda \frac{\partial g}{\partial x}, \]
\[ \frac{\partial f}{\partial y} = \lambda \frac{\partial g}{\partial y}. \]
En el punto critico, tenemos que tener $\nabla g \neq 0$.
Este método es esencial en optimización y se aplica en problemas donde se busca maximizar o minimizar una función sujeta a ciertas condiciones.
%Mapa topografica y recorrido. Encontrar los extremos en el recorrido. Minimisar $f(x,y)$ bajo la condicion $g(x,y)=k$. Entonces
%\[
%\nabla f= \lambda \nabla g.
%\]

\subsubsection{Ejemplo 1}

Para resolver el problema de minimizar la función \( f(x, y) = xy \) sujeta a la restricción \( g(x, y) = x^2 + y^2 - 1 = 0 \), utilizamos el método de los multiplicadores de Lagrange.
%
Definimos el Lagrangiano:
\[ \mathcal{L}(x, y, \lambda) = xy - \lambda (x^2 + y^2 - 1) \]
%
Tomamos las derivadas parciales del Lagrangiano con respecto a \(x\), \(y\) y \(\lambda\), y las igualamos a cero:

\paragraph{1.} Derivada parcial con respecto a \(x\):
\[ \frac{\partial \mathcal{L}}{\partial x} = y - 2 \lambda x = 0 \]
\[ y = 2 \lambda x \]

\paragraph{2.} Derivada parcial con respecto a \(y\):
\[ \frac{\partial \mathcal{L}}{\partial y} = x - 2 \lambda y = 0 \]
\[ x = 2 \lambda y \]

\paragraph{3.} Derivada parcial con respecto a \(\lambda\):
\[ \frac{\partial \mathcal{L}}{\partial \lambda} = -(x^2 + y^2 - 1) = 0 \]
\[ x^2 + y^2 = 1 \]

Resolvemos el sistema de ecuaciones:

De \( y = 2 \lambda x \) y \( x = 2 \lambda y \), obtenemos:
\[ \lambda = \frac{y}{2x} \]
\[ \lambda = \frac{x}{2y} \]

Igualando las dos expresiones para \(\lambda\):
\[ \frac{y}{2x} = \frac{x}{2y} \]
\[ y^2 = x^2 \]
\[ y = \pm x \]

Consideramos dos casos:

\paragraph{Caso 1:} \( y = x \)

Sustituyendo en la restricción:
\[ x^2 + x^2 = 1 \]
\[ 2x^2 = 1 \]
\[ x^2 = \frac{1}{2} \]
\[ x = \pm \frac{\sqrt{2}}{2} \]

Por lo tanto, \( y = \pm \frac{\sqrt{2}}{2} \).

\paragraph{Caso 2:} \( y = -x \)

Sustituyendo en la restricción:
\[ x^2 + (-x)^2 = 1 \]
\[ 2x^2 = 1 \]
\[ x^2 = \frac{1}{2} \]
\[ x = \pm \frac{\sqrt{2}}{2} \]

Por lo tanto, \( y = \mp \frac{\sqrt{2}}{2} \).

Así, los puntos críticos son:
1. \( \left( \frac{\sqrt{2}}{2}, \frac{\sqrt{2}}{2} \right) \)
2. \( \left( -\frac{\sqrt{2}}{2}, -\frac{\sqrt{2}}{2} \right) \)
3. \( \left( \frac{\sqrt{2}}{2}, -\frac{\sqrt{2}}{2} \right) \)
4. \( \left( -\frac{\sqrt{2}}{2}, \frac{\sqrt{2}}{2} \right) \)

Evaluamos la función \( f(x, y) = xy \) en estos puntos:

1. \( f\left( \frac{\sqrt{2}}{2}, \frac{\sqrt{2}}{2} \right) = \frac{1}{2} \)
2. \( f\left( -\frac{\sqrt{2}}{2}, -\frac{\sqrt{2}}{2} \right) = \frac{1}{2} \)
3. \( f\left( \frac{\sqrt{2}}{2}, -\frac{\sqrt{2}}{2} \right) = -\frac{1}{2} \)
4. \( f\left( -\frac{\sqrt{2}}{2}, \frac{\sqrt{2}}{2} \right) = -\frac{1}{2} \)

%Los valores mínimos de la función son \(-\frac{1}{2}\), y se alcanzan en los puntos \( \left( \frac{\sqrt{2}}{2}, -\frac{\sqrt{2}}{2} \right) \) y \( \left( -\frac{\sqrt{2}}{2}, \frac{\sqrt{2}}{2} \right) \).

%Por lo tanto, las soluciones que minimizan la función son:
%\[ (x, y) = \left( \frac{\sqrt{2}}{2}, -\frac{\sqrt{2}}{2} \right) \]
%y
%\[ (x, y) = \left( -\frac{\sqrt{2}}{2}, \frac{\sqrt{2}}{2} \right) \]

%Sin embargo, la pista proporcionada sugiere que \( x = y = \frac{\sqrt{2}}{2} \) es una solución, lo cual corresponde a un máximo de la función. Si el problema era encontrar el máximo, entonces la solución sería:

%Los valores máximos de la función son \(\frac{1}{2}\), y se alcanzan en los puntos \( \left( \frac{\sqrt{2}}{2}, \frac{\sqrt{2}}{2} \right) \) y \( \left( -\frac{\sqrt{2}}{2}, -\frac{\sqrt{2}}{2} \right) \).

%¿Quisiste decir maximizar en lugar de minimizar?
\subsubsection{Ejemplo 2}

%Ejemplo: minimizar $f(x,y)=xy$ bajo la restriccion  $g(x,y)=x^2+y^2-1=0$ (Hint:Se encuentra $y= \lambda 2x$ y $x = \lambda 2y$. Solucion: Mas o menos $x=y=\sqrt{2}/2$.)
%19

\noindent
Para maximizar la función \( f(x, y, z) = x + y + z \) sujeta a la restricción \( g(x, y, z) = x^2 + y^2 + z^2 - 1 = 0 \), utilizamos el método de los multiplicadores de Lagrange.
%
Definimos el Lagrangiano:
\[ \mathcal{L}(x, y, z, \lambda) = x + y + z - \lambda (x^2 + y^2 + z^2 - 1) \]
%
Tomamos las derivadas parciales del Lagrangiano con respecto a \(x\), \(y\), \(z\) y \(\lambda\), y las igualamos a cero:

\paragraph{1. Derivada parcial con respecto a \(x\):}
\[ \frac{\partial \mathcal{L}}{\partial x} = 1 - 2 \lambda x = 0 \]
\[ x = \frac{1}{2 \lambda} \]

\paragraph{2. Derivada parcial con respecto a \(y\):}
\[ \frac{\partial \mathcal{L}}{\partial y} = 1 - 2 \lambda y = 0 \]
\[ y = \frac{1}{2 \lambda} \]

\paragraph{3. Derivada parcial con respecto a \(z\):}
\[ \frac{\partial \mathcal{L}}{\partial z} = 1 - 2 \lambda z = 0 \]
\[ z = \frac{1}{2 \lambda} \]

4. Derivada parcial con respecto a \(\lambda\):
\[ \frac{\partial \mathcal{L}}{\partial \lambda} = -(x^2 + y^2 + z^2 - 1) = 0 \]
\[ x^2 + y^2 + z^2 = 1 \]

Sustituyendo \( x = y = z = \frac{1}{2 \lambda} \) en la ecuación de la restricción:
\[ \left( \frac{1}{2 \lambda} \right)^2 + \left( \frac{1}{2 \lambda} \right)^2 + \left( \frac{1}{2 \lambda} \right)^2 = 1 \]
\[ 3 \left( \frac{1}{2 \lambda} \right)^2 = 1 \]
\[ 3 = 4 \lambda^2 \]
\[ \lambda = \frac{\sqrt{3}}{2} \]

Por lo tanto, el punto que maximiza la función es:
\[ (x, y, z) = \left( \frac{\sqrt{3}}{3}, \frac{\sqrt{3}}{3}, \frac{\sqrt{3}}{3} \right) \]

El valor máximo de la función es:
\[ f\left( \frac{\sqrt{3}}{3}, \frac{\sqrt{3}}{3}, \frac{\sqrt{3}}{3} \right) = \sqrt{3} \]


%$f=x+y+z$ y $g=x^2+y^2+z^2=1$ Max en $\sqrt{3}/3$

\subsubsection{Ejemplo 3}
%***
Para maximizar la función \( f(x, y, z) = xyz \) sujeta a la restricción \( g(x, y, z) = x + y + z - 100 = 0 \), utilizamos el método de los multiplicadores de Lagrange.
%
Definimos el Lagrangiano:
\[ \mathcal{L}(x, y, z, \lambda) = xyz - \lambda (x + y + z - 100) \]
%
Tomamos las derivadas parciales del Lagrangiano con respecto a \(x\), \(y\), \(z\) y \(\lambda\), y las igualamos a cero:

\paragraph{1. Derivada parcial con respecto a \(x\):}
\[ \frac{\partial \mathcal{L}}{\partial x} = yz - \lambda = 0 \]
\[ yz = \lambda \]

\paragraph{2. Derivada parcial con respecto a \(y\):}
\[ \frac{\partial \mathcal{L}}{\partial y} = xz - \lambda = 0 \]
\[ xz = \lambda \]

\paragraph{3. Derivada parcial con respecto a \(z\):}
\[ \frac{\partial \mathcal{L}}{\partial z} = xy - \lambda = 0 \]
\[ xy = \lambda \]

\paragraph{4. Derivada parcial con respecto a \(\lambda\):}
\[ \frac{\partial \mathcal{L}}{\partial \lambda} = -(x + y + z - 100) = 0 \]
\[ x + y + z = 100 \]

De las primeras tres ecuaciones, obtenemos:
\[ yz = xz = xy \]

De esto, deducimos que:
\[ x = y = z \]

Sustituyendo \( x = y = z \) en la ecuación de la restricción:
\[ x + y + z = 100 \]
\[ 3x = 100 \]
\[ x = \frac{100}{3} \]

Por lo tanto, el punto que maximiza la función es:
\[ (x, y, z) = \left( \frac{100}{3}, \frac{100}{3}, \frac{100}{3} \right) \]

El valor máximo de la función es:
\[ f\left( \frac{100}{3}, \frac{100}{3}, \frac{100}{3} \right) = \left( \frac{100}{3} \right)^3 \]

\subsection{Dos condiciones}

\noindent
Admitimos el resultado siguiente.

\begin{teorema}
Sea \( f: \mathbb{R}^3 \to \mathbb{R} \) una función continua y diferenciable, y sean \( g_1, g_2: \mathbb{R}^3 \to \mathbb{R} \) funciones continuas y diferenciables. Supongamos que \( \nabla g_1 \) y \( \nabla g_2 \) son linealmente independientes en los puntos donde \( g_1(x, y, z) = 0 \) y \( g_2(x, y, z) = 0 \).
%
Si \( (x_0, y_0, z_0) \) es un punto donde \( f \) tiene un extremo local sujeto a las restricciones \( g_1(x, y, z) = 0 \) y \( g_2(x, y, z) = 0 \), entonces existen escalares \( \lambda_1 \) y \( \lambda_2 \) tales que:
%
\[ \nabla f(x_0, y_0, z_0) = \lambda_1 \nabla g_1(x_0, y_0, z_0) + \lambda_2 \nabla g_2(x_0, y_0, z_0) \]
\end{teorema}

\subsubsection{Ejemplo}
%**Ejemplo:**

Consideremos la función \( f(x, y, z) = z^3 - 3x^2 \) sujeta a las restricciones \( g_1(x, y, z) = x^2 + y^2 - 1 = 0 \) y \( g_2(x, y, z) = x + z - 1 = 0 \).
%
%**Solución:**
%
Definimos el Lagrangiano:
\[ \mathcal{L}(x, y, z, \lambda_1, \lambda_2) = z^3 - 3x^2 - \lambda_1 (x^2 + y^2 - 1) - \lambda_2 (x + z - 1) \]
%
Tomamos las derivadas parciales del Lagrangiano con respecto a \(x\), \(y\), \(z\), \(\lambda_1\) y \(\lambda_2\), y las igualamos a cero:

\paragraph{1. Derivada parcial con respecto a \(x\):}
\[ \frac{\partial \mathcal{L}}{\partial x} = -6x - 2\lambda_1 x - \lambda_2 = 0 \]
\[ -6x - 2\lambda_1 x - \lambda_2 = 0 \quad (1) \]

\paragraph{2. Derivada parcial con respecto a \(y\):}
\[ \frac{\partial \mathcal{L}}{\partial y} = -2\lambda_1 y = 0 \quad (2) \]

\paragraph{3. Derivada parcial con respecto a \(z\):}
\[ \frac{\partial \mathcal{L}}{\partial z} = 3z^2 - \lambda_2 = 0 \quad (3) \]

\paragraph{4. Derivada parcial con respecto a \(\lambda_1\):}
\[ \frac{\partial \mathcal{L}}{\partial \lambda_1} = -(x^2 + y^2 - 1) = 0 \]
\[ x^2 + y^2 = 1 \quad (4) \]

\paragraph{5. Derivada parcial con respecto a \(\lambda_2\):}
\[ \frac{\partial \mathcal{L}}{\partial \lambda_2} = -(x + z - 1) = 0 \]
\[ x + z = 1 \quad (5) \]

De la ecuación (2):
\[ -2\lambda_1 y = 0 \]

Esto implica que \(\lambda_1 = 0\) o \(y = 0\).

\paragraph{Caso 1:} \(\lambda_1 = 0\)

De la ecuación (1):
\[ -6x - \lambda_2 = 0 \]
\[ \lambda_2 = -6x \]

De la ecuación (3):
\[ 3z^2 - \lambda_2 = 0 \]
\[ 3z^2 = -6x \]
\[ z^2 = -2x \]

De la ecuación (5):
\[ x + z = 1 \]
\[ z = 1 - x \]

Sustituyendo \( z = 1 - x \) en \( z^2 = -2x \):
\[ (1 - x)^2 = -2x \]
\[ 1 - 2x + x^2 = -2x \]
\[ 1 + x^2 = 0 \]

Esta ecuación no tiene soluciones reales.

\paragraph{Caso 2:} \( y = 0 \)

De la ecuación (4):
\[ x^2 + 0^2 = 1 \]
\[ x^2 = 1 \]
\[ x = \pm 1 \]

De la ecuación (5):
\[ x + z = 1 \]
\[ z = 1 - x \]

Sustituyendo \( x = 1 \) y \( x = -1 \):

- Si \( x = 1 \):
\[ z = 1 - 1 = 0 \]

De la ecuación (3):
\[ 3(0)^2 - \lambda_2 = 0 \]
\[ \lambda_2 = 0 \]

De la ecuación (1):
\[ -6(1) - 2\lambda_1 (1) - 0 = 0 \]
\[ -6 - 2\lambda_1 = 0 \]
\[ \lambda_1 = -3 \]

- Si \( x = -1 \):
\[ z = 1 - (-1) = 2 \]

De la ecuación (3):
\[ 3(2)^2 - \lambda_2 = 0 \]
\[ 12 - \lambda_2 = 0 \]
\[ \lambda_2 = 12 \]

De la ecuación (1):
\[ -6(-1) - 2\lambda_1 (-1) - 12 = 0 \]
\[ 6 + 2\lambda_1 - 12 = 0 \]
\[ 2\lambda_1 = 6 \]
\[ \lambda_1 = 3 \]

Por lo tanto, los puntos críticos son:
1. \( (x, y, z) = (1, 0, 0) \)
2. \( (x, y, z) = (-1, 0, 2) \)

Evaluamos la función \( f(x, y, z) = z^3 - 3x^2 \) en estos puntos:

1. Para \( (1, 0, 0) \):
\[ f(1, 0, 0) = 0^3 - 3(1)^2 = -3 \]

2. Para \( (-1, 0, 2) \):
\[ f(-1, 0, 2) = 2^3 - 3(-1)^2 = 8 - 3 = 5 \]

Por lo tanto, los puntos críticos son \( (1, 0, 0) \) y \( (-1, 0, 2) \).

%\chapter{Funciones escalares Vbis de V}

%\newpage

%\section{Notas}
%

\subsection{Funcion Implicita}

\noindent 
Detallamos un poco del Teorema de la Función Implicita.

%\subsubsection{Teorema de la Función Implícita}

El teorema de la función implícita proporciona condiciones bajo las cuales una ecuación \( F(x, y) = 0 \) define implícitamente una función \( y = f(x) \). Si \( F(x, y) = 0 \) y \( \frac{\partial F}{\partial y} \neq 0 \), entonces se puede definir \( y \) como una función de \( x \), y la derivada de \( y \) con respecto a \( x \) es:
\[ \frac{dy}{dx} = -\frac{\frac{\partial F}{\partial x}}{\frac{\partial F}{\partial y}}. \]
Este teorema es fundamental en el estudio de curvas y superficies definidas implícitamente. Se puede extender a la dimensión 3.

\begin{example} Consideramos la función 
    $F(x,y)=x^2+y^2-1$ en el punto $p=(\sqrt{2}/2,\sqrt{2}/2)$. Con el teorema, obtenemos $g'(x)=-x/y$, y$g'(\sqrt{2}/2=-1)$.
\end{example}

\begin{example} Sea también la función
 $F(x,y) = x^2 + xy^2 -5$, con $p=(1,2)$. Obtenemos $g'(1)=-5/4$.   
\end{example}

\noindent
La fórmula general para la derivada de segundo orden de una función implícita \( F(x, y) = 0 \) es:
\[
\frac{d^2y}{dx^2} = -\frac{\frac{\partial^2 F}{\partial x^2} \left( \frac{\partial F}{\partial y} \right)^2 - 2 \frac{\partial F}{\partial x} \frac{\partial^2 F}{\partial x \partial y} \frac{\partial F}{\partial y} + \frac{\partial F}{\partial x}^2 \frac{\partial^2 F}{\partial y^2}}{\left( \frac{\partial F}{\partial y} \right)^3}
\]
Ejemplo: $F(x,y)=x^2+y^2-1$.
Diferenciando implícitamente con respecto a \( x \):
1. Primera derivada:
\[
2x + 2y \frac{dy}{dx} = 0 \implies \frac{dy}{dx} = -\frac{x}{y}.
\]
2. Segunda derivada:
\[
\frac{d^2y}{dx^2} = -\frac{y^2 + x^2}{y^3} = -\frac{1}{y^3}.
\]
Sea \( F(x, y, z) = 0 \) una función implícita. Si \( F \) es diferenciable y \( \frac{\partial F}{\partial z} \neq 0 \) en un punto \( (x_0, y_0, z_0) \), entonces existe una función \( z = f(x, y) \) tal que \( F(x, y, f(x, y)) = 0 \) en un entorno de \( (x_0, y_0, z_0) \). Las derivadas parciales de \( z \) con respecto a \( x \) y \( y \) pueden encontrarse utilizando las siguientes fórmulas:
\[
\frac{\partial z}{\partial x} = -\frac{\frac{\partial F}{\partial x}}{\frac{\partial F}{\partial z}}
\]
y también
\[
\frac{\partial z}{\partial y} = -\frac{\frac{\partial F}{\partial y}}{\frac{\partial F}{\partial z}}
\]
Además, si \( F \) es dos veces diferenciable, las segundas derivadas parciales pueden encontrarse utilizando:
\[
\frac{\partial^2 z}{\partial x^2} = -\frac{\frac{\partial^2 F}{\partial x^2} \left( \frac{\partial F}{\partial z} \right)^2 - 2 \frac{\partial F}{\partial x} \frac{\partial^2 F}{\partial x \partial z} \frac{\partial F}{\partial z} + \frac{\partial F}{\partial x}^2 \frac{\partial^2 F}{\partial z^2}}{\left( \frac{\partial F}{\partial z} \right)^3}
\]
ademas de 
\[
\frac{\partial^2 z}{\partial y^2} = -\frac{\frac{\partial^2 F}{\partial y^2} \left( \frac{\partial F}{\partial z} \right)^2 - 2 \frac{\partial F}{\partial y} \frac{\partial^2 F}{\partial y \partial z} \frac{\partial F}{\partial z} + \frac{\partial F}{\partial y}^2 \frac{\partial^2 F}{\partial z^2}}{\left( \frac{\partial F}{\partial z} \right)^3}
\]
y de
\[
\frac{\partial^2 z}{\partial x \partial y} = -\frac{\frac{\partial^2 F}{\partial x \partial y} \left( \frac{\partial F}{\partial z} \right)^2 - \frac{\partial F}{\partial x} \frac{\partial^2 F}{\partial x \partial z} \frac{\partial F}{\partial y} - \frac{\partial F}{\partial y} \frac{\partial^2 F}{\partial y \partial z} \frac{\partial F}{\partial x} + \frac{\partial F}{\partial x} \frac{\partial F}{\partial y} \frac{\partial^2 F}{\partial z^2}}{\left( \frac{\partial F}{\partial z} \right)^3}
\]

%\chapter{TEMP}

%\begin{example}
    
%\end{example}

\newpage

\section{Practico 6 (Lagrange y implicito)}

\paragraph{Multiplicadores de Lagrange}

\begin{enumerate}

\item
En los casos que siguen, hallar los extremos absolutos de la función $f:M\to\R$.
\begin{enumerate}
\item
$M$ es la circunferencia de  radio unidad y centro en el origen de $\R^2$, $f(x,y)=(x+y)e^{x+y}$.
\item
$M$ es  la esfera de radio unidad y centro en el origen de $\R^3$, $f(x,y,z)=ax+by+cz$, siendo $a,b,c\in \R$ no todos nulos.
\item
$M=\{(x,y,z)\in\R^3:2x-y-z=2\text{ y }x^2+y^2=1\}$, $f(x,y,z)=4y-2z$.
\end{enumerate}

\item
Hallar la m\'axima y mínima distancia desde el origen a la curva 
$5x^2 +6xy +5y^2=8$. 
Hallar adem\'as los puntos en los que esta curva tiene tangentes horizontales y verticales.
 

\item 
Hallar los puntos de la superficie $z^2-xy=1$ que est\'an m\'as pr\'oximos al origen.
 
\item
Calcular la base y altura del rectángulo de mayor área inscripto en la elipse $\frac{x^2}9+\frac{y^2}4=1$.


\item
Sea $f:\R^2\to\R$ definida por $f(x,y)= x^2+ y^2-xy$.
\vskip1mm\noindent
(a) Hallar los puntos críticos de $f$ en $\R^2$.
\vskip1mm\noindent

(b) Hallar los extremos absolutos de $f$ en el conjunto $\{(x,y)\in\R^2:\ x^2+y^2+xy\leq 1\}$.
 

\item
Hallar extremos absolutos de $f$ en $D$, siendo:
\begin{itemize}
\item[(a)]$D=\{ (x,y)\colon 5x^2 -6xy +5y^2 \leq 4 \}$, $f(x,y)=xy$.
\item[(b)] $D=\{ (x,y)\colon 2x^2 +y^2 \leq 1\}$, $f(x,y)=e^{(x-1)^2+y^2}$.
\end{itemize}

\item Hallar los puntos en la curva de intersección de las dos superficies
$x^2-xy+y^2-z^2=1$ y $x^2+y^2=1$ que están más cerca del origen.
\end{enumerate}

\paragraph{Funci\'on implicita}

\begin{enumerate} %[start=8]
\item
Probar que las siguientes ecuaciones determinan $y=f(x)$ en un entorno del punto $(a,b)$.
%que se indica en cada caso. 
Calcular $f'(a)$ y $f''(a)$.
\begin{itemize}
\item[(a)] $x^{2}y+\log(xy)=1$ , $(a,b)=(1,1)$.
\item[(b)] $y+\sinh(y)-\sen(x)=0$ , $(a,b)=(0,0)$. %Hallar $f'''(x_{0})$.
%\item[(c)] $x^2/a^2+y^2/b^2=1$, en $(x_0,y_0)$ gen\'erico con $x^2_{0}/a^2+y^2_{0}/b^2=1$, $y_0\neq 0 $.
%Verificar lo anterior despejando $y$ expl\'\i citamente.
\end{itemize}

\item Dada la ecuaci\'on $(x^2+y^2+2z^2)^{\frac{1}{2}}=x\cos(z)$, puede determinarse $z$ en funci\'on de $x,y$ alrededor del punto $(1,0,0)$ ?
De ser as\'i, determinar el primer desarrollo de Taylor para dicha funci\'on impl\'icita. 
\end{enumerate}

%\section*{Conceptos Avanzados de Cálculo y Análisis Matemático}

%\subsection{Derivadas de Orden Superior}

%\subsubsection{Serie de Taylor}

%La serie de Taylor es una representación de una función como una suma infinita de términos calculados a partir de los valores de sus derivadas en un punto \( a \). La serie de Taylor de una función \( f(x) \) alrededor de un punto \( a \) es:
%\[ f(x) = f(a) + f'(a)(x-a) + \frac{f''(a)}{2!}(x-a)^2 + \frac{f'''(a)}{3!}(x-a)^3 + \cdots \]
%Esta serie es útil para aproximar funciones y es ampliamente utilizada en análisis numérico y teoría de aproximación.

%\subsubsection{Criterio de la Hessiana}

%El criterio de la Hessiana se utiliza para determinar la naturaleza de los puntos críticos de una función de varias variables. Para una función \( f(x, y) \), la matriz Hessiana \( H \) es:
%\[ H = \begin{bmatrix}
%f_{xx} & f_{xy} \\
%f_{yx} & f_{yy}
%\end{bmatrix} \]

%El criterio de la Hessiana establece que:

%\begin{itemize}
%    \item Si el determinante de \( H \) es positivo y \( f_{xx} > 0 \), entonces el punto crítico es un mínimo local.
%    \item Si el determinante es positivo y \( f_{xx} < 0 \), entonces es un máximo local.
%    \item Si el determinante es negativo, el punto es un punto de silla.
%\end{itemize}


%\subsection{Extremos Condicionados (Multiplicadores de Lagrange)}

%El método de los multiplicadores de Lagrange se utiliza para encontrar los máximos y mínimos de una función sujeta a restricciones. Para encontrar los extremos de \( f(x, y) \) sujeta a la restricción \( g(x, y) = 0 \), se resuelve el sistema de ecuaciones:
%\[ \nabla f = \lambda \nabla g \]
%Esto es:
%\[ \frac{\partial f}{\partial x} = \lambda \frac{\partial g}{\partial x} \]
%\[ \frac{\partial f}{\partial y} = \lambda \frac{\partial g}{\partial y} \]
%Este método es esencial en optimización y se aplica en problemas donde se busca maximizar o minimizar una función sujeta a ciertas condiciones.

\chapter{Campos vectoriales}

\section{Temario}

\begin{itemize}
    \item Continuidad y diferenciabilidad. 
    \item Diferencial y matriz Jacobiana.
    \item Cambios de variable y determinante Jacobiano.
    \item Teorema de la Función Inversa.
\end{itemize}

\section{Notas}

\noindent Veamos las distintas nociones asociadas a un campo vectorial. Empezamos con definicion y ejemplos.

\subsection{Campo Vectorial}

\noindent
Un campo vectorial \( \mathbf{F} \) en \( \mathbb{R}^n \) es una función que asigna a cada punto \( (x_1, x_2, \ldots, x_n) \) un vector \( \mathbf{F}(x_1, x_2, \ldots, x_n) \).


Un campo de gradiente es un campo vectorial que es el gradiente de alguna función escalar \( f \). Es decir, si \( f: \mathbb{R}^n \to \mathbb{R} \) es una función escalar, entonces el campo de gradiente \( \mathbf{F} \) se define como:
\[ \mathbf{F} = \nabla f = \left( \frac{\partial f}{\partial x_1}, \frac{\partial f}{\partial x_2}, \ldots, \frac{\partial f}{\partial x_n} \right). \]

\paragraph{Ejemplo con la función paraboloide.} %\( f(x, y) = x^2 + y^2 \)
%.}
Consideremos la función \( f(x, y) = x^2 + y^2 \). El gradiente de \( f \) es:
\[ \nabla f = \left( \frac{\partial f}{\partial x}, \frac{\partial f}{\partial y} \right) = (2x, 2y). \]

\paragraph{Ejemplo en Variable Compleja.}
Consideremos la función \( f(z) = z^2 \). Si \( z = x + iy \), entonces:
\[ f(z) = (x + iy)^2 = x^2 - y^2 + i(2xy). \]
Por lo tanto, las partes real e imaginaria son:
\[ \text{Re}(f(z)) = x^2 - y^2 \]
\[ \text{Im}(f(z)) = 2xy \]
Ver \url{https://en.wikipedia.org/wiki/Mandelbrot_set}.

\paragraph{Otro ejemplo en Variable Compleja.}
Para la función \( g(z) = \exp(i\theta) z \), donde \( z = x + iy \), tenemos:
\[ g(z) = \exp(i\theta) (x + iy) = (x \cos \theta - y \sin \theta) + i(x \sin \theta + y \cos \theta) \]
Por lo tanto, las partes real e imaginaria son:
\[ \text{Re}(g(z)) = x \cos \theta - y \sin \theta, \]
\[ \text{Im}(g(z)) = x \sin \theta + y \cos \theta. \]
Corresponde a una rotacion de angulo $\theta$.

\subsection{Continuidad y Diferenciabilidad}

Un campo vectorial \(\mathbf{F}: \mathbb{R}^n \to \mathbb{R}^m\) es una función que asigna un vector de \(\mathbb{R}^m\) a cada punto en \(\mathbb{R}^n\). Un campo vectorial es continuo en un punto \(\mathbf{a}\) si:
\[
\lim_{\mathbf{x} \to \mathbf{a}} \mathbf{F}(\mathbf{x}) = \mathbf{F}(\mathbf{a}).
\]
Además, \(\mathbf{F}\) es diferenciable en \(\mathbf{a}\) si existe una transformación lineal \(D\mathbf{F}(\mathbf{a})\) tal que:
\[
\mathbf{F}(\mathbf{x}) = \mathbf{F}(\mathbf{a}) + D\mathbf{F}(\mathbf{a})(\mathbf{x} - \mathbf{a}) + \|\mathbf{x} - \mathbf{a}\| \mathbf{E}(\mathbf{x}, \mathbf{a})
\]
donde \(\mathbf{E}(\mathbf{x}, \mathbf{a}) \to 0\) cuando \(\mathbf{x} \to \mathbf{a}\).

\subsection{Diferencial y Matriz Jacobiana}

La diferencial de un campo vectorial \(\mathbf{F}\) en un punto \(\mathbf{a}\) está representada por la matriz Jacobiana \(J_{\mathbf{F}}(\mathbf{a})\), que es una matriz \(m \times n\) de las derivadas parciales de \(\mathbf{F}\):
\[
J_{\mathbf{F}} = \begin{bmatrix}
\frac{\partial f_1}{\partial x_1} & \frac{\partial f_1}{\partial x_2} & \cdots & \frac{\partial f_1}{\partial x_n} \\
\frac{\partial f_2}{\partial x_1} & \frac{\partial f_2}{\partial x_2} & \cdots & \frac{\partial f_2}{\partial x_n} \\
\vdots & \vdots & \ddots & \vdots \\
\frac{\partial f_m}{\partial x_1} & \frac{\partial f_m}{\partial x_2} & \cdots & \frac{\partial f_m}{\partial x_n}
\end{bmatrix}
\]
La matriz Jacobiana es fundamental para entender cómo se transforma un campo vectorial bajo cambios de coordenadas.


\paragraph{Jacobiana en Dimensión 2}
La matriz jacobiana de una función vectorial \( \mathbf{F}: \mathbb{R}^2 \to \mathbb{R}^2 \) se define como:
\[ J = \begin{pmatrix} \frac{\partial F_1}{\partial x} & \frac{\partial F_1}{\partial y} \\ \frac{\partial F_2}{\partial x} & \frac{\partial F_2}{\partial y} \end{pmatrix}. \]

\paragraph{Jacobiana para Coordenadas Polares}
Consideremos la transformación de coordenadas polares \( (r, \theta) \) a coordenadas cartesianas \( (x, y) \):
\[ x = r \cos \theta, \]
\[ y = r \sin \theta. \]
La matriz jacobiana es:
\[ J = \begin{pmatrix} \cos \theta & -r \sin \theta \\ \sin \theta & r \cos \theta \end{pmatrix}. \]

\paragraph{Jacobiana de \( (xy, \exp(z + x)) \)}
Consideremos la función \( \mathbf{F}(x, y, z) = (xy, \exp(z + x)) \). La matriz jacobiana es:
\[ J = \begin{pmatrix} y & x & 0 \\ \exp(z + x) & 0 & \exp(z + x) \end{pmatrix}. \]

\paragraph{Jacobiana de \( (x + y, x - y) \).}
Consideremos la función \( \mathbf{F}(x, y) = (x + y, x - y) \). La matriz jacobiana es:
\[ J = \begin{pmatrix} 1 & 1 \\ 1 & -1 \end{pmatrix}. \]

\paragraph{Regla de la Cadena Vectorial}
La regla de la cadena para funciones vectoriales establece que si \( \mathbf{F}: \mathbb{R}^n \to \mathbb{R}^m \) y \( \mathbf{G}: \mathbb{R}^m \to \mathbb{R}^p \) son funciones diferenciables, entonces la derivada de la composición \( \mathbf{G} \circ \mathbf{F} \) está dada por:
\[ D(\mathbf{G} \circ \mathbf{F}) = D\mathbf{G}(\mathbf{F}) \cdot D\mathbf{F}. \]

\paragraph{Ejemplo con \( f(x, y) = xy \) y \( G(r, \theta) = (r \cos \theta, r \sin \theta) \).}
Consideremos \( f(x, y) = xy \) y \( G(r, \theta) = (r \cos \theta, r \sin \theta) \). La composición \( f(G(r, \theta)) \) es:
\[ f(G(r, \theta)) = (r \cos \theta)(r \sin \theta) = r^2 \cos \theta \sin \theta. \]
La derivada de \( f \) con respecto a \( r \) y \( \theta \) es:
\[ \frac{\partial f}{\partial r} = 2r \cos \theta \sin \theta, \]
\[ \frac{\partial f}{\partial \theta} = r^2 (\cos^2 \theta - \sin^2 \theta). \]

%\section{Ejemplo de Composición de Funciones}
%Consideremos \( F(x, y, z) = (xyz, x + y) \) y \( G(s, t) = (st, s + t, s^2 - t^2) \). Evaluamos en \( s = t = 1 \):
%\[ G(1, 1) = (1 \cdot 1, 1 + 1, 1^2 - 1^2) = (1, 2, 0) \]
%Luego, \[ F(G(1, 1)) = F(1, 2, 0) = (1 \cdot 2 \cdot 0, 1 + 2) = (0, 3) \]

\paragraph{Otro ejemplo:} 
$F(x,y,z)=(xyz,x+y)$ y $ G(s,t) = (st,s+t,s^2-t^2)$ (en $s=t=1$).
%
%\documentclass{article}
%\usepackage{amsmath}
%
%\begin{document}
%\section*
{Jacobiana de \( F(x,y,z) = (xyz, x+y) \)}
\[
J_F = \begin{pmatrix}
yz & xz & xy \\
1 & 1 & 0
\end{pmatrix}.
\]
%\section*
{Jacobiana de \( G(s,t) = (st, s+t, s^2-t^2) \)}
\[
J_G = \begin{pmatrix}
t & s \\
1 & 1 \\
2s & -2t
\end{pmatrix}.
\]
%\section*
{Jacobiana compuesta \( (F \circ G)(1,1) \)}
\[
J_{F \circ G} = \begin{pmatrix}
4 & -4 \\
2 & 2
\end{pmatrix}.
\]

%\end{document}

\subsection{Difeomorfismo}
\noindent
Un difeomorfismo es una función diferenciable \( \mathbf{F}: \mathbb{R}^n \to \mathbb{R}^n \) que tiene una inversa diferenciable \( \mathbf{F}^{-1} \).

\paragraph{Jacobiano del Inverso de un Difeomorfismo}
Si \( \mathbf{F} \) es un difeomorfismo, entonces el jacobiano de \( \mathbf{F}^{-1} \) es la inversa del jacobiano de \( \mathbf{F} \):
\[ D\mathbf{F}^{-1}(\mathbf{F}(x)) = (D\mathbf{F}(x))^{-1}. \]

%\documentclass{article}
%\usepackage{amsmath}

%\begin{document}
\paragraph{Jacobiana del cambio a coordenadas polares \( F(r, \theta) = (r \cos \theta, r \sin \theta) \)}
\[
J_F = \begin{pmatrix}
\cos \theta & -r \sin \theta \\
\sin \theta & r \cos \theta
\end{pmatrix}
\]

\paragraph{Jacobiana del cambio inverso \( F^{-1}(x, y) = \left( \sqrt{x^2 + y^2}, \arctan\left(\frac{y}{x}\right) \right) \)}
\[
J_{F^{-1}} = \begin{pmatrix}
\frac{x}{\sqrt{x^2 + y^2}} & \frac{y}{\sqrt{x^2 + y^2}} \\
-\frac{y}{x^2 + y^2} & \frac{x}{x^2 + y^2}
\end{pmatrix}
\]

%\end{document}

\paragraph{Otro ejemplo.} Cambio de variable parabolico "$z^2$":
$x=u^2-v^2$ y $y=2uv$.
Jacobiano
\[
J = \begin{bmatrix}
2u & -2v \\
2v & 2u
\end{bmatrix}.
\]
Determinante del jacobiano
\[
\det(J) = 4(u^2 + v^2).
\]
Inversa del jacobiano
\[
J^{-1} = \frac{1}{2(u^2 + v^2)} \begin{bmatrix}
u & v \\
-v & u
\end{bmatrix}.
\]

%**Teorema de la Función Inversa (versión rioplatense):**

\subsection{Teorema de la Funcion Inversa}

\noindent
Veamos lo siguiente:

\begin{teorema}
Supongamos que tenés una función \( F: \mathbb{R}^n \to \mathbb{R}^n \) que es \( \mathcal{C}^1 \) (o sea, sus derivadas parciales existen y son continuas) en un conjunto abierto \( U \subseteq \mathbb{R}^n \). Si en un punto \( a \in U \) el \emph{determinante del jacobiano} de \( F \) no es cero:
\[
\det(J_F(a)) \neq 0,
\]
Entonces:
\begin{enumerate}
    \item \emph{Existencia de inversa local:} Hay un entorno \( V \) alrededor de \( a \) y otro \( W \) alrededor de \( F(a) \) donde \( F \) es biyectiva (o sea, tiene inversa).
\item \emph{La inversa es \( C^1 \):} La función inversa \( F^{-1}: W \to V \) también es continuamente diferenciable.
\item \emph{Relación entre jacobianas:} Para todo \( y \in W \), vale que:
   \[
   J_{F^{-1}}(y) = \left( J_F(F^{-1}(y)) \right)^{-1}.
   \]
\end{enumerate}
\end{teorema}
\noindent
Informalmente: si el jacobiano no se anula, cerca del punto \( a \) la función \( F \) tiene una inversa \emph{bien portada} (diferenciable y continua).

\subsection{Coordenadas Esféricas}

En \(\mathbb{R}^3\), las coordenadas esféricas \((r, \theta, \phi)\) se definen así:
- \( r \): distancia al origen (\( r \geq 0 \)),
- \( \theta \): ángulo en el plano \(xy\) desde el eje \(x\) (o sea, el ángulo polar, \(0 \leq \theta < 2\pi\)),
- \( \phi \): ángulo desde el eje \(z\) (o sea, el ángulo azimutal, \(0 \leq \phi \leq \pi\)).
La relación con las coordenadas cartesianas \((x, y, z)\) es:
\[
\begin{cases}
x = r \sin \phi \cos \theta, \\
y = r \sin \phi \sin \theta, \\
z = r \cos \phi.
\end{cases}
\]
Podemos introducir $\rho = r \sin \phi$.
\paragraph{Jacobiano de la Transformación}
El jacobiano \( J \) es la matriz de derivadas parciales de \((x, y, z)\) respecto de \((r, \theta, \phi)\):
\[
J = \begin{bmatrix}
\frac{\partial x}{\partial r} & \frac{\partial x}{\partial \theta} & \frac{\partial x}{\partial \phi} \\
\frac{\partial y}{\partial r} & \frac{\partial y}{\partial \theta} & \frac{\partial y}{\partial \phi} \\
\frac{\partial z}{\partial r} & \frac{\partial z}{\partial \theta} & \frac{\partial z}{\partial \phi}
\end{bmatrix}
\]
Calculando cada derivada:
\[
\frac{\partial x}{\partial r} = \sin \phi \cos \theta, \quad \frac{\partial x}{\partial \theta} = -r \sin \phi \sin \theta, \quad \frac{\partial x}{\partial \phi} = r \cos \phi \cos \theta,
\]
\[
\frac{\partial y}{\partial r} = \sin \phi \sin \theta, \quad \frac{\partial y}{\partial \theta} = r \sin \phi \cos \theta, \quad \frac{\partial y}{\partial \phi} = r \cos \phi \sin \theta,
\]
\[
\frac{\partial z}{\partial r} = \cos \phi, \quad \frac{\partial z}{\partial \theta} = 0, \quad \frac{\partial z}{\partial \phi} = -r \sin \phi.
\]
Entonces, el jacobiano queda:
\[
J = \begin{bmatrix}
\sin \phi \cos \theta & -r \sin \phi \sin \theta & r \cos \phi \cos \theta \\
\sin \phi \sin \theta & r \sin \phi \cos \theta & r \cos \phi \sin \theta \\
\cos \phi & 0 & -r \sin \phi
\end{bmatrix}.
\]

\paragraph{Determinante del Jacobiano}
El determinante de \( J \) se calcula como:
\[
\det(J) = - r^2 \sin \phi.
\]
\emph{Ojo}: Este determinante es clave para integrar en coordenadas esféricas, porque es el factor que "deforma" el volumen. Si \(\sin \phi = 0\) (o sea, \(\phi = 0\) o \(\pi\)), el jacobiano se anula (estamos en el eje \(z\)).

%\paragraph{}
%### **¿Por qué esto es útil?**
%- **Integrales**: Para cambiar de cartesianas a esféricas, multiplicás por \( |\det(J)| = r^2 \sin \phi \).
%- **Simetrías**: Ideal para problemas con simetría radial (ej: campos eléctricos, fluidos, etc.).

\paragraph{Nota:}
%**Nota rioplatense**: 
Si alguna vez te confundís con los ángulos, recordá que \(\theta\) es como la "longitud" (en el plano \(xy\)) y \(\phi\) es la "latitud" (desde el eje \(z\)). ¡Y listo! 

\subsection{Operadores diferenciales}

\noindent Veamos ahora a nivel introductivo operadores diferenciales y ecuaciones en derivadas parciales simples.

\paragraph{Divergencia}
La divergencia de un campo vectorial \( \mathbf{F}: \mathbb{R}^n \to \mathbb{R}^n \) se define como:
\[ \nabla \cdot \mathbf{F} = \sum_{i=1}^n \frac{\partial F_i}{\partial x_i}. \]

\paragraph{Laplaciano Definido por Divergencia}
El laplaciano de una función escalar \( f: \mathbb{R}^n \to \mathbb{R} \) se define como la divergencia del gradiente de \( f \):
\[ \Delta f = \nabla \cdot (\nabla f) = \sum_{i=1}^n \frac{\partial^2 f}{\partial x_i^2}. \]

\paragraph{Expresión del Laplaciano en Coordenadas Polares}

El laplaciano de una función escalar \( u(r, \theta) \) en 2D se escribe como:
\[
\nabla^2 u = \frac{1}{r} \frac{\partial}{\partial r}\left(r \frac{\partial u}{\partial r}\right) + \frac{1}{r^2} \frac{\partial^2 u}{\partial \theta^2}.
\]
%### **Explicación Detallada**
%
%1. **Coordenadas Polares**:
En coordenadas polares, un punto en el plano se representa por \((r, \theta)\), donde:
\begin{itemize}
    \item \( r \) es la distancia al origen,
    \item \( \theta \) es el ángulo con respecto al eje \( x \).
\end{itemize}

\noindent
Derivadas Parciales:
   - El término \(\frac{1}{r} \frac{\partial}{\partial r}\left(r \frac{\partial u}{\partial r}\right)\) representa la derivada radial, que incluye el factor \( r \) para tener en cuenta la variación del área con \( r \).
   - El término \(\frac{1}{r^2} \frac{\partial^2 u}{\partial \theta^2}\) representa la derivada angular, dividida por \( r^2 \) para normalizar la variación angular.

\medskip

\noindent Interpretación Física:
   El laplaciano mide la diferencia entre el valor de la función en un punto y el promedio de sus valores en un entorno cercano. En coordenadas polares, esta diferencia debe considerar la geometría del sistema (círculos concéntricos y ángulos).

\paragraph{Gradiente en Coordenadas Polares}
%
%## 1. **Gradiente en Coordenadas Polares**

El gradiente de una función escalar \( u(r, \theta) \) en coordenadas polares es:
%
\[
\nabla u = \frac{\partial u}{\partial r} \hat{\mathbf{r}} + \frac{1}{r} \frac{\partial u}{\partial \theta} \hat{\boldsymbol{\theta}}
\]
%
Derivación:
en coordenadas cartesianas, el gradiente es:
  \[
  \nabla u = \frac{\partial u}{\partial x} \hat{\mathbf{x}} + \frac{\partial u}{\partial y} \hat{\mathbf{y}}
  \]
La relación entre coordenadas polares y cartesianas es:
  \[
  x = r \cos \theta, \quad y = r \sin \theta
  \]
Aplicando la regla de la cadena:
  \[
  \frac{\partial u}{\partial x} = \frac{\partial u}{\partial r} \frac{\partial r}{\partial x} + \frac{\partial u}{\partial \theta} \frac{\partial \theta}{\partial x} = \frac{\partial u}{\partial r} \cos \theta - \frac{1}{r} \frac{\partial u}{\partial \theta} \sin \theta
  \]
  \[
  \frac{\partial u}{\partial y} = \frac{\partial u}{\partial r} \frac{\partial r}{\partial y} + \frac{\partial u}{\partial \theta} \frac{\partial \theta}{\partial y} = \frac{\partial u}{\partial r} \sin \theta + \frac{1}{r} \frac{\partial u}{\partial \theta} \cos \theta
  \]
Sustituyendo en el gradiente cartesiano y agrupando términos:
  \[
  \nabla u = \left( \frac{\partial u}{\partial r} \cos \theta - \frac{1}{r} \frac{\partial u}{\partial \theta} \sin \theta \right) \hat{\mathbf{x}} + \left( \frac{\partial u}{\partial r} \sin \theta + \frac{1}{r} \frac{\partial u}{\partial \theta} \cos \theta \right) \hat{\mathbf{y}}
  \]
Usando la transformación de los vectores unitarios:
  \[
  \hat{\mathbf{r}} = \cos \theta \hat{\mathbf{x}} + \sin \theta \hat{\mathbf{y}}, \quad \hat{\boldsymbol{\theta}} = -\sin \theta \hat{\mathbf{x}} + \cos \theta \hat{\mathbf{y}}
  \]
Se obtiene finalmente:
  \[
  \nabla u = \frac{\partial u}{\partial r} \hat{\mathbf{r}} + \frac{1}{r} \frac{\partial u}{\partial \theta} \hat{\boldsymbol{\theta}}
  \]

\paragraph{Divergencia en Coordenadas Polares}
%## 2. **Divergencia en Coordenadas Polares**

La divergencia de un campo vectorial \( \mathbf{F} = F_r \hat{\mathbf{r}} + F_\theta \hat{\boldsymbol{\theta}} \) en coordenadas polares es:
\[
\nabla \cdot \mathbf{F} = \frac{1}{r} \frac{\partial}{\partial r} (r F_r) + \frac{1}{r} \frac{\partial F_\theta}{\partial \theta}
\]
Derivación:
en coordenadas cartesianas, la divergencia es:
  \[
  \nabla \cdot \mathbf{F} = \frac{\partial F_x}{\partial x} + \frac{\partial F_y}{\partial y}
  \]
Expresando \( F_x \) y \( F_y \) en términos de \( F_r \) y \( F_\theta \):
  \[
  F_x = F_r \cos \theta - F_\theta \sin \theta, \quad F_y = F_r \sin \theta + F_\theta \cos \theta
  \]
Aplicando la regla de la cadena y simplificando, se llega a:
  \[
  \nabla \cdot \mathbf{F} = \frac{1}{r} \frac{\partial}{\partial r} (r F_r) + \frac{1}{r} \frac{\partial F_\theta}{\partial \theta}
  \]

%## 3. **Laplaciano como Divergencia del Gradiente**
\subsection{Laplaciano, divergencia, gradiente}

\noindent El laplaciano de \( u \) es la divergencia del gradiente de \( u \):
\[
\nabla^2 u = \nabla \cdot (\nabla u)
\]
Sustituyendo el gradiente en coordenadas polares:
\[
\nabla u = \frac{\partial u}{\partial r} \hat{\mathbf{r}} + \frac{1}{r} \frac{\partial u}{\partial \theta} \hat{\boldsymbol{\theta}}
\]
Aplicando la divergencia al gradiente:
\[
\nabla \cdot (\nabla u) = \frac{1}{r} \frac{\partial}{\partial r} \left( r \frac{\partial u}{\partial r} \right) + \frac{1}{r} \frac{\partial}{\partial \theta} \left( \frac{1}{r} \frac{\partial u}{\partial \theta} \right)
\]
Simplificando:
\[
\nabla^2 u = \frac{1}{r} \frac{\partial}{\partial r} \left( r \frac{\partial u}{\partial r} \right) + \frac{1}{r^2} \frac{\partial^2 u}{\partial \theta^2}
\]

\subsection{Soluci\'on fundamental del laplaciano 2D}

\noindent Recordemos que la solución fundamental \( \Phi(\mathbf{x}) \) satisface la ecuación:
\[
\nabla^2 \Phi = \delta(\mathbf{x}),
\]
donde \( \delta(\mathbf{x}) \) es la delta de Dirac en el origen. En coordenadas polares \((r, \theta)\), la solución fundamental del laplaciano en 2D es:
\[
\Phi(r) = \frac{1}{2\pi} \ln r.
\]

\paragraph{Derivación}

%1. Ecuación del Laplaciano en Coordenadas Polares
El laplaciano de una función \( \Phi(r, \theta) \) en coordenadas polares es:
\[
\nabla^2 \Phi = \frac{1}{r} \frac{\partial}{\partial r}\left(r \frac{\partial \Phi}{\partial r}\right) + \frac{1}{r^2} \frac{\partial^2 \Phi}{\partial \theta^2}
\]
Como la solución fundamental solo depende de \( r \) (por simetría radial), se simplifica a:
\[
\nabla^2 \Phi = \frac{1}{r} \frac{d}{dr}\left(r \frac{d\Phi}{dr}\right)
\]
%2. Ecuación a Resolver
Queremos resolver:
\[
\frac{1}{r} \frac{d}{dr}\left(r \frac{d\Phi}{dr}\right) = \delta(\mathbf{x})
\]
Fuera del origen (\( r \neq 0 \)), la ecuación se reduce a:
\[
\frac{1}{r} \frac{d}{dr}\left(r \frac{d\Phi}{dr}\right) = 0
\]
%
%3. Resolviendo la Ecuación Diferencial
Multiplicamos ambos lados por \( r \):
\[
\frac{d}{dr}\left(r \frac{d\Phi}{dr}\right) = 0
\]
Integrando una vez:
\[
r \frac{d\Phi}{dr} = C_1
\]
Donde \( C_1 \) es una constante. Integrando nuevamente:
\[
\Phi(r) = C_1 \ln r + C_2
\]
%
%4. Determinando las Constantes
%Condición de regularidad en el infinito**: \( \Phi(r) \to 0 \) cuando \( r \to \infty \). Esto implica que \( C_1 = 0 \) no es válido, ya que \( \ln r \to \infty \) cuando \( r \to \infty \). Sin embargo, 
En el contexto de la solución fundamental, nos interesa el comportamiento cerca del origen.
%- Condición de singularidad en el origen: 
La solución fundamental debe tener una singularidad en \( r = 0 \) que reproduzca el efecto de la delta de Dirac. La solución clásica es:
\[
\Phi(r) = \frac{1}{2\pi} \ln r
\]
La constante \( \frac{1}{2\pi} \) se determina integrando el laplaciano de \( \Phi \) sobre un círculo de radio \( \epsilon \) centrado en el origen y aplicando el teorema de la divergencia. El resultado debe ser 1 (la integral de la delta de Dirac).
Verificamos.
Aplicando el laplaciano a \( \Phi(r) = \frac{1}{2\pi} \ln r \):
\[
\frac{d\Phi}{dr} = \frac{1}{2\pi r}
\]
\[
\frac{d}{dr}\left(r \frac{d\Phi}{dr}\right) = \frac{d}{dr}\left(r \cdot \frac{1}{2\pi r}\right) = \frac{d}{dr}\left(\frac{1}{2\pi}\right) = 0 \quad \text{para } r \neq 0
\]
En \( r = 0 \), la función \( \ln r \) es singular. % y su laplaciano es proporcional a la delta de Dirac.

\subsection{Ejemplo de Campo con Divergencia Nula}

\noindent
Considerá el campo vectorial en \(\mathbb{R}^3\):
\[
\mathbf{F}(x, y, z) = (-y, x, 0)
\]
La divergencia de \(\mathbf{F}\) es:
\[
\nabla \cdot \mathbf{F} = \frac{\partial}{\partial x}(-y) + \frac{\partial}{\partial y}(x) + \frac{\partial}{\partial z}(0) = 0.
\]
Este campo representa un movimiento de rotación alrededor del eje \(z\) y no tiene fuentes ni sumideros (divergencia nula).

\subsection{Ejemplo de Campo con Divergencia No Nula}

\noindent Considerá el campo vectorial en \(\mathbb{R}^3\):
\[
\mathbf{F}(x, y, z) = (x, y, z)
\]
La divergencia de \(\mathbf{F}\) es:
\[
\nabla \cdot \mathbf{F} = \frac{\partial}{\partial x}(x) + \frac{\partial}{\partial y}(y) + \frac{\partial}{\partial z}(z) = 3.
\]
Este campo tiene divergencia no nula, lo que indica que hay un "flujo" que se expande (o contrae) en el espacio.

\subsection{Potencial y Rotacional}

\paragraph{Campos que Derivan de un Potencial \(C^2\) y su Rotacional (2D)}

\noindent Supongamos que tenés un campo vectorial \(\mathbf{F}(x, y) = (P, Q)\) que deriva de un potencial \(\phi(x, y)\) de clase \(C^2\). Esto significa que:
\[
\mathbf{F} = \nabla \phi = \left( \frac{\partial \phi}{\partial x}, \frac{\partial \phi}{\partial y} \right).
\]
En este caso, el rotacional (o curl) de \(\mathbf{F}\) es:
\[
\nabla \times \mathbf{F} = \frac{\partial Q}{\partial x} - \frac{\partial P}{\partial y} = \frac{\partial^2 \phi}{\partial x \partial y} - \frac{\partial^2 \phi}{\partial y \partial x} = 0,
\]
por el teorema de Schwarz/Clairaut (igualdad de las derivadas cruzadas).

\noindent Ejemplo.
Considerá el campo:
\[
\mathbf{F}(x, y) = (2x, 2y).
\]
Este campo deriva del potencial \(\phi(x, y) = x^2 + y^2\). El rotacional de \(\mathbf{F}\) es:
\[
\nabla \times \mathbf{F} = \frac{\partial}{\partial x}(2y) - \frac{\partial}{\partial y}(2x) = 0.
\]

\paragraph{En 3D}
Supongamos que tenés un campo vectorial \(\mathbf{F}(x, y, z) = (P, Q, R)\) que deriva de un potencial \(\phi(x, y, z)\) de clase \(C^2\). Esto significa que:
\[
\mathbf{F} = \nabla \phi = \left( \frac{\partial \phi}{\partial x}, \frac{\partial \phi}{\partial y}, \frac{\partial \phi}{\partial z} \right).
\]
En este caso, el rotacional de \(\mathbf{F}\) es:
\[
\nabla \times \mathbf{F} = \left( \frac{\partial R}{\partial y} - \frac{\partial Q}{\partial z}, \frac{\partial P}{\partial z} - \frac{\partial R}{\partial x}, \frac{\partial Q}{\partial x} - \frac{\partial P}{\partial y} \right) = (0, 0, 0),
\]
de nuevo, por la igualdad de las derivadas cruzadas.

\noindent Ejemplo.
Considerá el campo:
\[
\mathbf{F}(x, y, z) = (2x, 2y, 2z).
\]
Este campo deriva del potencial \(\phi(x, y, z) = x^2 + y^2 + z^2\). El rotacional de \(\mathbf{F}\) es:
\[
\nabla \times \mathbf{F} = \left( \frac{\partial}{\partial y}(2z) - \frac{\partial}{\partial z}(2y), \frac{\partial}{\partial z}(2x) - \frac{\partial}{\partial x}(2z), \frac{\partial}{\partial x}(2y) - \frac{\partial}{\partial y}(2x) \right) = (0, 0, 0).
\]

\paragraph{Conclusión}
Si un campo vectorial deriva de un potencial \(C^2\), entonces su rotacional es cero. Esto es válido tanto en 2D como en 3D.


\subsection{EDPs Básicas}

\noindent
Las ecuaciones diferenciales parciales (EDPs) son ecuaciones que involucran derivadas parciales de una función. Algunas EDPs básicas incluyen:
\begin{enumerate}
    \item Ecuación de Laplace: \( \Delta f = 0 \)
    \item Ecuación de Poisson: \( \Delta f = g \)
    \item Ecuación de Onda: \( \frac{\partial^2 f}{\partial t^2} = c^2 \Delta f \)
    \item Ecuación de Calor: \( \frac{\partial f}{\partial t} = k \Delta f \)
\end{enumerate}

\paragraph{Ecuación de Laplace}
La ecuación de Laplace es:
\[
\nabla^2 u = 0.
\]
Solución en 1D.
En una dimensión, la ecuación se reduce a:
\[
\frac{d^2 u}{dx^2} = 0.
\]
La solución general es:
\[
u(x) = Ax + B,
\]
donde \(A\) y \(B\) son constantes determinadas por las condiciones de borde.
Solución en 2D.
En dos dimensiones, la ecuación es:
\[
\frac{\partial^2 u}{\partial x^2} + \frac{\partial^2 u}{\partial y^2} = 0.
\]
Una solución básica (armónica) es:
\[
u(x, y) = x^2 - y^2.
\]

\paragraph{Ecuación de Poisson}
La ecuación de Poisson es:
\[
\nabla^2 u = f(x, y).
\]
Solución en 1D.
En una dimensión, la ecuación se reduce a:
\[
\frac{d^2 u}{dx^2} = f(x).
\]
La solución general es:
\[
u(x) = \int \int f(x) \, dx^2 + Ax + B,
\]
donde \(A\) y \(B\) son constantes determinadas por las condiciones de borde.

\medskip

\noindent
Solución en 2D.
En dos dimensiones, la ecuación es:
\[
\frac{\partial^2 u}{\partial x^2} + \frac{\partial^2 u}{\partial y^2} = f(x, y).
\]
Una solución básica para \(f(x, y) = 0\) es la misma que para Laplace. Para \(f(x, y) = x^2 + y^2\), una solución particular es:
\[
u(x, y) = \frac{x^4}{12} + \frac{y^4}{12}.
\]

\paragraph{Ecuación del Calor}
La ecuación del calor es:
\[
\frac{\partial u}{\partial t} = \alpha \nabla^2 u.
\]
Solución en 1D.
En una dimensión, la ecuación se reduce a:
\[
\frac{\partial u}{\partial t} = \alpha \frac{\partial^2 u}{\partial x^2}.
\]
Una solución básica es:
\[
u(x, t) = e^{-\alpha k^2 t} \sin(kx),
\]
donde \(k\) es una constante.

\medskip

\noindent
Solución en 2D.
En dos dimensiones, la ecuación es:
\[
\frac{\partial u}{\partial t} = \alpha \left( \frac{\partial^2 u}{\partial x^2} + \frac{\partial^2 u}{\partial y^2} \right).
\]
Una solución básica es:
\[
u(x, y, t) = e^{-\alpha (k_x^2 + k_y^2) t} \sin(k_x x) \sin(k_y y),
\]
donde \(k_x\) y \(k_y\) son constantes.

\paragraph{Ecuación de Onda}
La ecuación de onda es:
\[
\frac{\partial^2 u}{\partial t^2} = c^2 \nabla^2 u.
\]
Solución en 1D.
En una dimensión, la ecuación se reduce a:
\[
\frac{\partial^2 u}{\partial t^2} = c^2 \frac{\partial^2 u}{\partial x^2}.
\]
La solución general es:
\[
u(x, t) = f(x - ct) + g(x + ct),
\]
donde \(f\) y \(g\) son funciones arbitrarias determinadas por las condiciones iniciales.

\medskip

\noindent Solución en 2D.
En dos dimensiones, la ecuación es:
\[
\frac{\partial^2 u}{\partial t^2} = c^2 \left( \frac{\partial^2 u}{\partial x^2} + \frac{\partial^2 u}{\partial y^2} \right).
\]
Una solución básica es:
\[
u(x, y, t) = \sin(k_x x + k_y y - \omega t),
\]
donde \(k_x\), \(k_y\) y \(\omega\) son constantes que satisfacen la relación de dispersión \(\omega^2 = c^2 (k_x^2 + k_y^2)\).

\newpage

\section{Practico 6}

\paragraph{Funciones vectoriales}

\begin{enumerate} %[start=10]
\item Calcular las derivadas parciales de $f\circ g(p)$, en los ejemplos siguientes.

  \begin{enumerate}
  \item  $f(x,y)=e^{xy}\cos(x)$, $g(x,y,z)=(x^3,yz^2)$, $p=(1,1,-1)$.
  \item $f(x,y)=x^2-y^2$, $g(r,\theta)=(r\cos\theta, r\sen\theta)$, $p=(\sqrt{3},\pi)$.
  \item $f(x,y,z)=z+\ln(x+y)$, $g(x,y)=(x+y, xy, y^2), p=(0,-1)$.
    \end{enumerate}
\item Consideremos los cambios de coordenadas a cil\'indricas y esf\'ericas :
  \[f(\rho,\theta,z)=(\rho\cos\theta, \rho\sen\theta,z), \ g(\rho,\theta,\varphi)=(\rho\sen\varphi\cos\theta, \rho\sen\varphi\sen\theta, \rho\cos\varphi),\]
  donde $\rho>0, 0<\theta<2\pi, z\in\R$ y $0<\varphi<\pi$.
  \begin{itemize}
  \item[(a)] Calcular las matrices Jacobianas de $f$ y $g$ en un punto general del dominio.
  \item[(b)] Si $h(x,y,z)=(x^2+y^2, z^2xy)$ verificar la regla de la cadena para la jacobiana de $h\circ f$ en el punto $(1,\pi/2,0)$.
    \item[(c)] Si $h(x,y,z)=x^2+y^2+z^2$, verificar la regla de la cadena para el gradiente de $h\circ g$ en el punto $(1,\pi,\pi/2)$. 
    \end{itemize}

\item Exprese una elipse de la forma $ax^2+by^2=1$ usando las coordenadas polares, y lo mismo para un elipsoide $ax^2+by^2+cz^2=1$ usando las coordenadas esféricas.
\item Verificar en qué puntos del dominio las funciones siguientes definen difeomorfismos locales:
  \begin{itemize}
  %\item[(a)] $f:\R^2\to\R^2$, $f(x,y)=(xy,y+x^2)$.
  \item[(a)] $f:\R^2-\{(0,0)\}\to\R^2$, $f(x,y)=(\frac{-y}{x^2+y^2},\frac{x}{x^2+y^2})$.
   
  \item[(b)] $f:\R^3\to\R^3$, $f(x,y,z)=(x^2,yz,z^2)$.
    \item[(c)] $f:\R^3\to\R^3$, $f(x,y,z)=(xz,y,x^2z^2-y^2)$.
    \end{itemize}
\end{enumerate}



\chapter{Integrales múltiples}

\section{Temario}

\begin{itemize}
    \item Definición y aplicaciones. 
    \item Cálculo por métodos de integrales iteradas y cambio de variable.
\end{itemize}

\section{Notas}

\subsection{Definición y Aplicaciones}

Las integrales múltiples son una extensión natural de las integrales simples a funciones de varias variables. Se utilizan para calcular volúmenes, masas, centros de masa, y muchas otras cantidades en espacios de dos, tres o más dimensiones.

\subsubsection{Definición}

\noindent Una integral doble sobre una región \( D \) en el plano \( xy \) se define como:
\[
\iint_D f(x, y) \, dA = \lim_{n \to \infty} \sum_{i=1}^{n} f(x_i, y_i) \Delta A_i
\]
donde \( \Delta A_i \) es el área de un pequeño subrectángulo de \( D \), y \( (x_i, y_i) \) es un punto en ese subrectángulo. Similarmente, una integral triple sobre una región \( E \) en el espacio \( xyz \) se define como:
\[
\iiint_E f(x, y, z) \, dV = \lim_{n \to \infty} \sum_{i=1}^{n} f(x_i, y_i, z_i) \Delta V_i
\]

%\begin{proposition}
Sea \( f: D \subseteq \mathbb{R}^2 \to \mathbb{R} \) una función acotada definida en un dominio \( D \) (rectángulo o conjunto medible).
Si el conjunto de puntos de discontinuidad de \( f \) tiene área nula (medida de Lebesgue cero),
entonces \( f \) es integrable según Riemann en \( D \).

En particular, si \( f \) es continua en \( D \), el conjunto de discontinuidades es vacío (y por tanto de área nula),
lo que implica que toda función continua en un dominio compacto es integrable.
%\end{position}

\subsubsection{Una propiedad}

%\begin{proposition}
Sea \( A \subseteq \mathbb{R}^2 \) un dominio acotado y \( R \) un rectángulo tal que \( A \subseteq R \).
Dada una función \( f: A \to \mathbb{R} \) acotada, se define su \textbf{extensión por cero} fuera de \( A \) como:
\[
f_R(x,y) =
\begin{cases}
f(x,y) & \text{si } (x,y) \in A, \\
0 & \text{si } (x,y) \in R \setminus A.
\end{cases}
\]
Si \( f \) es integrable en \( A \), entonces su extensión \( f_R \) es integrable en \( R \) y se cumple:
\[
\iint_A f(x,y) \, dA = \iint_R f_R(x,y) \, dA.
\]
%\end{proposition}


\subsubsection{Aplicaciones}

\noindent Las integrales múltiples tienen numerosas aplicaciones en física e ingeniería, incluyendo:

\begin{itemize}
    \item Cálculo de áreas y volúmenes.
    \item Cálculo de masas y momentos de inercia.
    \item Determinación de centros de masa y gravedad.
    \item Cálculo de probabilidades en estadística.
\end{itemize}

\subsection{Cálculo por Métodos de Integrales Iteradas y Cambio de Variable}

\subsubsection{Integrales Iteradas}

El cálculo de integrales múltiples se puede simplificar utilizando integrales iteradas. Para una integral doble, esto significa integrar primero con respecto a una variable y luego con respecto a la otra:
\[
\iint_D f(x, y) \, dA = \int_{a}^{b} \left( \int_{g_1(x)}^{g_2(x)} f(x, y) \, dy \right) dx
\]
donde \( D \) es la región limitada por \( a \leq x \leq b \) y \( g_1(x) \leq y \leq g_2(x) \).

\subsubsection{Fubini}

\begin{theorem}[Fubini]
Sea \( f: [a,b] \times [c,d] \to \mathbb{R} \) una función integrable en el rectángulo \( R = [a,b] \times [c,d] \).
Entonces, las integrales iteradas existen y se cumple:
\[
\int_{R} f(x,y) \, dA = \int_{a}^{b} \left( \int_{c}^{d} f(x,y) \, dy \right) dx = \int_{c}^{d} \left( \int_{a}^{b} f(x,y) \, dx \right) dy.
\]
\end{theorem}
\noindent Ejemplo:
\[
\begin{aligned}
\int_{-1}^{1} \int_{0}^{1} \sen(x^2) y \, dx \, dy
&= \int_{-1}^{1} y \left( \int_{0}^{1} \sen(x^2) \, dx \right) dy \quad \text{(por Fubini)} \\
&= \left( \int_{-1}^{1} y \, dy \right) \left( \int_{0}^{1} \sen(x^2) \, dx \right) \quad \text{(por linealidad)} \\
&= \left[ \frac{y^2}{2} \right]_{-1}^{1} \cdot \int_{0}^{1} \sen(x^2) \, dx \\
&= \left( \frac{1}{2} - \frac{1}{2} \right) \cdot \int_{0}^{1} \sen(x^2) \, dx \\
&= 0.
\end{aligned}
\]

\subsubsection{Cambio de Variable}

El cambio de variable es una técnica poderosa para simplificar el cálculo de integrales múltiples. Utilizando un cambio de coordenadas, podemos transformar una integral compleja en una más simple. La fórmula general para el cambio de variable en una integral doble es:
\[
\iint_D f(x, y) \, dx \, dy = \iint_{D'} f(u, v) \left| \frac{\partial(x, y)}{\partial(u, v)} \right| \, du \, dv
\]
donde \( \left| \frac{\partial(x, y)}{\partial(u, v)} \right| \) es el determinante Jacobiano de la transformación.

\subsubsection{Ejemplo: Coordenadas Polares}

Un ejemplo común de cambio de variable es la transformación a coordenadas polares, donde \( x = r \cos \theta \) y \( y = r \sin \theta \). El determinante Jacobiano para esta transformación es \( r \), por lo que la integral doble en coordenadas polares se escribe como:
\[
\iint_D f(x, y) \, dx \, dy = \iint_{D'} f(r \cos \theta, r \sin \theta) \, r \, dr \, d\theta
\]
Este método es especialmente útil para regiones circulares o anulares. La idea clave:
\[
dA_C = dxdy; \quad dA_P = rdrd\theta . 
\]

%\section{Notas}

\subsection{Ejemplos} 

\begin{enumerate}

\item \textbf{Ejemplo 1: Integral doble sobre un rectángulo}

Calcular la integral doble:
\[
\iint_D (2x + 3y) \, dA
\]
donde \( D \) es el rectángulo definido por \( 0 \leq x \leq 1 \) y \( 0 \leq y \leq 2 \).

\textbf{Solución:}
\[
\int_{0}^{1} \int_{0}^{2} (2x + 3y) \, dy \, dx = \int_{0}^{1} \left[ 2xy + \frac{3}{2}y^2 \right]_{0}^{2} \, dx = \int_{0}^{1} (4x + 6) \, dx = \left[ 2x^2 + 6x \right]_{0}^{1} = 8.
\]

\item \textbf{Ejemplo 2: Integral doble sobre una región triangular}

Calcular la integral doble:
\[
\iint_D xy \, dA
\]
donde \( D \) es la región triangular con vértices en \( (0,0) \), \( (1,0) \) y \( (1,1) \).

\textbf{Solución:}
\[
\int_{0}^{1} \int_{0}^{x} xy \, dy \, dx = \int_{0}^{1} \left[ \frac{x y^2}{2} \right]_{0}^{x} \, dx = \int_{0}^{1} \frac{x^3}{2} \, dx = \left[ \frac{x^4}{8} \right]_{0}^{1} = \frac{1}{8}.
\]

\item \textbf{Ejemplo 3. Integral doble en coordenadas polares}

Calcular la integral doble:
\[
\iint_D e^{-x^2-y^2} \, dA
\]
donde \( D \) es el disco unitario.

\textbf{Solución:}
\[
\int_{0}^{2\pi} \int_{0}^{1} e^{-r^2} r \, dr \, d\theta = \int_{0}^{2\pi} \left[ -\frac{1}{2} e^{-r^2} \right]_{0}^{1} \, d\theta = \int_{0}^{2\pi} \frac{1}{2} (1 - e^{-1}) \, d\theta = \pi (1 - e^{-1}).
\]

\item \textbf{Ejemplo 4: Integral triple sobre un paralelepípedo}

Calcular la integral triple:
\[
\iiint_E xyz \, dV
\]
donde \( E \) es el paralelepípedo definido por \( 0 \leq x \leq 1 \), \( 0 \leq y \leq 2 \) y \( 0 \leq z \leq 3 \).

\textbf{Solución:}
\[
\int_{0}^{1} \int_{0}^{2} \int_{0}^{3} xyz \, dz \, dy \, dx = \int_{0}^{1} \int_{0}^{2} \left[ \frac{xyz^2}{2} \right]_{0}^{3} \, dy \, dx = \int_{0}^{1} \int_{0}^{2} \frac{9xy}{2} \, dy \, dx %= \int_{0}^{1} \left[ \frac{9xy^2}{4} \right]_{0}^{2} \, dx = \int_{0}^{1} 9x \, dx = \frac{9}{2}
\]
y
\[
%\int_{0}^{1} \int_{0}^{2} \int_{0}^{3} xyz \, dz \, dy \, dx = \int_{0}^{1} \int_{0}^{2} \left[ \frac{xyz^2}{2} \right]_{0}^{3} \, dy \, dx = 
\int_{0}^{1} \int_{0}^{2} \frac{9xy}{2} \, dy \, dx = \int_{0}^{1} \left[ \frac{9xy^2}{4} \right]_{0}^{2} \, dx = \int_{0}^{1} 9x \, dx = \frac{9}{2}.
\]

\item \textbf{Ejemplo 5: Integral triple en coordenadas cilíndricas}

Calcular la integral triple:
\[
\iiint_E z \, dV
\]
donde \( E \) es el cilindro definido por \( x^2 + y^2 \leq 1 \) y \( 0 \leq z \leq 1 \).

\textbf{Solución:}
\[
\int_{0}^{2\pi} \int_{0}^{1} \int_{0}^{1} z r \, dz \, dr \, d\theta = \int_{0}^{2\pi} \int_{0}^{1} \left[ \frac{z^2 r}{2} \right]_{0}^{1} \, dr \, d\theta = \int_{0}^{2\pi} \int_{0}^{1} \frac{r}{2} \, dr \, d\theta %= \int_{0}^{2\pi} \left[ \frac{r^2}{4} \right]_{0}^{1} \, d\theta = \int_{0}^{2\pi} \frac{1}{4} \, d\theta = \frac{\pi}{2}
\]
y
\[
%\int_{0}^{2\pi} \int_{0}^{1} \int_{0}^{1} z r \, dz \, dr \, d\theta = \int_{0}^{2\pi} \int_{0}^{1} \left[ \frac{z^2 r}{2} \right]_{0}^{1} \, dr \, d\theta = \int_{0}^{2\pi} \int_{0}^{1} \frac{r}{2} \, dr \, d\theta = 
\int_{0}^{2\pi} \left[ \frac{r^2}{4} \right]_{0}^{1} \, d\theta = \int_{0}^{2\pi} \frac{1}{4} \, d\theta = \frac{\pi}{2}.
\]
\item \textbf{Ejemplo 6: Integral doble con cambio de orden de integración}

Calcular la integral doble cambiando el orden de integración:
\[
\int_{0}^{1} \int_{x}^{1} \sin(y) \, dy \, dx
\]

\textbf{Solución:}
\[
\int_{0}^{1} \int_{0}^{y} \sin(y) \, dx \, dy = \int_{0}^{1} \left[ x \sin(y) \right]_{0}^{y} \, dy = \int_{0}^{1} y \sin(y) \, dy
\]
Usando integración por partes:
\[
\int_{0}^{1} y \sin(y) \, dy = \left[ -y \cos(y) \right]_{0}^{1} + \int_{0}^{1} \cos(y) \, dy = -\cos(1) + \sin(1).
\]

\item \textbf{Ejemplo 7: Integral doble sobre una región acotada por curvas}

Calcular la integral doble:
\[
\iint_D (x + y) \, dA
\]
donde \( D \) es la región acotada por \( y = x^2 \) y \( y = \sqrt{x} \).

\textbf{Solución:}
\[
\int_{0}^{1} \int_{x^2}^{\sqrt{x}} (x + y) \, dy \, dx = \int_{0}^{1} \left[ xy + \frac{y^2}{2} \right]_{x^2}^{\sqrt{x}} \, dx = \int_{0}^{1} \left( x^{3/2} + \frac{x}{2} - x^3 - \frac{x^4}{2} \right) \, dx %= \left[ \frac{2}{5}x^{5/2} + \frac{x^2}{4} - \frac{x^4}{4} - \frac{x^5}{10} \right]_{0}^{1} = \frac{8}{20} + \frac{1}{4} - \frac{1}{4} - \frac{1}{10} = \frac{3}{10}
\]
y
\[
%\int_{0}^{1} \int_{x^2}^{\sqrt{x}} (x + y) \, dy \, dx = \int_{0}^{1} \left[ xy + \frac{y^2}{2} \right]_{x^2}^{\sqrt{x}} \, dx %= 
\int_{0}^{1} \left( x^{3/2} + \frac{x}{2} - x^3 - \frac{x^4}{2} \right) \, dx = 
\left[ \frac{2}{5}x^{5/2} + \frac{x^2}{4} - \frac{x^4}{4} - \frac{x^5}{10} \right]_{0}^{1} = \frac{8}{20} + \frac{1}{4} - \frac{1}{4} - \frac{1}{10} = \frac{3}{10}.
\]

\item \textbf{Ejemplo 8: Integral triple sobre una esfera}

Calcular la integral triple:
\[
\iiint_E z^2 \, dV
\]
donde \( E \) es la esfera unitaria.

\textbf{Solución:}
Usando coordenadas esféricas:
\[
\int_{0}^{2\pi} \int_{0}^{\pi} \int_{0}^{1} (r \cos \phi)^2 r^2 \sin \phi \, dr \, d\phi \, d\theta = \int_{0}^{2\pi} \int_{0}^{\pi} \int_{0}^{1} r^4 \cos^2 \phi \sin \phi \, dr \, d\phi \, d\theta 
\]
\[
= \int_{0}^{2\pi} \int_{0}^{\pi} \left[ \frac{r^5}{5} \right]_{0}^{1} \cos^2 \phi \sin \phi \, d\phi \, d\theta 
= \frac{1}{5} \int_{0}^{2\pi} \int_{0}^{\pi} \cos^2 \phi \sin \phi \, d\phi \, d\theta
\]
\[
= \frac{1}{5} \int_{0}^{2\pi} \left[ -\frac{\cos^3 \phi}{3} \right]_{0}^{\pi} \, d\theta = \frac{2}{15} \int_{0}^{2\pi} d\theta = \frac{4\pi}{15}.
\]

\item \textbf{Ejemplo 9: Integral doble con cambio de variable}

Calcular la integral doble usando el cambio de variable \( u = x - y \) y \( v = x + y \):
\[
I = \iint_D f(x+y,x-y) \, dA
\]
donde \( D \) es el diamante con vértices en \( (-1,0) \), \( (0,1) \), \( (1,0) \) y \( (0,-1) \) (rotacion de un cuadrado).

\textbf{Solución:}
Tenemos el Jacobiano siguiente
\[
\frac{\partial(u, v)}{\partial(x, y)} = \begin{vmatrix}
1 & -1 \\
1 & 1
\end{vmatrix} = 2.
\]
Entonces:
%de la transformación es:
\[
\frac{\partial(x, y)}{\partial(u, v)} = \begin{vmatrix}
%\frac{1}{2} & \frac{1}{2} \\
%-\frac{1}{2} & \frac{1}{2}
1 & -1 \\
1 & 1
\end{vmatrix}^{-1} = \frac{1}{2}.
\]
%Tenemos 
%\[
%0 \leq x \leq 1, \quad, 
%0 \leq y \leq 1,
%\]
%entonces
%\[
%0 \leq x+y \leq 2, \quad, 
%-1 \leq x-y \leq 1.
%\]
El dominio se parametriza con
\[
-1 \leq u \leq 1, \quad 
-1 \leq v \leq 1. 
\]
La integral se transforma en:
\[
I = \frac{1}{2}\int_{u=-1}^{1} \int_{v=-1}^{1} f(u,v) \, du \, dv.
\]
En el caso $f=1$, obtenemos
\[
I = 2 (=(\sqrt{2})^2).
\]
En el caso
\[
f(x-y,x+y)=(x-y)(x+y)=x^2 - y^2
\]
obtenemos
\[
I = \frac{1}{2}\int_{u=-1}^{1} \int_{v=-1}^{1} uv \, du \, dv = 0.
\]
%Calcular la integral doble usando el cambio de variable \( u = x - y \) y \( v = x + y \):
%\[
%\iint_D (x - y) \sin(x + y) \, dA
%\]
%donde \( D \) es el cuadrado con vértices en \( (0,0) \), \( (1,0) \), \( (1,1) \) y \%( (0,1) \).

%\textbf{Solución:}
%Tenemos el Jacobiano siguiente
%\[
%\frac{\partial(u, v)}{\partial(x, y)} = %\begin{vmatrix}
%1 & -1 \\
%1 & 1
%\end{vmatrix} = 2.
%\]
%Entonces:
%de la transformación es:
%\[
%\frac{\partial(x, y)}{\partial(u, v)} = %\begin{vmatrix}
%%\frac{1}{2} & \frac{1}{2} \\
%%-\frac{1}{2} & \frac{1}{2}
%1 & -1 \\
%1 & 1
%\end{vmatrix}^{-1} = \frac{1}{2}.
%\]
%Tenemos 
%\[
%0 \leq x \leq 1, \quad, 
%0 \leq y \leq 1,
%\]
%entonces
%\[
%0 \leq x+y \leq 2, \quad, 
%-1 \leq x-y \leq 1.
%\]
%La integral se transforma en:
%\[
%\int_{0}^{1} \int_{-1}^{1} u \sin(v) \, du %\, dv = \int_{0}^{1} \left[ \frac{u^2}{2} %\sin(v) \right]_{-1}^{1} \, dv = %\int_{0}^{1} 0 \, dv = 0.
%\]

\item \textbf{Ejemplo 10: Integral triple sobre un tetraedro}

Calcular la integral triple:
\[
\iiint_E z \, dV
\]
donde \( E \) es el tetraedro con vértices en \( (0,0,0) \), \( (1,0,0) \), \( (0,1,0) \) y \( (0,0,1) \).

\textbf{Solución:}
\[
\int_{0}^{1} \int_{0}^{1-x} \int_{0}^{1-x-y} z \, dz \, dy \, dx = \int_{0}^{1} \int_{0}^{1-x} \left[ \frac{z^2}{2} \right]_{0}^{1-x-y} \, dy \, dx 
\]
\[
= \int_{0}^{1} \int_{0}^{1-x} \frac{(1-x-y)^2}{2} \, dy \, dx
\]
\[
= \int_{0}^{1} \left[ -\frac{(1-x-y)^3}{6} \right]_{0}^{1-x} \, dx = \int_{0}^{1} \frac{(1-x)^3}{6} \, dx = \left[ -\frac{(1-x)^4}{24} \right]_{0}^{1} = \frac{1}{24}.
\]

\item \textbf{Ejemplo 11: Integral doble sobre un triangulo}

\[
\begin{aligned}
\iint_A 2xy \, dA &= \int_{x=0}^{1} \left( \int_{y=0}^{2-2x} 2xy \, dy \right) dx \\
&= \int_{0}^{1} 2x \left( \int_{0}^{2-2x} y \, dy \right) dx \\
&= \int_{0}^{1} 2x \left[ \frac{y^2}{2} \right]_{0}^{2-2x} dx \\
&= \int_{0}^{1} 2x \cdot \frac{(2-2x)^2}{2} \, dx \\
&= \int_{0}^{1} x (4 - 8x + 4x^2) \, dx \\
&= \int_{0}^{1} (4x - 8x^2 + 4x^3) \, dx \\
&= \left[ 2x^2 - \frac{8}{3}x^3 + x^4 \right]_{0}^{1} \\
&= 2 - \frac{8}{3} + 1 = \frac{6}{3} - \frac{8}{3} + \frac{3}{3} = \frac{1}{3} \cdot 3 = \boxed{\frac{1}{2}}.
\end{aligned}
\]

\[
\begin{aligned}
\iint_A 2xy \, dA &= \int_{y=0}^{2} \left( \int_{x=0}^{1 - \frac{y}{2}} 2xy \, dx \right) dy \\
&= \int_{0}^{2} y \left( \int_{0}^{1 - \frac{y}{2}} 2x \, dx \right) dy \\
&= \int_{0}^{2} y \left[ x^2 \right]_{0}^{1 - \frac{y}{2}} dy \\
&= \int_{0}^{2} y \left(1 - \frac{y}{2}\right)^2 dy \\
&= \int_{0}^{2} y \left(1 - y + \frac{y^2}{4}\right) dy \\
&= \int_{0}^{2} \left(y - y^2 + \frac{y^3}{4}\right) dy \\
&= \left[ \frac{y^2}{2} - \frac{y^3}{3} + \frac{y^4}{16} \right]_{0}^{2} \\
&= \left( \frac{4}{2} - \frac{8}{3} + \frac{16}{16} \right) - 0 \\
&= 2 - \frac{8}{3} + 1 = \boxed{\frac{1}{2}}.
\end{aligned}
\]

\item \textbf{Ejemplo 12: Integral doble entre 2 curvas}

\[
\begin{aligned}
\iint_A (2y - x) \, dA &= \int_{x=0}^{1} \left( \int_{y=x^2}^{x} (2y - x) \, dy \right) dx \\
&= \int_{0}^{1} \left[ y^2 - xy \right]_{y=x^2}^{x} dx \\
&= \int_{0}^{1} \left( (x^2 - x \cdot x) - ((x^2)^2 - x \cdot x^2) \right) dx \\
&= \int_{0}^{1} \left( x^2 - x^2 - x^4 + x^3 \right) dx \\
&= \int_{0}^{1} (-x^4 + x^3) \, dx \\
&= \left[ -\frac{x^5}{5} + \frac{x^4}{4} \right]_{0}^{1} \\
&= -\frac{1}{5} + \frac{1}{4} = \frac{-4 + 5}{20} = \boxed{\frac{1}{20}}.
\end{aligned}
\]

\[
\begin{aligned}
\iint_A (2y - x) \, dA &= \int_{y=0}^{1} \left( \int_{x=y}^{\sqrt{y}} (2y - x) \, dx \right) dy \\
&= \int_{0}^{1} \left[ 2y x - \frac{x^2}{2} \right]_{x=y}^{\sqrt{y}} dy \\
&= \int_{0}^{1} \left( 2y \sqrt{y} - \frac{y}{2} - 2y^2 + \frac{y^2}{2} \right) dy \\
&= \int_{0}^{1} \left( 2y^{3/2} - \frac{y}{2} - \frac{3y^2}{2} \right) dy \\
&= \left[ \frac{4}{5} y^{5/2} - \frac{y^2}{4} - y^3 \right]_{0}^{1} \\
&= \frac{4}{5} - \frac{1}{4} - 1 = \frac{16}{20} - \frac{5}{20} - \frac{20}{20} = \boxed{\frac{1}{20}}.
\end{aligned}
\]

\item 
\textbf{Ejemplo 13: integral con $\exp(x^2)$}
%Otro ejemplo interesante.

Descripción del dominio \(A\) reparametrizado:
La región \(A\) está definida por:
- \(y \leq x \leq 2\),
- \(0 \leq y \leq 2\).
Pero si queremos integrar primero en \(y\) y luego en \(x\), debemos expresar los límites de \(y\) en función de \(x\):
- Para \(0 \leq x \leq 2\), \(y\) varía desde \(0\) hasta \(x\).
Integral iterada en el orden \(dy \, dx\).
La integral se expresa como:
\[
\iint_A e^{x^2} \, dA = \int_{x=0}^{2} \left( \int_{y=0}^{x} e^{x^2} \, dy \right) dx.
\]
Aquí, \(e^{x^2}\) no depende de \(y\), por lo que podemos factorizarlo:
\[
\int_{x=0}^{2} e^{x^2} \left( \int_{y=0}^{x} dy \right) dx = \int_{0}^{2} e^{x^2} \cdot x \, dx.
\]
La integral se simplifica a:
\[
\int_{0}^{2} x e^{x^2} \, dx.
\]
Hacemos el cambio de variable \(u = x^2\), entonces \(du = 2x \, dx\) y \(x \, dx = \frac{du}{2}\). Los límites de integración cambian a:
- Cuando \(x = 0\), \(u = 0\),
- Cuando \(x = 2\), \(u = 4\).
Por lo tanto, la integral se convierte en:
\[
\int_{0}^{4} e^{u} \cdot \frac{du}{2} = \frac{1}{2} \int_{0}^{4} e^{u} \, du = \frac{1}{2} \left[ e^{u} \right]_{0}^{4} = \frac{1}{2} (e^{4} - 1).
\]
Resultado final:
\[
\boxed{\frac{e^{4} - 1}{2}}.
\]

\end{enumerate}

\subsection{Cambio de variables (complemento)}

\paragraph{Enunciado: Fórmula de Cambio de Variable en Integrales Dobles}
%
Hipótesis:
1. Sea \( T: \Omega \subset \mathbb{R}^2 \to \mathbb{R}^2 \) una transformación biyectiva y de clase \( C^1 \) (continuamente diferenciable) en un abierto \( \Omega \).
2. Sea \( \Phi = (u(x,y), v(x,y)) \) la transformación, con jacobiano:
   \[
   J_T(x,y) = \det \begin{pmatrix}
   \frac{\partial u}{\partial x} & \frac{\partial u}{\partial y} \\
   \frac{\partial v}{\partial x} & \frac{\partial v}{\partial y}
   \end{pmatrix} \neq 0 \quad \forall (x,y) \in \Omega.
   \]
3. Sea \( f: T(\Omega) \to \mathbb{R} \) una función continua.
Conclusion:
%**Tesis:**
Bajo estas condiciones, si \( R \subset \Omega \) es una región medible y \( T(R) \) es su imagen bajo \( T \), entonces:
\[
\iint_{T(R)} f(u,v) \, du \, dv = \iint_{R} f(T(x,y)) |J_T(x,y)| \, dx \, dy.
\]

\paragraph{Demostración. Paso 1.}
%
%**Paso 1: Descomposición en regiones pequeñas**
Dividimos \( R \) en subrectángulos \( R_{ij} \) de área \( \Delta A_{ij} = \Delta x_i \Delta y_j \). La imagen de cada \( R_{ij} \) bajo \( T \) es aproximadamente un paralelogramo en el plano \( (u,v) \), cuya área es \( |J_T(x_i,y_j)| \Delta A_{ij} \) (por la interpretación geométrica del jacobiano).

\paragraph{Demostración. Paso 2.}
%**Paso 2: Suma de Riemann**
Para cada subrectángulo, evaluamos \( f \) en un punto \( (u_i,v_j) = T(x_i,y_j) \):
\[
\iint_{T(R)} f(u,v) \, du \, dv \approx \sum_{i,j} f(u_i,v_j) |J_T(x_i,y_j)| \Delta A_{ij}.
\]

\paragraph{Demostración. Paso 3.}
%**Paso 3: 
Límite cuando \( \Delta x_i, \Delta y_j \to 0 \).
Al refinar la partición, la suma de Riemann converge a la integral:
\begin{eqnarray*}
&&\iint_{T(R)} f(u,v) \, du \, dv \\&=& \lim_{\Delta x_i, \Delta y_j \to 0} \sum_{i,j} f(T(x_i,y_j)) |J_T(x_i,y_j)| \Delta A_{ij}\\ &=& \iint_{R} f(T(x,y)) |J_T(x,y)| \, dx \, dy.    
\end{eqnarray*}

\paragraph{Demostración. Paso 4.}
Justificación del valor absoluto del jacobiano.
El valor absoluto \( |J_T| \) garantiza que el área sea positiva, independientemente de si \( T \) preserva o invierte la orientación.

\paragraph{Ejemplo Ilustrativo}
Transformación a coordenadas polares:
Sea \( T(x,y) = (r \cos \theta, r \sin \theta) \), con \( r \geq 0 \) y \( \theta \in [0, 2\pi) \). El jacobiano es:
\[
J_T(r,\theta) = \det \begin{pmatrix}
\cos \theta & -r \sin \theta \\
\sin \theta & r \cos \theta
\end{pmatrix} = r.
\]
Por lo tanto:
\[
\iint_{T(R)} f(x,y) \, dx \, dy = \iint_{R} f(r \cos \theta, r \sin \theta) \, r \, dr \, d\theta.
\]

\paragraph{Ver tambien:}
\url{https://www.fing.edu.uy/~eleonora/dvi/IntegralesDobles4.pdf}

\subsection{Aplicaciones en fisica}

\subsubsection{Electromagnetismo}

\paragraph{Ejemplo 1: Potencial eléctrico de un disco cargado}
Contexto: sea un disco de radio \( R \) con densidad superficial de carga uniforme \( \sigma \). Se calcula el potencial eléctrico \( V \) en un punto a lo largo del eje \( z \) (a distancia \( h \) del centro del disco).

\medskip

%
\noindent
\textbf{Integral doble:}
\[
V(h) = \frac{\sigma}{4\pi\epsilon_0} \int_{0}^{2\pi} \int_{0}^{R} \frac{r \, dr \, d\theta}{\sqrt{r^2 + h^2}}.
\]
%
Cálculo paso a paso:

\noindent
1. Integración en \( \theta \):
   \[
   \int_{0}^{2\pi} d\theta = 2\pi.
   \]
2. Integración en \( r \):
   \[
   V(h) = \frac{\sigma}{4\pi\epsilon_0} 2\pi \int_{0}^{R} \frac{r \, dr}{\sqrt{r^2 + h^2}}.
   \]
   Sea \( u = r^2 + h^2 \), entonces \( du = 2r \, dr \):
   \[
   V(h) = \frac{\sigma}{2\epsilon_0} \int_{h^2}^{R^2 + h^2} \frac{du/2}{\sqrt{u}} = \frac{\sigma}{2\epsilon_0} \left[ \sqrt{u} \right]_{h^2}^{R^2 + h^2}.
   \]
3. Resultado:
   \[
   V(h) = \frac{\sigma}{2\epsilon_0} \left( \sqrt{R^2 + h^2} - h \right).
   \]

\paragraph{Ejemplo 2: Campo magnético de una corriente superficial en un disco}
Contexto: Un disco de radio \( R \) con corriente superficial uniforme \( \mathbf{K} = K \hat{\phi} \). Se calcula el campo magnético \( \mathbf{B} \) en un punto del eje \( z \).

\medskip

\noindent En realidad, la sustitución correcta es:
   \[
   \int_{0}^{R} \frac{r^2 \, dr}{(r^2 + h^2)^{3/2}} = \int_{0}^{R} \frac{(r^2 + h^2 - h^2) \, dr}{(r^2 + h^2)^{3/2}} = \int_{0}^{R} \frac{dr}{(r^2 + h^2)^{1/2}} - h^2 \int_{0}^{R} \frac{dr}{(r^2 + h^2)^{3/2}}.
   \]
%
%3. **Integración por partes o %sustitución directa:**
La primera integral es estándar (\url{https://www.physics.umd.edu/hep/drew/IntegralTable.pdf}):
     \[
     \int \frac{dr}{(r^2 + h^2)^{1/2}} = \ln\left(r + \sqrt{r^2 + h^2}\right).
     \]
%
%\begin{verbatim}
%    from sympy import symbols, integrate, sqrt, log
%
%# Define the variables
%r, h = symbols('r h', positive=True)
%
%# Define the expression to integrate
%expression = 1 / sqrt(r**2 + h**2)
%
%# Perform the integration
%integral_result = integrate(expression, r)
%
%# Define the expected result
%expected_result = log(r + sqrt(r**2 + h**2))
%
%# Print the results
%print("Integral result:", integral_result)
%print("Expected result:", expected_result)
%
%# Check if the results are equivalent
%if integral_result.simplify() == expected_result.simplify():
%    print("The results are equivalent.")
%else:
%    print("The results are not equivalent.")
%
%Integral result: asinh(r/h)
%Expected result: log(r + sqrt(h**2 + r**2))
%The results are not equivalent.
%    
%\end{verbatim}
%     
La segunda integral se resuelve con \( r = h \tan \theta \):
     \[
     \int \frac{dr}{(r^2 + h^2)^{3/2}} = \frac{r}{h^2 \sqrt{r^2 + h^2}} + C.
     \]
%
%\begin{verbatim}
%    from sympy import symbols, integrate, sqrt, log
%
%# Define the variables
%r, h = symbols('r h', positive=True)
%
%# Define the expression to integrate
%expression = 1 / (r**2 + h**2)**(3/2)
%
%# Perform the integration
%integral_result = integrate(expression, r)
%
%# Define the expected result
%expected_result = r / (h**2 * sqrt(r**2 + h**2))
%
%# Print the results
%print("Integral result:", integral_result)
%print("Expected result:", expected_result)
%
%# Check if the results are equivalent
%if integral_result.simplify() == expected_result.simplify():
%    print("The results are equivalent.")
%else:
%    print("The results are not equivalent.")
%
%Integral result: 0.564189583547756*sqrt(pi)*r/(h**2.0*sqrt(h**2 + r**2))
%Expected result: r/(h**2*sqrt(h**2 + r**2))
%The results are not equivalent.
%
%    
%\end{verbatim}
%
Evaluación:
   \[
   \left[ \ln\left(r + \sqrt{r^2 + h^2}\right) \right]_{0}^{R} - h^2 \left[ \frac{r}{h^2 \sqrt{r^2 + h^2}} \right]_{0}^{R}.
   \]
   Simplificando:
   \[
   \ln\left(R + \sqrt{R^2 + h^2}\right) - \ln(h) - \frac{R}{\sqrt{R^2 + h^2}}.
   \]
Integración en \( \theta \)
\[
\int_{0}^{2\pi} d\theta = 2\pi.
\]
Resultado final
\[
\mathbf{B}(h) = \frac{\mu_0 K}{2} \left( \ln\left(\frac{R + \sqrt{R^2 + h^2}}{h}\right) - \frac{R}{\sqrt{R^2 + h^2}} \right) \hat{z}.
\]


\subsubsection{Optica} 

\paragraph{Ejemplo 1: Patrón de difracción de Fraunhofer para una apertura rectangular}
Contexto: Una apertura rectangular de dimensiones \( a \times b \) iluminada por luz monocromática de longitud de onda \( \lambda \). Se calcula la amplitud del campo eléctrico en un punto lejano.

\medskip

\noindent \textbf{Integral doble:}
\[
E(\theta, \phi) = E_0 \int_{-a/2}^{a/2} \int_{-b/2}^{b/2} e^{-i k (x \sen\theta + y \sen\phi)} \, dx \, dy.
\]
donde \( k = \frac{2\pi}{\lambda} \).

\medskip

\noindent Cálculo paso a paso:

\noindent 1. Separación de variables:
   \[
   E(\theta, \phi) = E_0 \left( \int_{-a/2}^{a/2} e^{-i k x \sen\theta} \, dx \right) \left( \int_{-b/2}^{b/2} e^{-i k y \sen\phi} \, dy \right).
   \]
2. Cada integral es de la forma:
   \[
   \int_{-L/2}^{L/2} e^{-i \alpha y} dy = L \frac{\sen(\alpha L/2)}{\alpha L/2}.
   \]
3. Resultado:
   \[
   \displaystyle
   E(\theta, \phi) = E_0 a b \frac{\sen\left(\frac{k a \sen\theta}{2}\right)}{\frac{k a \sen\theta}{2}} \frac{\sen\left(\frac{k b \sen\phi}{2}\right)}{\frac{k b \sen\phi}{2}}.
   \]

\paragraph{Ejemplo 2: Intensidad de una lente con aberración esférica}
Contexto: Una lente circular de radio \( R \) con aberración esférica. La fase introducida es \( \psi(r) = \alpha r^4 \).

\medskip

\noindent
\textbf{Integral doble para la amplitud:}
\[
A(P) = \int_{0}^{R} \int_{0}^{2\pi} e^{i \alpha r^4} r \, dr \, d\theta
\]

\medskip

\noindent
Cálculo paso a paso:
1. Integración en \( \theta \):
   \[
   \int_{0}^{2\pi} d\theta = 2\pi
   \]
2. Integración en \( r \):
   \[
   A(P) = 2\pi \int_{0}^{R} e^{i \alpha r^4} r \, dr.
   \]
   %Sea \( u = \alpha r^4 \), \( du = 4\alpha r^3 dr \), pero 
   No tiene solución analítica simple en términos de funciones elementales. %Para aberraciones pequeñas, 
   Se usa aproximación numérica o series.

\subsubsection{Fluidos}

\paragraph{Ejemplo 1: Circulación de un vórtice potencial en 2D} 
Contexto: Un vórtice potencial con vorticidad \( \omega \) y circulación \( \Gamma \). El campo de velocidades es:
\[
\mathbf{v}(r, \theta) = \frac{\Gamma}{2\pi r} \hat{\theta}.
\]
Se calcula el flujo a través de un anillo de radio \( a \) y ancho \( dr \).

\medskip

\noindent \textbf{Integral doble (flujo):}
\[
Q = \int_{0}^{2\pi} \int_{a}^{a + dr} \mathbf{v} \cdot \hat{r} \, r \, dr \, d\theta = 0.
\]
Interpretación: El flujo es cero porque \( \mathbf{v} \) es tangencial.


\paragraph{Ejemplo 2: Energía cinética de un fluido en un dominio rectangular}
Contexto: Flujo en 2D con velocidad \( \mathbf{v}(x,y) = (u(x,y), v(x,y)) \). La energía cinética por unidad de profundidad es:
\[
E = \frac{\rho}{2} \int_{0}^{L_y} \int_{0}^{L_x} (u^2 + v^2) \, dx \, dy.
\]
Ejemplo concreto: Flujo de Couette entre placas paralelas:
\[
u(y) = \frac{U y}{h}, \quad v = 0,
\]
\[
E = \frac{\rho}{2} \int_{0}^{L_x} \int_{0}^{h} \left( \frac{U y}{h} \right)^2 \, dy \, dx = \frac{\rho U^2 L_x h}{6}.
\]


%\begin{center}
%    $\star$ versi\'on 3 -- 27/10/25 $\star$
%\end{center}

%\chapter{Electromagnetismo (epilogo)}

%
%### Comparación con el resultado clásico
%El resultado clásico para un **disco de corriente superficial** es:
%\[
%\mathbf{B}(h) = \frac{\mu_0 K}{2} \left( 1 - \frac{h}{\sqrt{R^2 + h^2}} \right) \hat{z},
%\]
%pero este corresponde a una **espira circular**, no a un disco.
%
%### Conclusión
%El error en tu desarrollo original está en la sustitución: **no puedes reemplazar \( r^2 \, dr \) directamente por \( (u - h^2) \, du/2 \)**, porque \( r^2 \, dr = r \cdot r \, dr \), y \( r \, dr = du/2 \), pero \( r \) no es constante. La integral correcta para un disco es más compleja y no da el resultado que mencionaste.
%
%**Si el problema es para una espira circular (corriente concentrada en \( r = R \)), entonces la integral es distinta y el resultado es el que mencionaste.**

\subsection{Integracion numerica}
\subsubsection{Introducción}
La integración numérica es una herramienta fundamental en el análisis numérico, especialmente cuando se trata de integrales que no admiten una solución analítica cerrada o cuando los integrandos son complejos. En este documento, exploramos dos familias de métodos clásicos: el \textbf{método del trapecio} (simple y compuesto) y los \textbf{métodos de cuadratura de Gauss} (Gauss-Legendre), tanto en una como en dos dimensiones. Además, analizamos las cotas de error asociadas y proporcionamos ejemplos ilustrativos.

\subsubsection{Métodos de Integración en 1D}

\paragraph{Método del Trapecio Simple}
Dada una función \( f \) integrable en \([a, b]\), la regla del trapecio aproxima la integral como el área de un trapecio:
\[
\int_{a}^{b} f(x) \, dx \simeq T(f) = (b - a) \cdot \frac{f(a) + f(b)}{2}
\]

\paragraph{Cota de Error}
El error de truncamiento \( E_T \) para el método del trapecio simple está dado por:
\[
E_T = -\frac{(b - a)^3}{12} f''(\xi), \quad \xi \in (a, b)
\]
Esto requiere que \( f \) sea dos veces diferenciable en \([a, b]\).

\paragraph{Método del Trapecio Compuesto}
Para mejorar la precisión, se divide el intervalo \([a, b]\) en \( n \) subintervalos de igual longitud \( h = \frac{b - a}{n} \). La regla compuesta es:
\[
T_n(f) = \frac{h}{2} \left[ f(a) + 2 \sum_{i=1}^{n-1} f(a + ih) + f(b) \right]
\]

\paragraph{Cota de Error}
El error para el trapecio compuesto es:

\[
E_T = -\frac{(b - a) h^2}{12} f''(\xi), \quad \xi \in (a, b)
\]

\paragraph{Cuadratura de Gauss-Legendre}
La cuadratura de Gauss aproxima la integral utilizando nodos y pesos específicos:
\[
\int_{-1}^{1} f(x) \, dx \simeq \sum_{i=1}^{n} w_i f(x_i)
\]
donde \( x_i \) son las raíces del polinomio de Legendre de grado \( n \) y \( w_i \) son los pesos asociados. Para un intervalo \([a, b]\), se aplica un cambio de variable:
\[
\int_{a}^{b} f(x) \, dx = \frac{b - a}{2} \int_{-1}^{1} f\left(\frac{b - a}{2} t + \frac{a + b}{2}\right) dt
\]
\paragraph{Cota de Error}
El error para la cuadratura de Gauss de \( n \) puntos es:

\[
E_G = \frac{2^{2n+1} (n!)^4}{(2n+1) [(2n)!]^3} f^{(2n)}(\xi), \quad \xi \in (-1, 1)
\]

\subsubsection{Métodos de Integración en 2D}

\paragraph{Regla del Trapecio en 2D}
Para una función \( f(x, y) \) sobre un dominio rectangular \([a, b] \times [c, d]\), la regla del trapecio en 2D se construye aplicando el trapecio en cada dirección:
\[
\iint_{[a,b] \times [c,d]} f(x, y) \, dx dy \simeq \frac{(b - a)(d - c)}{4} \left[ f(a, c) + f(a, d) + f(b, c) + f(b, d) \right]
\]
\paragraph{Cota de Error}
El error es análogo al caso 1D, pero depende de las derivadas parciales mixtas de \( f \).

\paragraph{Cuadratura de Gauss en 2D}
Se extiende la cuadratura de Gauss a dos dimensiones utilizando productos tensoriales de las reglas 1D:
\[
\iint_{[-1,1] \times [-1,1]} f(x, y) \, dx dy \simeq \sum_{i=1}^{n} \sum_{j=1}^{m} w_i w_j f(x_i, y_j)
\]
Para dominios rectangulares \([a, b] \times [c, d]\), se aplica un cambio de variable similar al caso 1D.

\subsubsection{Ejemplos}

\paragraph{Ejemplo 1D: Método del Trapecio}
Consideremos \( f(x) = e^x \) en \([0, 1]\). La integral exacta es \( e - 1 \approx 1.71828 \).
\textbf{Trapecio simple}:
  \[
  T(f) = \frac{1 - 0}{2} (f(0) + f(1)) = \frac{1}{2} (1 + e) \approx 1.85914
  \]
Error: \( |1.85914 - 1.71828| \approx 0.14086 \).
\textbf{Trapecio compuesto} con \( n = 4 \):
  \[
  T_4(f) = \frac{0.25}{2} \left[ f(0) + 2(f(0.25) + f(0.5) + f(0.75)) + f(1) \right] \approx 1.7254
  \]
Error: \( |1.7254 - 1.71828| \approx 0.00712 \).

\paragraph{Ejemplo 1D: Cuadratura de Gauss}
Para \( f(x) = e^x \) en \([0, 1]\), aplicamos Gauss-Legendre con \( n = 2 \).
Nodos y pesos: \( x_1 = -0.57735 \), \( x_2 = 0.57735 \), \( w_1 = w_2 = 1 \).
Cambio de variable: \( t = 2x - 1 \), \( x = \frac{t + 1}{2} \).
Aproximación:
  \[
  \int_{0}^{1} e^x \, dx \approx \frac{1}{2} \sum_{i=1}^{2} w_i e^{\frac{x_i + 1}{2}} \approx 1.71828
  \]
Error: \( \approx 0 \). % (exacto para polinomios de grado \( \leq 3 \)).

\paragraph{Ejemplo 2D: Trapecio}
Para \( f(x, y) = xy \) en \([0, 1] \times [0, 1]\), la integral exacta es \( \frac{1}{4} \).
\textbf{Trapecio simple}:
  \[
  T(f) = \frac{1}{4} (f(0,0) + f(0,1) + f(1,0) + f(1,1)) = \frac{1}{4} (0 + 0 + 0 + 1) = 0.25
  \]
Error: \( 0 \) (exacto para funciones bilineales).

\subsubsection{Observaciones} 
\noindent
El método del trapecio es simple pero menos preciso para funciones con alta curvatura.
La cuadratura de Gauss es más eficiente para funciones suaves, especialmente con pocos nodos.
En 2D, la extensión de estos métodos es directa, pero el costo computacional aumenta con la dimensionalidad.

\medskip 

\noindent
\textbf{Nota:} Para implementaciones prácticas, se recomienda utilizar bibliotecas numéricas como \texttt{scipy.integrate} en Python, que implementan estos métodos de manera optimizada.

\chapter{Curvas en el espacio} % y geometria diferencial de curvas}

\section{Temario}

\begin{itemize}
    \item Continuidad y derivabilidad. 
    \item Regla de la cadena. Longitud.
    \item Triedro de Frenet. \item Curvatura y torsión. \item Ecuaciones de Frenet.

\end{itemize}

\section{Notas}

\noindent{Notas sobre curvas en el espacio, triedro de Frenet y ecuaciones de Frenet.}

\subsection{Motivación}

Curvas en el espacio sirven para muchisimas cosas.
de forma general, pueden ser un primer escalon para el estudio de la geometria diferencial, ver el libro de Mandredo Do Carmo (Geometria de Curvas y Superficies) por ejemplo. El primer capitulo del Do Carmo es la base de estos apuntes. %~\cite{abraham,doCarmo1976,pressley2010,banchoff2016,spivak1999,shifrin2016}. % la %geometría diferencial es fundamental para el estudio de curvas y superficies.
%TAmbien es la base geometrica para muchos modelos de
%
%(ref). 
%\cite{abraham}
%
%
%Marsden
%DO CARMO CURVES
%https://books.google.com.uy/books/about/%Differential_Geometry_of_Curves_and_Surf.html?id=uXF6DQAAQBAJ&redir_esc=y
%
Ademas la geometria de curvas en tres dimensiones es la base geometrica de muchos modelos de varillas
(\emph{rods} en ingles)\footnote{Ver \url{https://hal.science/cel-01023392/}}. Estos modelos son de mucha utilidad, por ejemplo 
%\begin{itemize}
%    \item 
en computacion grafica, para simular el pelo de personajes (\url{https://elan.inrialpes.fr/people/bertails/}).
%    \item 
%\end{itemize}



%\url{http://www.lmm.jussieu.fr/~neukirch/}

%\url{http://goriely.com/}

%\url{https://cde.nus.edu.sg/me/staff/cecilia-laschi/}

%Cristian Duriez

%Robots Blandos

%\emph{soft robots}"** se traduce como **"robots blandos"** .

%En cuanto a investigadores destacados en el modelado de robots blandos que utilizan estructuras tipo "rod" (barras), uno de los nombres más relevantes es *
%*Cecilia Laschi**, profesora de la Universidad de Pisa y experta en robótica blanda, conocida por su trabajo en robots inspirados en pulpos y el uso de estructuras flexibles y actuadores blandos. También destacan **Prashant Mehta** y **Mattia Gazzola** de la Universidad de Illinois, quienes han desarrollado modelos computacionales avanzados para robots blandos inspirados en la biomecánica de los tentáculos de pulpos, incluyendo el uso de elementos estructurales flexibles .

%Computación grafica, biofisica
%Neukirsch Goriely
%, robotica Kanty, etc.

%\begin{itemize}
%    \item 
%\end{itemize}

%Introduccion a la geometria diferencial. 

\subsection{Definiciones y nociones basicas}
\noindent Una curva en el espacio se puede definir como una función vectorial de un parámetro \(t\):
\[
\mathbf{r} : I \rightarrow \mathbb{R}^3
\]
donde $I$ es un intervalo de $\mathbb{R}$. Lo %\dots
podemos escribir
\[ \mathbf{r}(t) = (x(t), y(t), z(t)) \]
donde \(x(t), y(t), z(t)\) son funciones reales de \(t\).

\begin{example}
Un helice se puede definir como
\[
x(t)=\sqrt{t} \cos(5t), \quad
y(t)=\sqrt{t} \sen(5t), \quad
z(t)=\sqrt{t}
\] con
$I=[0;2\pi]$.
\end{example}

\begin{example}
Otro tipo de helice, mas simple
\[
x(t)=\cos(t), \quad
y(t)=\sen(t), \quad
z(t)=t,
\]
con
$I=[0;2\pi]$.
\end{example}

\noindent Nos van a interesar curvas
$\mathbf{r}$ regulares: derivables por lo menos y con derivada
$\mathbf{r}'(t)\neq \mathbf{0}$,
donde \[ \mathbf{r}'(t) = (x'(t), y'(t), z'(t)). \]

\subsection{Parametrizaci\'on}

\noindent
Sea $J$ otro intervalo de $\mathbb{R}$ y sea $\tau:J\rightarrow I$ un difeomorfismo creciente. Definimos
$t=\tau(s)$ y
$$\boldsymbol{\rho}(s)=\mathbf{r}(\tau(s))$$ es una reparametrizaci\'on de la curva. Deducimos de la relaci\'on anterior
$$
\boldsymbol{\rho}'(s)=\mathbf{r}'(\tau(s))\tau'(s)$$
donde hemos utilizado la regla de la cadena.
Con la regla de la cadena, notamos también que
\[
\mathrm{d}t=\tau'(s) \mathrm{d}s, \quad \tau'(s) = \frac{\mathrm{d}t}{\mathrm{d}s}.
\]
De esto se deduce
\[
\boldsymbol{\rho}'(s)=\mathbf{r}'(t) \frac{\mathrm{d}t}{\mathrm{d}s}.
\]
\begin{example}
Con la helice
\[
x(t)=\cos(t), \quad
y(t)=\sen(t), \quad
z(t)=t,
\]
y
$I=[0;2\pi]$, hacemos el cambio siguiente
\[
t = 2\pi s
\]
y $s\in [0;1]$. Uno nota que las derivadas son diferentes: son las mismas curvas, pero recorridas a velocidad distinta.
\end{example}

\begin{example}
El semi-circulo se puede definir con coordenadas cartesianas
\[
\mathbf{r}(t) = ( t , \sqrt{1-t^2} , 0 )
\]
con $t\in [-1;1]$
o coordenadas polares
\[
\boldsymbol{\rho}(s) = ( -\cos(s) , \sen(s) , 0 )
\]
y $s\in [0;\pi]$, dado que
\[
\cos(\pi-s)=-\cos(s), \quad \cos(0)=1, \quad \cos(\pi)=-1.
\]
Tenemos
\[
\tau(s) = -\cos(s),
\]
que es creciente en $[0;\pi]$.
\end{example}

\subsection{Longitud de arco}

\noindent La longitud de arco se define a nivel infinitesimal asi
\[
\mathrm{d}l = 
\| \mathrm{d} \mathbf{r} \|
= 
\| \mathbf{r}'(t) \mathrm{d}t  \|
= 
\| \mathbf{r}'(t)   \| \mathrm{d}t,
\]
para una variaci\'on infinitesimal $\mathrm{d}t$ del parametro $t$. Se puede considerar $\mathbf{r}'(t)$ como una velocidad instantanea. Despues se puede integrar asi:
\[
l = \int \mathrm{d}l = \int_I \| \mathbf{r}'(t)   \| \mathrm{d}t.
\]

\begin{ejemplo}
    Para un circulo de radio $R$, tenemos $t\in [0,2\pi]$ y
    \[ \mathbf{r}(t) = ( R \cos(t), R\sen(t) , 0 )\]
    y entonces
    \[
    \| \mathbf{r}'(t) \| = R.
    \]
    Si integramos volvemos a obtener
    \[
    l = 2\pi R.
    \]
\end{ejemplo}

\paragraph{La longitud de arco no depende de la parametrizci\'on}
Sea de nuevo 
$$\boldsymbol{\rho}(s)=\mathbf{r}(\tau(s))$$ una reparametrizaci\'on de la curva.
Utilizamos la regla de la cadena y observamos que
\[
\| \boldsymbol{\rho}'(s)   \| \mathrm{d}s
= 
\| \mathbf{r}'(\tau(s)) \| \tau'(s) \mathrm{d}s
%= 
%\| \mathbf{r}'(\tau(s)) \| \tau'(s) \mathrm{d}s
=
\| \mathbf{r}'(t) \| \mathrm{d}t =
\mathrm{dl}
\]
donde hemos utilizado $\tau(s)=t$ y $\tau'(s) \mathrm{d}s = \mathrm{d}t$.
%\noindent La longitud de arco $dl = \| \gamma'(t) \| dt$ es
%independiente de la parametrizacion. 

\paragraph{Parametrizaci\'on de longitud de arco (PDLA)} Definimos
$$
h(s) = \int_0^s \| \mathbf{r}'(t) \| \mathrm{d}t,$$ que es derivable y creciente. Su deriva es
\[
h'(s) = \| \mathbf{r}'(s) \|
\]
Tomamos ahora la reparametrizaci\'on $\tau(s)=h^{-1}(s)$ ($h$ es invertible), que verifica tambi\'en las condiciones de derivabilidad. Calculamos
$$\tau'(s)=(h^{-1})'(s)=\frac{1}{h'(\tau(s))}
=
\frac{1}{\| \mathbf{r}'(\tau(s)) \|}.$$
Sea ahora 
$$
\boldsymbol{\alpha}(s) = \mathbf{r} (\tau(s)).$$
%
Calculamos de nuevo con la regla de la cadena
\[
\| \boldsymbol{\alpha}'(s) \|= \| \mathbf{r}'  (\tau(s))\| \tau'(s).
\]
Y utilizando la formula anterior, obtenemos
\[
\| \boldsymbol{\alpha}'(s) \|= \| \mathbf{r}'  (\tau(s))\| \frac{1}{\| \mathbf{r}'(\tau(s)) \|} =1.
\]
Permite tener una parametrizaci\'on con velocidad constante y equal a $1$, y donde el parametro es la distancia recorrida. En general, esta parametrizacion no se puede calcular de forma explicita, pero es util para definir la curvatura y la torsion. 

\begin{ejemplo}
    Para una recta o un circulo, es facil definir la PDLA.
\end{ejemplo}

%La longitud de arco no depende de la 
%\dots Es unica \dots Permite definir la longitud (mapea la curva en una linea de mismo tamano) \dots
%\medskip

%\noindent 
%
%Parametrizacion de longitud de arco (plda) corresponde, para todo $s$: $\| \alpha'(s) \|=1$.
%
%Teorema: siempre se puede definir. Tomamos

%A insertar despues: long de un grafico, catenaria, elipse y integral eliptica

%\medskip
%Espiral log
%$(e^t \cos(t), e^t \sen(t))$ 
%$\tau(s)= \log(s/\sqrt(2))$. VER SEBASTIAN

%\begin{example}
%    Catenaria
%\end{example}

%
%\noindent 
%Ejemplos:
%semi circulo, catenaria. % $f(t)=\cosh{t}$.

\subsection{Vector tangente y triedro de Frenet}

\noindent
Sea de nuevo una curva $\boldsymbol{\alpha}$
definida con parametrización por longitud de arco.

\subsubsection{Vector tangente}

\noindent
Podemos definir 
%Curvatura 
$$\mathbf{T}(s) = \boldsymbol{\alpha}'(s)$$el vector tangente 
a la curva en el punto $s$. 
Debido a la PDLA,
este vector es unitario
\[
\| \mathbf{T}(s) \|
=
\| \boldsymbol{\alpha}'(s) \|
=1.
\]
De hecho, si hacemos un Taylor cerca de un punto $s$ en la curva, obtenemos:
\[
\boldsymbol{\alpha} (s+\delta s) =
\boldsymbol{\alpha} (s)
+ \boldsymbol{\alpha}'(s) \delta s
+ \mathcal{O}(\delta s^2)
= 
\boldsymbol{\alpha} (s)
+ \mathbf{T}(s) \delta s
+ \mathcal{O}(\delta s^2)
\]
lo que da $\boldsymbol{\alpha} (s+\delta s) \simeq \boldsymbol{\alpha} (s)
+ \mathbf{T}(s) \delta s$ y permite visualizar que $\mathbf{T}$ es geometricamente tangente a la curva.

\subsubsection{Vector acceleración y vector normal}

\noindent
El vector acceleraci\'on se define simplemente como la derivada del vector tangente (que es un vector de velocidad en cierto sentido):
$$\mathbf{A}(s) = \mathbf{T}'(s) = \boldsymbol{\alpha}"(s).$$ Dado que la velocidad es constante en norma, la acceleracion es perpendicular a la curva. De hecho, lo podemos mostrar de la siguiente forma. 
%
Partemos de
$$\frac{1}{2}=\frac{\| \boldsymbol{\alpha}'(s) \|^2}{2}$$ y derivamos. Obtenemos
\[
0 = \langle \boldsymbol{\alpha}'(s) , \boldsymbol{\alpha}" (s) \rangle
\]
y entonces
\[
\langle \mathbf{A}(s) , \mathbf{N}(s) \rangle = 0.
\]
%Prop N orto a T porque 
%*
El vector acceleraci\'on puede tener cualquier norma.
Si es diferente de cero, definimos el vector normal como el vector acceleraci\'on normalizado:
\[
\mathbf{N}(s) = \frac{\mathbf{A}(s)}{\| \mathbf{A}(s) \|}
\]
y asi
\[
\| \mathbf{N}(s) \| = 1.
\]

\subsubsection{Vector binormal y triedro de Frenet}

\noindent
El vector binormal se define como:
\[ \mathbf{B}(s) = \mathbf{T}(s) \times \mathbf{N}(s) \]
%
Es tambi\'en de norma 1.
Estos tres vectores forman una base ortonormal local conocida como \emph{el triedro de Frenet}.

\subsubsection{Ejemplos}

\noindent Para una recta o un circulo, es facil calcular el triedro de Frenet.

\subsection{Curvatura y torsión}

\noindent La curvatura y la torsi\'on son las nociones adecuadas para medir la propriedades geometricas (y entonces kinematicas) de la curva. Los vamos a definir primero con la PDLA.

\subsubsection{Curvatura}

\noindent La curvatura \(\kappa\) mide cuánto se curva localmente la curva y se define como:
\[ \kappa(s) = \|\mathbf{A}(s)\| \]
lo que es equivalente a 
\[ \kappa(s) = \|\mathbf{T}'(s)\| 
= \|\boldsymbol{\alpha}"(s)\|. 
\]
%

\begin{example}
    La curvatura de una recta es $0$ y la de un circulo de radio $R$ es $1/R$.
\end{example}

\subsubsection{Torsi\'on}

\noindent 
Se puede establecer que
\[
\mathbf{B}'(s) = -\tau(s) \mathbf{N}(s)
\]
y la cuantidad $\tau(s)$ se llama la torsi\'on de la curva.
La torsión \(\tau (s)\) mide cuánto se retuerce la curva fuera del plano definido por \(\mathbf{T}\) y \(\mathbf{N}\).

\noindent Partemos de $\| \mathbf{B}(s) \|=1$ y derivamos. Obtenemos
\[
\langle \mathbf{B}'(s) , \mathbf{B}(s) \rangle = 0.
\]
Por otro lado, retomamos la definici\'on del vector binormal
\[ \mathbf{B}(s) = \mathbf{T}(s) \times \mathbf{N}(s) \]
y la derivamos utilizando la regla de Leibniz
\[ \mathbf{B}'(s) = \mathbf{T}'(s) \times \mathbf{N}(s)
+
\mathbf{T}(s) \times \mathbf{N}'(s).
\]
Como los vectores $\mathbf{T}'$ y $\mathbf{N}$ son colineales,
se obtiene
\[ \mathbf{B}'(s) = 
\mathbf{T}(s) \times \mathbf{N}'(s).
\]
Entonces
\[
\langle \mathbf{B}'(s) , \mathbf{T}(s) \rangle =
\langle \mathbf{T}(s) \times \mathbf{N}'(s) , \mathbf{T}(s) \rangle
= 0.
\]
y tenemos $\mathbf{B}'$ que tiene que ser colineal a $\mathbf{N}$.
\begin{example}
Curvas en el plano tienen una torsi\'on nula. Una helice tiene una torsi\'on constante. 
\end{example}

\subsection{Ecuaciones de Frenet-Serret}
\noindent Las ecuaciones de Frenet-Serret describen cómo cambian los vectores del triedro a lo largo de la curva:
\[
\begin{cases}
\mathbf{T}'(s) = \kappa(s) \mathbf{N}(s) \\
\mathbf{N}'(s) = -\kappa(s) \mathbf{T}(s) + \tau(s) \mathbf{B}(s) \\
\mathbf{B}'(s) = -\tau(s) \mathbf{N}(s)
\end{cases}
\]
La primera ecuaci\'on es una reformulaci\'on de
\[
\mathbf{N}(s) = \frac{\mathbf{A}(s)}{\| \mathbf{A}(s) \|}  = \frac{\mathbf{T}'(s)}{\| \mathbf{T}'(s) \|} =
\frac{\mathbf{T}'(s)}{\kappa(s)}.
\]
La tercera es por definici\'on de la torsi\'on.
En el triedro de Frenet:
\[
\mathbf{N}(s) = \mathbf{B}(s) \times \mathbf{T}(s).
\]
Derivamos utilizando la regla de Leibniz:
\[
\mathbf{N}'(s) = \mathbf{B}'(s) \times \mathbf{T}(s)
+ \mathbf{B}(s) \times \mathbf{T}'(s).
\]
Y aplicamos las otras ecuaciones:
\[
\mathbf{N}'(s) = -\tau(s) \mathbf{N}(s) \times \mathbf{T}(s)
+ \kappa(s) \mathbf{B}(s) \times \mathbf{N}(s).
\]

\[
\mathbf{N}'(s) =
-\kappa(s) \mathbf{T}(s) + \tau(s) \mathbf{B}(s).
\]

\subsection{Ejemplo: Hélice circular}

\noindent
Consideremos la hélice circular:
\[ \mathbf{r}(t) = (\cos t, \sin t, t). \]

\subsubsection*{Vectores del triedro de Frenet}
1. Vector tangente:
\[ \mathbf{r}'(t) = (-\sin t, \cos t, 1) \]
\[ |\mathbf{r}'(t)| = \sqrt{\sin^2 t + \cos^2 t + 1} = \sqrt{2} \]
\[ \mathbf{T}(t) = \frac{1}{\sqrt{2}} (-\sin t, \cos t, 1) \]

2. Vector normal:
\[ \mathbf{T}'(t) = \frac{1}{\sqrt{2}} (-\cos t, -\sin t, 0) \]
\[ |\mathbf{T}'(t)| = \frac{1}{\sqrt{2}} \]
\[ \mathbf{N}(t) = (-\cos t, -\sin t, 0) \]

3. Vector binormal:
\[ \mathbf{B}(t) = \mathbf{T}(t) \times \mathbf{N}(t) = \frac{1}{\sqrt{2}} (\sin t, -\cos t, 1) \]

\subsubsection*{Curvatura y torsión}
\noindent La curvatura es:
\[ \kappa = \frac{1}{2}. \]
La torsión es:
\[ \tau = \frac{1}{2}. \]

% Histoire et applications

%\begin{center}
%    $\star$ versi\'on 2 -- 10/11/25 $\star$
%\end{center}

%\chapter{Curvas (bis)}

%\section{Ejemplos}
%Tres ejemplos. Recta, circulo, espiral (Rcos t, R sen t, b t). Dar:
%Parametrizacion de longitud de arco, expresion del triedro de frenet, de la curvatura y de la torsion

%\noindent
%Se detallan las parametrizaciones por longitud de arco, el triedro de Frenet, y los cálculos explícitos de la curvatura y la torsión para tres curvas clásicas: una recta, un círculo de radio \(R\), y una espiral cilíndrica.

\subsection{Recta en \(\mathbb{R}^3\)}

\subsubsection{Parametrización general}
\[
\mathbf{r}(t) = \mathbf{r}_0 + t\mathbf{v}, \quad t \in \mathbb{R}
\]
donde \(\mathbf{r}_0\) es un punto de la recta y \(\mathbf{v}\) es el vector dirección.

\subsubsection{Parametrización por longitud de arco}
\noindent Si \(\|\mathbf{v}\| = 1\), entonces la parametrización por longitud de arco es:
\[
\mathbf{r}(s) = \mathbf{r}_0 + s\mathbf{v}
\]

\subsubsection{Triedro de Frenet}
\begin{itemize}
    \item \(\mathbf{T}(s) = \mathbf{v}\)
    \item \(\mathbf{N}(s)\): Cualquier vector unitario perpendicular a \(\mathbf{v}\) (no está definido de forma única).
    \item \(\mathbf{B}(s) = \mathbf{v} \times \mathbf{N}(s)\)
\end{itemize}

\subsubsection{Curvatura y torsión}
\[
\kappa(s) = 0, \quad \tau(s) = 0
\]

\subsection{Círculo de radio \(R\) en el plano \(xy\)}

\subsubsection{Parametrización general}
\[
\mathbf{r}(t) = (R\cos t, R\sin t, 0), \quad t \in [0, 2\pi]
\]

\subsubsection{Parametrización por longitud de arco}
\noindent La velocidad es:
\[
\mathbf{r}'(t) = (-R\sin t, R\cos t, 0)
\]
Su norma es:
\[
\|\mathbf{r}'(t)\| = R
\]
Por lo tanto, la longitud de arco es:
\[
s(t) = \int_0^t R \, du = Rt
\]
Invirtiendo, \(t = \frac{s}{R}\), y la parametrización por longitud de arco es:
\[
\mathbf{r}(s) = \left(R\cos\left(\frac{s}{R}\right), R\sin\left(\frac{s}{R}\right), 0\right)
\]

\subsubsection{Triedro de Frenet}
\begin{itemize}
    \item \( \displaystyle \mathbf{T}(s) = \left(-\sin\left(\frac{s}{R}\right), \cos\left(\frac{s}{R}\right), 0\right)\)
    \item \(\displaystyle \mathbf{N}(s) = \left(-\cos\left(\frac{s}{R}\right), -\sin\left(\frac{s}{R}\right), 0\right)\)
    \item \(\mathbf{B}(s) = (0, 0, 1)\)
\end{itemize}

\subsubsection{Cálculo de la curvatura}
\noindent La curvatura se define como:
\[
\kappa(s) = \|\mathbf{T}'(s)\|
\]
Calculamos \(\mathbf{T}'(s)\):
\[
\mathbf{T}'(s) = \left(-\frac{1}{R}\cos\left(\frac{s}{R}\right), -\frac{1}{R}\sin\left(\frac{s}{R}\right), 0\right)
\]
Su norma es:
\[
\|\mathbf{T}'(s)\| = \frac{1}{R}
\]
Por lo tanto:
\[
\kappa(s) = \frac{1}{R}
\]

\subsubsection{Cálculo de la torsión}
\noindent La torsión se define como:
\[
\tau(s) = -\mathbf{N}'(s) \cdot \mathbf{B}(s)
\]
Calculamos \(\mathbf{N}'(s)\):
\[
\mathbf{N}'(s) = \left(\frac{1}{R}\sin\left(\frac{s}{R}\right), -\frac{1}{R}\cos\left(\frac{s}{R}\right), 0\right)
\]
El producto punto con \(\mathbf{B}(s)\) es:
\[
\mathbf{N}'(s) \cdot \mathbf{B}(s) = 0
\]
Por lo tanto:
\[
\tau(s) = 0
\]

\subsection{Espiral cilíndrica: \(\mathbf{r}(t) = (R\cos t, R\sin t, bt)\)}

\subsubsection{Parametrización general}
\[
\mathbf{r}(t) = (R\cos t, R\sin t, bt), \quad t \in \mathbb{R}
\]

\subsubsection{Parametrización por longitud de arco}
\noindent
La velocidad es:
\[
\mathbf{r}'(t) = (-R\sin t, R\cos t, b)
\]
Su norma es:
\[
\|\mathbf{r}'(t)\| = \sqrt{R^2 + b^2}
\]
Por lo tanto, la longitud de arco es:
\[
s(t) = \int_0^t \sqrt{R^2 + b^2} \, du = t\sqrt{R^2 + b^2}
\]
Invirtiendo, \(t = \frac{s}{\sqrt{R^2 + b^2}}\), y la parametrización por longitud de arco es:
\[
\mathbf{r}(s) = \left(R\cos\left(\frac{s}{\sqrt{R^2 + b^2}}\right), R\sin\left(\frac{s}{\sqrt{R^2 + b^2}}\right), \frac{bs}{\sqrt{R^2 + b^2}}\right)
\]

\subsubsection{Triedro de Frenet}
\begin{itemize}
    \item \(\displaystyle \mathbf{T}(s) = \frac{1}{\sqrt{R^2 + b^2}}(-R\sin t, R\cos t, b)\)
    \item \(\displaystyle \mathbf{N}(s) = (-\cos t, -\sin t, 0)\)
    \item \(\displaystyle \mathbf{B}(s) = \frac{1}{\sqrt{R^2 + b^2}}(b\sin t, -b\cos t, R)\)
\end{itemize}

\subsubsection{Cálculo de la curvatura}
\noindent
Calculamos \(\mathbf{T}'(s)\):
\[
\mathbf{T}'(s) = \frac{d\mathbf{T}}{dt} \cdot \frac{dt}{ds} = \frac{1}{R^2 + b^2}(-R\cos t, -R\sin t, 0)
\]
Su norma es:
\[
\|\mathbf{T}'(s)\| = \frac{R}{R^2 + b^2}
\]
Por lo tanto:
\[
\kappa(s) = \frac{R}{R^2 + b^2}
\]

\subsubsection{Cálculo de la torsión}
\noindent
Calculamos \(\mathbf{B}'(s)\):
\[
\mathbf{B}'(s) = \frac{d\mathbf{B}}{dt} \cdot \frac{dt}{ds} = \frac{1}{R^2 + b^2}(b\cos t, b\sin t, 0)
\]
La torsión es:
\[
\tau(s) = -\mathbf{N}'(s) \cdot \mathbf{B}(s)
\]
Calculamos \(\mathbf{N}'(s)\):
\[
\mathbf{N}'(s) = \frac{d\mathbf{N}}{dt} \cdot \frac{dt}{ds} = \frac{1}{\sqrt{R^2 + b^2}}(\sin t, -\cos t, 0)
\]
El producto punto \(\mathbf{N}'(s) \cdot \mathbf{B}(s)\) es:
\[
\mathbf{N}'(s) \cdot \mathbf{B}(s) = \frac{1}{R^2 + b^2}(b\sin^2 t + b\cos^2 t) = \frac{b}{R^2 + b^2}
\]
Por lo tanto:
\[
\tau(s) = -\frac{b}{R^2 + b^2}
\]

\subsection{Espiral logarítmica: \(\mathbf{r}(t) = (e^t \cos t, e^t \sin t)\)}

\subsubsection{Parametrización general}
\[
\mathbf{r}(t) = (e^t \cos t, e^t \sin t), \quad t \in \mathbb{R}
\]

\subsubsection{Parametrización por longitud de arco}
\noindent La velocidad es:
\[
\mathbf{r}'(t) = (e^t \cos t - e^t \sin t, e^t \sin t + e^t \cos t) = e^t (\cos t - \sin t, \sin t + \cos t)
\]
Su norma es:
\[
\|\mathbf{r}'(t)\| = e^t \sqrt{(\cos t - \sin t)^2 + (\sin t + \cos t)^2} = e^t \sqrt{2}
\]
Por lo tanto, la longitud de arco es:
\[
s(t) = \int_0^t e^u \sqrt{2} \, du = \sqrt{2} (e^t - 1)
\]
Invirtiendo, \(t = \ln\left(\frac{s}{\sqrt{2}} + 1\right)\), y la parametrización por longitud de arco es:
\[
\mathbf{r}(s) = \left(\left(\frac{s}{\sqrt{2}} + 1\right) \cos\left(\ln\left(\frac{s}{\sqrt{2}} + 1\right)\right), \left(\frac{s}{\sqrt{2}} + 1\right) \sin\left(\ln\left(\frac{s}{\sqrt{2}} + 1\right)\right)\right)
\]

\subsubsection{Triedro de Frenet}
\begin{itemize}
    \item \(\mathbf{T}(s) = \frac{\mathbf{r}'(t)}{\|\mathbf{r}'(t)\|} = \frac{1}{\sqrt{2}} (\cos t - \sin t, \sin t + \cos t)\)
    \item \(\mathbf{N}(s)\): Se calcula como \(\frac{\mathbf{T}'(s)}{\|\mathbf{T}'(s)\|}\)
    \item \(\mathbf{B}(s) = \mathbf{T}(s) \times \mathbf{N}(s)\) (en 3D, pero como es plana, \(\mathbf{B}(s) = (0,0,1)\))
\end{itemize}

\subsubsection{Cálculo de la curvatura}
\noindent
La curvatura se define como:
\[
\kappa(s) = \|\mathbf{T}'(s)\|
\]
Primero calculamos \(\mathbf{T}'(t)\):
\[
\mathbf{T}'(t) = %\frac{d}{dt}\left(\frac{\mathbf{r}'(t)}{\|\mathbf{r}'(t)\|}\right) = \frac{1}{\sqrt{2}} (-\sin t - \cos t, \cos t - \sin t) - \frac{e^t}{\sqrt{2} e^t} \mathbf{T}(t) 
= \frac{1}{\sqrt{2}} (-\sin t - \cos t, \cos t - \sin t)
\]
Su norma es:
\[
\|\mathbf{T}'(t)\| %= \frac{1}{\sqrt{2}} \sqrt{(-\sin t - \cos t)^2 + (\cos t - \sin t)^2} 
= \frac{1}{\sqrt{2}} \sqrt{2 + 2 \sin t \cos t - 2 \cos t \sin t} = 1
\]
Pero como \(\|\mathbf{T}'(s)\| = \frac{\|\mathbf{T}'(t)\|}{\|\mathbf{r}'(t)\|} = \frac{1}{e^t \sqrt{2}}\), y usando la relación \(ds = e^t \sqrt{2} dt\), obtenemos:
\[
\kappa(s) = \frac{1}{e^t \sqrt{2}} = \frac{1}{\sqrt{2} \left(\frac{s}{\sqrt{2}} + 1\right)} = \frac{1}{s + \sqrt{2}}
\]

\subsubsection{Cálculo de la torsión}
\noindent Como la curva es plana, la torsión es:
\[
\tau(s) = 0.
\]

%\chapter{Curvas (ter)}

%\section{Introducción}

%La curvatura de una curva en el espacio se define como la norma de la derivada segunda de la parametrización por longitud de arco. Ahora, derivamos la fórmula de la curvatura para una parametrización arbitraria.

\subsection{Curvatura para cualquier parametrizaci\'on}

%\section{Definiciones}
\subsubsection{Preliminarios}

Sea $\boldsymbol{\alpha}(s)$ una parametrización por longitud de arco de una curva, donde $s$ es la longitud de arco. La curvatura $\kappa$ se define como:
\begin{equation*}
\kappa = \left\| \frac{d^2 \boldsymbol{\alpha}}{ds^2} \right\|
\end{equation*}
%
%\section{Parametrización General}
%
Consideremos una parametrización general $\mathbf{r}(t)$, donde $t$ es un parámetro arbitrario. La relación entre la derivada respecto a $t$ y la derivada respecto a $s$ es:
\begin{equation*}
\frac{\mathrm{d}\boldsymbol{\alpha}}{\mathrm{d}s} = \frac{d\mathbf{r}/\mathrm{d}t}{\mathrm{d}s/\mathrm{d}t}.
\end{equation*}

\subsubsection{Derivación de la Fórmula de Curvatura}

Queremos expresar ${\mathrm{d}^2 \boldsymbol{s}}/{\mathrm{d}s^2}$ en términos de %$\mathbf{r}(t)$, 
$\mathbf{r}'(t)$ y $\mathbf{r}''(t)$. Primero, notemos que la relación anterior combinada a $\| \boldsymbol{\alpha}(s)=1\|$ implica:
\begin{equation*}
\frac{\mathrm{d}s}{\mathrm{d}t} = \left\| \frac{d\mathbf{r}}{\mathrm{d}t} \right\|.
\end{equation*}
Despues calculamos con la regla de la cadena (ver antes):
\[
\mathbf{r}'(t) = \boldsymbol{\alpha}'(s) \frac{\mathrm{d}s}{\mathrm{d}t}.
\]
Derivamos otra vez:
\[
\mathbf{r}"(t) = 
\boldsymbol{\alpha}"(s) \left ( \frac{\mathrm{d}s}{\mathrm{d}t} \right )^2
+
\boldsymbol{\alpha}'(s) \frac{\mathrm{d}^2s}{\mathrm{d}t^2}.
\]
Combinamos las dos formulas:
\[
\mathbf{r}'(t) \times \mathbf{r}"(t) 
= \boldsymbol{\alpha}'(s) \frac{\mathrm{d}s}{\mathrm{d}t} \times
\left (
\boldsymbol{\alpha}"(s) \left ( \frac{\mathrm{d}s}{\mathrm{d}t} \right )^2
+
\boldsymbol{\alpha}'(s) \frac{\mathrm{d}^2s}{\mathrm{d}t^2} \right ).
\]
Dado que \( \boldsymbol{\alpha}'(s) \times \boldsymbol{\alpha}'(s) = 0\), obtenemos
\[
\mathbf{r}'(t) \times \mathbf{r}"(t) 
= \boldsymbol{\alpha}'(s) \frac{\mathrm{d}s}{\mathrm{d}t} \times
\boldsymbol{\alpha}"(s) \left ( \frac{\mathrm{d}s}{\mathrm{d}t} \right )^2.
\]
Utilizamos el hecho de que $\boldsymbol{\alpha}'(s)$ y $\boldsymbol{\alpha}"(s)$ son ortogonales, y tomamos la norma:
\[
\| \mathbf{r}'(t) \times \mathbf{r}"(t) \|
= \frac{\mathrm{d}s}{\mathrm{d}t} 
\kappa(s) \left ( \frac{\mathrm{d}s}{\mathrm{d}t} \right )^2.
\]
%
%
%Luego, la derivada segunda respecto a $s$ es:
%
%\begin{align*}
%\frac{d^2 \mathbf{r}}{ds^2} &= \frac{d}{ds}\left(\frac{d\mathbf{r}/dt}{ds/dt}\right) \\
%&= \frac{d}{dt}\left(\frac{d\mathbf{r}/dt}{ds/dt}\right) \cdot \frac{dt}{ds} \\
%&= \left[\frac{d^2\mathbf{r}/dt^2}{ds/dt} - \frac{d\mathbf{r}/dt}{(ds/dt)^2} \cdot \frac{d^2s}{dt^2}\right] \cdot \frac{dt}{ds}
%\end{align*}
%
%Simplificando y usando el hecho de que $\frac{dt}{ds} = \frac{1}{ds/dt}$, obtenemos:
%
%\begin{align*}
%\frac{d^2 \mathbf{r}}{ds^2} &= \frac{d^2\mathbf{r}/dt^2}{(ds/dt)^2} - \frac{d\mathbf{r}/dt}{(ds/dt)^3} \cdot \frac{d^2s}{dt^2} \\
%&= \frac{\mathbf{r}''(t) \|\mathbf{r}'(t)\|^2 - \mathbf{r}'(t) (\mathbf{r}''(t) \cdot \mathbf{r}'(t))}{\|\mathbf{r}'(t)\|^4}
%\end{align*}
%
Finalmente, la curvatura $\kappa$ es:
\begin{equation*}
\kappa %= \left\| \frac{d^2 \mathbf{r}}{ds^2} \right\| 
= \frac{\|\mathbf{r}'(t) \times \mathbf{r}''(t)\|}{\|\mathbf{r}'(t)\|^3}.
\end{equation*}

\subsection{La elipse}

\noindent
Vamos a detallar las fórmulas para una elipse con semiejes \(a\) (eje mayor) y \(b\) (eje menor), centrada en el origen y alineada con los ejes coordenados.
%
%---
%
%### **1. Parametrización de la elipse**
La parametrización estándar de una elipse es:
\[
x(t) = a \cos t, \quad y(t) = b \sin t, \quad t \in [0, 2\pi].
\]
%
%---
%
%### **2. Longitud de arco \(s(t)\)**
La longitud de arco \(s(t)\) desde \(t=0\) hasta un parámetro \(t\) se calcula como:
\[
s(t) = \int_{0}^{t} \sqrt{\left(\frac{dx}{dt}\right)^2 + \left(\frac{dy}{dt}\right)^2} \, dt
\]
Sustituyendo las derivadas:
\[
\frac{dx}{dt} = -a \sin t, \quad \frac{dy}{dt} = b \cos t
\]
Por lo tanto:
\[
s(t) = \int_{0}^{t} \sqrt{a^2 \sin^2 t + b^2 \cos^2 t} \, dt.
\]
Esta integral no tiene una solución analítica cerrada en términos de funciones elementales. Se suele expresar en términos de integrales elípticas de segunda especie \(E(t|k)\):
\[
s(t) = a \, E(t|k), \quad \text{donde } k = \sqrt{1 - \frac{b^2}{a^2}}
\]
%
%---
%
%### **3. Curvatura \(\kappa(t)\)**
La curvatura de una curva parametrizada \((x(t), y(t))\) está dada por:
\[
\kappa(t) = \frac{|x'(t)y''(t) - y'(t)x''(t)|}{\left( (x'(t))^2 + (y'(t))^2 \right)^{3/2}}
\]
Para la elipse:
\[
x'(t) = -a \sin t, \quad y'(t) = b \cos t
\]
\[
x''(t) = -a \cos t, \quad y''(t) = -b \sin t
\]
Sustituyendo:
\[
\kappa(t) = \frac{|(-a \sin t)(-b \sin t) - (b \cos t)(-a \cos t)|}{\left(a^2 \sin^2 t + b^2 \cos^2 t\right)^{3/2}} = \frac{ab}{\left(a^2 \sin^2 t + b^2 \cos^2 t\right)^{3/2}}
\]

\subsection{Plano osculador y plano normal}

\noindent
Consideremos una curva parametrizada en 3D, dada por la función vectorial:
\begin{equation*}
\mathbf{r}(t) = (x(t), y(t), z(t))
\end{equation*}
donde \( t \) es el parámetro.
\subsubsection{Plano Osculador}

\noindent
El plano osculador en un punto \(\mathbf{r}(t)\) de la curva es el plano que mejor aproxima a la curva en ese punto. Está definido por los vectores tangente y normal principal a la curva en ese punto.
\begin{itemize}
\item El vector tangente \(\mathbf{T}(t)\) está dado por:
\begin{equation*}
\mathbf{T}(t) = \frac{\mathbf{r}'(t)}{\|\mathbf{r}'(t)\|}
\end{equation*}
\item El vector normal principal \(\mathbf{N}(t)\) está dado por:
\begin{equation*}
\mathbf{N}(t) = \frac{\mathbf{T}'(t)}{\|\mathbf{T}'(t)\|}
\end{equation*}
\end{itemize}
El plano osculador está definido por el punto \(\mathbf{r}(t)\) y los vectores \(\mathbf{T}(t)\) y \(\mathbf{N}(t)\). Su ecuación es:
\begin{equation*}
\det(\mathbf{X} - \mathbf{r}(t), \mathbf{T}(t), \mathbf{N}(t)) = 0
\end{equation*}
donde \(\mathbf{X}\) es un punto genérico del plano.

\subsubsection{Plano Normal}

\noindent
El plano normal en un punto \(\mathbf{r}(t)\) de la curva es el plano que contiene al vector normal principal \(\mathbf{N}(t)\) y es perpendicular al vector tangente \(\mathbf{T}(t)\). Está definido por el punto \(\mathbf{r}(t)\) y el vector \(\mathbf{N}(t)\). Su ecuación es:
\begin{equation*}
\mathbf{N}(t) \cdot (\mathbf{X} - \mathbf{r}(t)) = 0
\end{equation*}
donde \(\mathbf{X}\) es un punto genérico del plano.

\subsection{Fórmula de la Torsión para Cualquier Parametrizac\'ion}

\paragraph{Formula}
Consideremos una curva parametrizada en 3D, dada por la función vectorial:
\begin{equation*}
\mathbf{r}(t) = (x(t), y(t), z(t))
\end{equation*}
donde \( t \) es el parámetro.
La torsión \(\tau(t)\) de la curva en un punto \(\mathbf{r}(t)\) se define como:
\begin{equation*}
\tau(t) = \frac{(\mathbf{r}'(t) \times \mathbf{r}''(t)) \cdot \mathbf{r}'''(t)}{\|\mathbf{r}'(t) \times \mathbf{r}''(t)\|^2}
\end{equation*}

\paragraph{Ejemplo: Curva \((\cos t, \sin t, t^2/8)\)}

\noindent
Consideremos la curva parametrizada por:
\begin{equation*}
\mathbf{r}(t) = (\cos t, \sin t, \frac{t^2}{8})
\end{equation*}
Calculamos las derivadas:
\begin{align*}
\mathbf{r}'(t) &= (-\sin t, \cos t, \frac{t}{4}), \\
\mathbf{r}''(t) &= (-\cos t, -\sin t, \frac{1}{4}), \\
\mathbf{r}'''(t) &= (\sin t, -\cos t, 0).
\end{align*}
Calculamos el producto cruz \(\mathbf{r}'(t) \times \mathbf{r}''(t)\):
\begin{align*}
\mathbf{r}'(t) \times \mathbf{r}''(t) &= \begin{vmatrix}
\mathbf{i} & \mathbf{j} & \mathbf{k} \\
-\sin t & \cos t & \frac{t}{4} \\
-\cos t & -\sin t & \frac{1}{4}
\end{vmatrix} \\
&= \left( \cos t \cdot \frac{1}{4} - \frac{t}{4} \cdot (-\sin t) \right) \mathbf{i} - \left( -\sin t \cdot \frac{1}{4} - \frac{t}{4} \cdot (-\cos t) \right) \mathbf{j} \\\ & + \left( -\sin t \cdot (-\sin t) - \cos t \cdot (-\cos t) \right) \mathbf{k} \\
&= \left( \frac{\cos t + t \sin t}{4} \right) \mathbf{i} + \left( \frac{\sin t - t \cos t}{4} \right) \mathbf{j} + (\sin^2 t + \cos^2 t) \mathbf{k} \\
&= \left( \frac{\cos t + t \sin t}{4} \right) \mathbf{i} + \left( \frac{\sin t - t \cos t}{4} \right) \mathbf{j} + \mathbf{k}
\end{align*}
Calculamos el producto escalar \((\mathbf{r}'(t) \times \mathbf{r}''(t)) \cdot \mathbf{r}'''(t)\):
\begin{align*}
(\mathbf{r}'(t) \times \mathbf{r}''(t)) \cdot \mathbf{r}'''(t) &= \left( \frac{\cos t + t \sin t}{4} \right) \sin t + \left( \frac{\sin t - t \cos t}{4} \right) (-\cos t) + 1 \cdot 0 \\
&= \frac{(\cos t + t \sin t) \sin t + (-\sin t + t \cos t) \cos t}{4} \\
&= \frac{\cos t \sin t + t \sin^2 t - \sin t \cos t + t \cos^2 t}{4} \\
&= \frac{t (\sin^2 t + \cos^2 t)}{4} \\
&= \frac{t}{4}
\end{align*}
Calculamos la norma del producto cruz:
\begin{align*}
\|\mathbf{r}'(t) \times \mathbf{r}''(t)\|^2 &= \left( \frac{\cos t + t \sin t}{4} \right)^2 + \left( \frac{\sin t - t \cos t}{4} \right)^2 + 1^2 \\
&= \frac{(\cos t + t \sin t)^2 + (\sin t - t \cos t)^2}{16} + 1 \\
&= \frac{\cos^2 t + 2 t \cos t \sin t + t^2 \sin^2 t + \sin^2 t - 2 t \sin t \cos t + t^2 \cos^2 t}{16} + 1 \\
&= \frac{\cos^2 t + \sin^2 t + t^2 (\sin^2 t + \cos^2 t)}{16} + 1 \\
&= \frac{1 + t^2}{16} + 1 \\
&= \frac{17 + t^2}{16}
\end{align*}
Por lo tanto, la torsión es:
\begin{equation*}
\tau(t) = \frac{\frac{t}{4}}{\frac{17 + t^2}{16}} = \frac{4t}{17 + t^2}.
\end{equation*}

\subsection{Más sobre Frenet-Serret}

%\section{Introducción}
\noindent En geometría diferencial, el {Teorema Fundamental de las Curvas} establece que una curva en \( \mathbb{R}^3 \) está completamente determinada (salvo movimientos rígidos) por sus funciones de {curvatura} \( \kappa(s) \) y {torsión} \( \tau(s) \), donde \( s \) es el parámetro de arco. Esto permite {construir} curvas directamente a partir de estas propiedades.

\subsubsection{Teorema Fundamental de las Curvas}
\begin{quote}
    Dadas dos funciones continuas \( \kappa(s) > 0 \) y \( \tau(s) \) definidas para \( s \in [0, L] \), existe una curva única \( \mathbf{r}(s) \) (salvo rotaciones y traslaciones) cuya curvatura y torsión son \( \kappa(s) \) y \( \tau(s) \), respectivamente.
\end{quote}

\subsubsection{Ecuaciones de Frenet-Serret}
\noindent La curva se reconstruye resolviendo el siguiente sistema de ecuaciones diferenciales para el {triedro de Frenet} \( \{\mathbf{T}(s), \mathbf{N}(s), \mathbf{B}(s)\} \):
\[
\begin{cases}
\mathbf{T}'(s) = \kappa(s) \mathbf{N}(s), \\
\mathbf{N}'(s) = -\kappa(s) \mathbf{T}(s) + \tau(s) \mathbf{B}(s), \\
\mathbf{B}'(s) = -\tau(s) \mathbf{N}(s).
\end{cases}
\]
La curva \( \mathbf{r}(s) \) se obtiene integrando:
\[
\mathbf{r}(s) = \mathbf{r}(0) + \int_0^s \mathbf{T}(u) \, du.
\]

\subsubsection{Ejemplo Clásico: La Hélice Circular}
\noindent Para una hélice con radio \( R \) y paso \( h \):
\begin{itemize}
    \item Curvatura: \( \kappa(s) = \frac{R}{R^2 + (h/2\pi)^2} \) (constante).
    \item Torsión: \( \tau(s) = \frac{h/2\pi}{R^2 + (h/2\pi)^2} \) (constante).
\end{itemize}
La curva resultante es:
\[
\mathbf{r}(s) = \left( R \cos\left(\frac{s}{c}\right), R \sin\left(\frac{s}{c}\right), \frac{h s}{2\pi c} \right), \quad c = \sqrt{R^2 + \left(\frac{h}{2\pi}\right)^2}.
\]

\subsubsection{Implementación Numérica en Python}
\noindent Para reconstruir una curva a partir de \( \kappa(s) \) y \( \tau(s) \), se puede resolver numéricamente el sistema de Frenet-Serret. A continuación, se muestra un ejemplo en Python:

\begin{lstlisting}[language=Python, caption=Reconstrucción de una curva a partir de \( \kappa(s) \) y \( \tau(s) \)]
import numpy as np
from scipy.integrate import solve_ivp
import matplotlib.pyplot as plt
from mpl_toolkits.mplot3d import Axes3D

# Definir curvatura y torsion (ejemplo: constantes)
kappa = lambda s: 0.5
tau = lambda s: 0.2

# Ecuaciones de Frenet-Serret
def frenet_serret(s, y):
    T, N, B = y.reshape(3, 3)  # Vectores T, N, B apilados
    dTds = kappa(s) * N
    dNds = -kappa(s) * T + tau(s) * B
    dBds = -tau(s) * N
    return np.concatenate([dTds, dNds, dBds])

# Condiciones iniciales (T(0) = [1,0,0], N(0) = [0,1,0], B(0) = [0,0,1])
y0 = np.array([1, 0, 0, 0, 1, 0, 0, 0, 1])

# Resolver el sistema de EDOs
sol = solve_ivp(frenet_serret, [0, 10], y0, t_eval=np.linspace(0, 10, 100), method='RK45')

# Reconstruir la curva integrando T(s)
T = sol.y[:3, :]  # Vector tangente en cada punto
r = np.zeros_like(T)
for i in range(1, len(sol.t)):
    r[:, i] = r[:, i-1] + T[:, i] * (sol.t[i] - sol.t[i-1])

# Graficar
fig = plt.figure()
ax = fig.add_subplot(111, projection='3d')
ax.plot(r[0, :], r[1, :], r[2, :], label='Curva reconstruida')
ax.set_xlabel('X')
ax.set_ylabel('Y')
ax.set_zlabel('Z')
plt.legend()
plt.show()
\end{lstlisting}

\subsubsection{Aplicaciones}
\begin{itemize}
    \item \textbf{Robótica}: Planificación de trayectorias suaves para brazos robóticos.
    \item \textbf{Diseño de carreteras}: Curvas con curvatura controlada para seguridad.
    \item \textbf{Biología}: Modelado de estructuras como el ADN o proteínas.
\end{itemize}

\subsubsection{Referencias}
\begin{itemize}
    \item do Carmo, M. (1976). \textit{Differential Geometry of Curves and Surfaces}. Prentice-Hall.
    \item Kreyszig, E. (1991). \textit{Differential Geometry}. Dover.
\end{itemize}

\subsection{Curvas de Bézier en 3D: Cálculo de Curvatura y Torsión}

\noindent Las {curvas de Bézier} son herramientas fundamentales en el diseño geométrico asistido por computadora (CAGD) y la computación gráfica. Una curva de Bézier en 3D se define mediante un conjunto de puntos de control \( \mathbf{P}_0, \mathbf{P}_1, \ldots, \mathbf{P}_n \) y está dada por la ecuación paramétrica:
\[
\mathbf{B}(t) = \sum_{i=0}^n \binom{n}{i} (1-t)^{n-i} t^i \mathbf{P}_i, \quad t \in [0,1].
\]

\subsubsection{Cálculo de la Curvatura}
\noindent La {curvatura} \( \kappa(t) \) de una curva paramétrica \( \mathbf{B}(t) \) se calcula como:
\[
\kappa(t) = \frac{\|\mathbf{B}'(t) \times \mathbf{B}''(t)\|}{\|\mathbf{B}'(t)\|^3},
\]
donde \( \mathbf{B}'(t) \) y \( \mathbf{B}''(t) \) son la primera y segunda derivada de la curva, respectivamente.

\subsubsection{Ejemplo en Python}
\noindent A continuación, se muestra cómo calcular la curvatura para una curva de Bézier cúbica en 3D:

\begin{lstlisting}[language=Python, caption=Cálculo de la curvatura en Python]
import numpy as np
import sympy as sp

# Puntos de control de una curva de Bezier cubica en 3D
P = np.array([[0, 0, 0], [1, 2, 1], [3, 1, 2], [4, 0, 0]])

# Funcion de Bezier
def bezier_curve(P, t):
    n = len(P) - 1
    curve = np.zeros(3)
    for i in range(len(P)):
        curve += sp.binomial(n, i) * (t**i) * ((1-t)**(n-i)) * P[i]
    return curve

t = sp.symbols('t')
curve = bezier_curve(P, t)

# Primera y segunda derivada
d1 = sp.diff(curve, t)
d2 = sp.diff(d1, t)

# Producto cruz y norma
cross = d1.cross(d2)
norm_cross = sp.sqrt(cross.dot(cross))
norm_d1 = sp.sqrt(d1.dot(d1))

# Curvatura
curvature = norm_cross / norm_d1**3
sp.simplify(curvature)  # Simplifica la expresion
\end{lstlisting}

\subsubsection{Cálculo de la Torsión}
\noindent La {torsión} \( \tau(t) \) mide cómo la curva se desvía de un plano y se calcula como:
\[
\tau(t) = \frac{(\mathbf{B}'(t) \times \mathbf{B}''(t)) \cdot \mathbf{B}'''(t)}{\|\mathbf{B}'(t) \times \mathbf{B}''(t)\|^2}.
\]

\subsubsection{Ejemplo en Python}
\noindent Cálculo de la torsión para la misma curva de Bézier:

\begin{lstlisting}[language=Python, caption=Cálculo de la torsión en Python]
# Tercera derivada
d3 = sp.diff(d2, t)

# Producto punto y norma
dot_product = (d1.cross(d2)).dot(d3)
norm_cross_squared = norm_cross**2

# Torsion
torsion = dot_product / norm_cross_squared
sp.simplify(torsion)  # Simplifica la expresion
\end{lstlisting}

\subsubsection{Control de Curvatura y Torsión}
\noindent Para {controlar} la curvatura y torsión, se pueden ajustar los puntos de control \( \mathbf{P}_i \). Por ejemplo:
- Aumentar la distancia entre puntos de control suele aumentar la curvatura.
- La torsión se ve afectada por la no coplanaridad de los puntos de control.

\subsubsection{Ejemplo Visual}
\noindent Puedes visualizar la curva y sus propiedades con \texttt{matplotlib}:

\begin{lstlisting}[language=Python, caption=Visualización de la curva de Bézier]
import matplotlib.pyplot as plt
from mpl_toolkits.mplot3d import Axes3D

# Evaluar la curva en puntos discretos
t_vals = np.linspace(0, 1, 100)
curve_vals = np.array([bezier_curve(P, ti).evalf(subs={t: ti}) for ti in t_vals])

# Graficar
fig = plt.figure()
ax = fig.add_subplot(111, projection='3d')
ax.plot(curve_vals[:, 0], curve_vals[:, 1], curve_vals[:, 2], label='Curva de Bezier')
ax.scatter(P[:, 0], P[:, 1], P[:, 2], c='red', label='Puntos de control')
ax.legend()
plt.show()
\end{lstlisting}

\subsubsection{Referencias}
\begin{itemize}
    \item Farin, G. (2002). \textit{Curves and Surfaces for CAGD}. Morgan Kaufmann.
    \item Piegl, L. \& Tiller, W. (1997). \textit{The NURBS Book}. Springer.
    \item do Carmo, M. (1976). \textit{Differential Geometry of Curves and Surfaces}. Prentice-Hall.
\end{itemize}

\newpage

\section{Practico Curvas}

\begin{enumerate}

\item Pruebe que la curva $x(t)=t^{2}$, $y(t)=1-3t$, $z(t)=1+t^{3}$ pasa por los puntos $(1,4,0)$ y $(9,-8,28)$ pero no por $(4,7,-6)$.

\item Graficar las siguientes curvas, y dibujar la posición y el vector tangente en los puntos que se indican:
\begin{enumerate}
\item $x(t)=\cos(t)$, $y(t)=\sen(t)$ en  $t=\pi/4$
\item $x(t)=t^{3}$, $y(t)=t^{2}$, en  $t=1$
\item $x(t)=2 \sen(t)$, $y(t)=3\cos(t)$, en  $t=\pi/3$.     
 \end{enumerate}
 
\item Calcular la longitud de las siguientes curvas:
\begin{enumerate}

\item $x(t)=t^2$, $y(t)=t$ con $t\in[0,1]$.
\item $x(t)=2t$, $y(t)=3\sen(t)$, $z(t)=3\cos(t)$ para $0\leq t \leq \pi$.
\item $x(t)=e^{t}$, $y(t)=e^{t}\sen(t)$, $z(t)=e^{t }\cos(t)$ para $0\leq t \leq 2\pi$.

\end{enumerate}


\item Hallar la curvatura de las siguientes curvas y hacer un bosquejo local de la curva en dos puntos a elecci\'on.

\begin{enumerate}
 \item      $x(t)=1$, $y(t)=t$, $z(t)=t^{2}$
  \item      $x(t)=t^{2}+2$, $y(t)=t^{2}-4t$, $z(t)=2t$
   \item      $x(t)=\sen(t)$, $y(t)=\cos(t)$, $z(t)=\sen(t)$
   \item $y= sen x$
   \item $y=x^{3}$.
\end{enumerate}


\item Hallar el radio de curvatura y la torsi\'on de $x(t)=t^{2}$, $y(t)=2t^{3}$ y hacer un bosquejo local en dos puntos a elecci\'on.

\item Calcular el triedro de Frenet, curvatura y torsi\'on de la \emph{h\'elice}: $\alpha(t)=(\cos(t),\sen(t),t)$. Bosquejar.

\item Hallar curvatura y torsi\'on de $\alpha(t)=(t^2,t^3,t^4)$ y hacer un bosquejo local en dos puntos a elecci\'on.

\item Hallar la curvatura y la torsi\'on de las siguientes curvas y hacer bosquejos locales en un punto a elecci\'on:
\begin{enumerate}
 \item $x(t)=e^{t}$, $y(t)=e^{-t}$, $z(t)=t\sqrt{2}$
  \item $x(t)=e^{-t}\sen(t)$, $y(t)=e^{-t}cost$, $z(t)=e^{t}$.
\end{enumerate}

\item Hallar el triedro de Frenet en un punto gen\'erico para las siguientes curvas.
\begin{enumerate}
   \item      $x(t)=\senh t$, $y(t)=\cosh t$, $z(t)=t$
  \item      $x(t)=t$, $y(t)=t/\sqrt{2}$, $z(t)=t^{3}/3$.

\end{enumerate}

\end{enumerate}


\bibliographystyle{siam}
\bibliography{biblio}

\end{document}